\documentclass[11pt,a4paper]{article}
\usepackage[T1]{fontenc}
\usepackage[utf8]{inputenc}
\usepackage[italian]{babel}
\usepackage{amsmath,amssymb,enumitem,xcolor}
\usepackage[a4paper,margin=2.2cm]{geometry}
\newcounter{quiz}
\newenvironment{quiz}{\par\medskip\noindent\refstepcounter{quiz}\textbf{\thequiz)}\enspace\ignorespaces}{\par\bigskip}
\newcommand{\registraSoluzione}[2]{}
\begin{document}
\section*{Quiz Analisi 1 -- Paola Suria}
\input{quiz/quiz_pagine_001_010.tex}
\input{quiz/quiz_pagine_011_020.tex}
\input{quiz/quiz_pagine_021_030.tex}
\input{quiz/quiz_pagine_031_040.tex}
\input{quiz/quiz_pagine_041_050.tex}
\input{quiz/quiz_pagine_051_060.tex}
\input{quiz/quiz_pagine_061_070.tex}
\input{quiz/quiz_pagine_071_080.tex}
\input{quiz/quiz_pagine_081_090.tex}
\input{quiz/quiz_pagine_091_100.tex}
\input{quiz/quiz_pagine_101_110.tex}
\input{quiz/quiz_pagine_111_120.tex}
\input{quiz/quiz_pagine_121_130.tex}
\input{quiz/quiz_pagine_131_140.tex}
\input{quiz/quiz_pagine_141_150.tex}
\input{quiz/quiz_pagine_151_160.tex}
\input{quiz/quiz_pagine_161_170.tex}
\input{quiz/quiz_pagine_171_180.tex}
\input{quiz/quiz_pagine_181_190.tex}
\input{quiz/quiz_pagine_191_200.tex}
\input{quiz/quiz_pagine_201_210.tex}
\input{quiz/quiz_pagine_211_220.tex}
\input{quiz/quiz_pagine_221_230.tex}
\input{quiz/quiz_pagine_231_240.tex}
\input{quiz/quiz_pagine_241_250.tex}
\input{quiz/quiz_pagine_251_260.tex}
\input{quiz/quiz_pagine_261_270.tex}
\input{quiz/quiz_pagine_271_280.tex}
\input{quiz/quiz_pagine_281_290.tex}
\input{quiz/quiz_pagine_291_300.tex}
\input{quiz/quiz_pagine_301_310.tex}
\input{quiz/quiz_pagine_311_320.tex}
\input{quiz/quiz_pagine_321_330.tex}
\input{quiz/quiz_pagine_331_340.tex}
\input{quiz/quiz_pagine_341_350.tex}
\input{quiz/quiz_pagine_351_360.tex}
\input{quiz/quiz_pagine_361_370.tex}
\input{quiz/quiz_pagine_371_380.tex}
\input{quiz/quiz_pagine_381_390.tex}
\input{quiz/quiz_pagine_391_400.tex}
\input{quiz/quiz_pagine_401_410.tex}
\input{quiz/quiz_pagine_411_420.tex}
\input{quiz/quiz_pagine_421_430.tex}
\input{quiz/quiz_pagine_431_440.tex}
\input{quiz/quiz_pagine_441_450.tex}
\input{quiz/quiz_pagine_451_460.tex}
\input{quiz/quiz_pagine_461_470.tex}
\input{quiz/quiz_pagine_471_480.tex}
\input{quiz/quiz_pagine_481_490.tex}
\input{quiz/quiz_pagine_491_500.tex}
\input{quiz/quiz_pagine_501_510.tex}
\input{quiz/quiz_pagine_511_520.tex}
\input{quiz/quiz_pagine_521_521.tex}
\clearpage\section*{Soluzioni e spiegazioni per capitolo (AM1 Morandotti)}
\subsection*{Capitolo 1}
\subsubsection*{1.1 -- Insiemi}
\noindent\textbf{quiz\_pagine\_111\_120\_017}: C. I termini $-n/(n+3)$ appartengono a $(-1,0]$, includendo $0$; inoltre $-3\in A$. Pertanto $\max A=0$ e $\min A=-3$.\par\medskip
\noindent\textbf{quiz\_pagine\_111\_120\_025}: D. L'esistenza di un estremo inferiore finito garantisce solo limitatezza inferiore, non superiore.\par\medskip
\noindent\textbf{quiz\_pagine\_131\_140\_021}: D. Dalla disuguaglianza tra estremo superiore e inferiore si scelgono elementi $x\in A$ sufficientemente vicini a $\sup A$ e $y\in B$ vicini a $\inf B$, ottenendo $x>y$. Gli estremi possono essere estesi; le altre opzioni non sono garantite.\par\medskip
\noindent\textbf{quiz\_pagine\_161\_170\_021}: E. Avere lo stesso estremo superiore non implica che sia raggiunto: ad esempio $B=(0,3)$ non ha massimo.\par\medskip
\noindent\textbf{quiz\_pagine\_161\_170\_022}: B. L'infimo e\textbackslash{}'e il massimo dei minoranti, ma non necessariamente appartiene all'insieme; ad esempio $A=(5,6)$.\par\medskip
\noindent\textbf{quiz\_pagine\_171\_180\_017}: E. Caratterizzazione del supremo: per ogni $\varepsilon>0$ entrambi gli insiemi hanno un elemento in $(4-\varepsilon,4]$.\par\medskip
\noindent\textbf{quiz\_pagine\_191\_200\_005}: B. $\min A=-1$ (per $n=1$) e $\max A=2$.\par\medskip
\noindent\textbf{quiz\_pagine\_221\_230\_005}: E. Ogni minorante di $B$ e' anche minorante di $A$, dunque $\inf A\geq\inf B$ (quando gli estremi inferiori sono definiti).\par\medskip
\noindent\textbf{quiz\_pagine\_271\_280\_001}: E. Ogni minorante di $B$ e' minorante di $A$: quindi $\inf A\geq\inf B$ (per insiemi con estremo inferiore definito).\par\medskip
\noindent\textbf{quiz\_pagine\_281\_290\_007}: A. Un infimo finito assicura un limite inferiore ma non superiore: $[0,+\infty)$ ha infimo $0$ ed e' illimitato.\par\medskip
\noindent\textbf{quiz\_pagine\_291\_300\_016}: E. Il massimo e' $0$ e il minimo e' $-3$. Anche l'opzione D e' vera: il quesito presenta due risposte corrette.\par\medskip
\noindent\textbf{quiz\_pagine\_321\_330\_024}: A. Per inclusione, ogni maggiorante di $A$ \textbackslash{}`e maggiorante di $B$, quindi $\sup B\leq\sup A$ quando i supremi sono definiti. Il testo non specifica che gli insiemi siano non vuoti.\par\medskip
\noindent\textbf{quiz\_pagine\_391\_400\_016}: B. L\textbackslash{}textquotesingle elemento $5$ e massimo; $1+1/(3n)$ tende a $1$ dall\textbackslash{}textquotesingle alto senza mai raggiungerlo, quindi non esiste minimo.\par\medskip
\noindent\textbf{quiz\_pagine\_391\_400\_017}: D. Esempio $A=(0,1)$, $B=[0,1]$: sono distinti e $A\subset B$, ma entrambi hanno infimo $0$ e supremo $1$.\par\medskip
\noindent\textbf{quiz\_pagine\_401\_410\_018}: E. Per inclusione, ogni minorante di $B$ \textbackslash{}`e minorante di $A$, quindi $\inf A\geq\inf B$ (quando gli estremi inferiori sono definiti).\par\medskip
\noindent\textbf{quiz\_pagine\_491\_500\_001}: E. Per definizione di estremo superiore, ogni intorno sinistro di $4$ contiene elementi di entrambi gli insiemi. Non occorre che $A$ e $B$ abbiano elementi comuni.\par\medskip
\subsubsection*{1.2 -- Logica delle proposizioni e dei predicati}
\noindent\textbf{quiz\_pagine\_181\_190\_025}: non verificata / nessuna alternativa corretta. La E esprime $A=\varnothing$, ossia la negazione di $A\neq\varnothing$. Attenzione: la proposizione originale e' una equivalenza logicamente sempre vera; la sua negazione formale e' sempre falsa, e nessuna alternativa rappresenta tale negazione. Risposta ufficiale E riportata con riserva.\par\medskip
\noindent\textbf{quiz\_pagine\_191\_200\_012}: D. $e^{\sqrt n}\to+\infty$; pertanto supera definitivamente ogni soglia positiva. Il minimo $1$ e' assunto per $n=0$.\par\medskip
\noindent\textbf{quiz\_pagine\_191\_200\_018}: non verificata / nessuna alternativa corretta. La limitatezza di $b_n$ e' sufficiente ma non necessaria: $b_n=n$, $c_n=1/n^2$ danno $a_n\to0$. Nessuna alternativa descrive correttamente una condizione necessaria e sufficiente. La risposta ufficiale E va intesa solo come condizione sufficiente.\par\medskip
\noindent\textbf{quiz\_pagine\_221\_230\_007}: C. Per $z\in A$, $-1<\operatorname{Re}z<1$; per $z\in B$, $\operatorname{Re}z>5$. Non possono intersecarsi.\par\medskip
\noindent\textbf{quiz\_pagine\_221\_230\_008}: B. Essendo i coefficienti reali, anche $2-i$ e' radice. Le tre radici distinte impongono grado almeno tre.\par\medskip
\noindent\textbf{quiz\_pagine\_261\_270\_008}: E. $w=(5+i)(3+i)/10=(14+8i)/10$; parte reale $7/5$. Nel testo compare $\operatorname{Re}(z)$ invece di $\operatorname{Re}(w)$.\par\medskip
\noindent\textbf{quiz\_pagine\_261\_270\_020}: D. $|z|\overline z=\sqrt{13}(2-3i)$ non e' reale; le altre espressioni hanno le proprieta' indicate.\par\medskip
\noindent\textbf{quiz\_pagine\_451\_460\_013}: D. Per definizione di estremo superiore e inferiore, $\sup A>\inf B$ garantisce elementi $x\in A$ e $y\in B$ con $x>y$. Le altre alternative non sono necessarie; A e C possono fallire per insiemi non limitati.\par\medskip
\noindent\textbf{quiz\_pagine\_461\_470\_016}: A. Il valore 5 appartiene ad $A$ ed \textbackslash{}`e il massimo; $\arctan(-n)\to-\pi/2$ senza raggiungerlo. ATTENZIONE: anche B \textbackslash{}`e vera, poich\textbackslash{}'e ogni massimo \textbackslash{}`e un estremo superiore; la domanda non ha risposta unica.\par\medskip
\noindent\textbf{quiz\_pagine\_461\_470\_017}: B. $\cosh x\ge1$ e non \textbackslash{}`e limitata superiormente. Aggiungendo $-9$ si ottiene il minimo $-9$ e il supremo $+\infty$.\par\medskip
\noindent\textbf{quiz\_pagine\_511\_520\_021}: E. Il sottoinsieme $A$ non pu\textbackslash{}`o contenere valori inferiori all'estremo inferiore di $B$: $\inf A\ge\inf B$, quando gli estremi sono definiti.\par\medskip
\subsubsection*{1.3 -- Relazioni}
\noindent\textbf{quiz\_pagine\_161\_170\_001}: C. $|w|=|1+3i|/|4+2i|=\sqrt{10}/\sqrt{20}=\sqrt2/2$. Nel testo le alternative (a) ed (e) sono identiche.\par\medskip
\noindent\textbf{quiz\_pagine\_301\_310\_010}: B. Dominio $x>-2$, $x\ne3$. Il segno del prodotto dei fattori mostra che la frazione e' non positiva in $(-2,-1]\cup[0,3)$.\par\medskip
\noindent\textbf{quiz\_pagine\_481\_490\_019}: B. $|w|=|1+i|/|3+4i|=\sqrt2/5$.\par\medskip
\subsubsection*{1.4 -- Funzioni}
\noindent\textbf{quiz\_pagine\_001\_010\_003}: B. Spiegazione non ancora disponibile.\par\medskip
\noindent\textbf{quiz\_pagine\_001\_010\_004}: B. Spiegazione non ancora disponibile.\par\medskip
\noindent\textbf{quiz\_pagine\_001\_010\_009}: E. Spiegazione non ancora disponibile.\par\medskip
\noindent\textbf{quiz\_pagine\_001\_010\_015}: A. Spiegazione non ancora disponibile.\par\medskip
\noindent\textbf{quiz\_pagine\_001\_010\_016}: E. Spiegazione non ancora disponibile.\par\medskip
\noindent\textbf{quiz\_pagine\_001\_010\_020}: E. Spiegazione non ancora disponibile.\par\medskip
\noindent\textbf{quiz\_pagine\_001\_010\_028}: E. Spiegazione non ancora disponibile.\par\medskip
\noindent\textbf{quiz\_pagine\_001\_010\_030}: B. Spiegazione non ancora disponibile.\par\medskip
\noindent\textbf{quiz\_pagine\_001\_010\_031}: A. Spiegazione non ancora disponibile.\par\medskip
\noindent\textbf{quiz\_pagine\_001\_010\_035}: D. Spiegazione non ancora disponibile.\par\medskip
\noindent\textbf{quiz\_pagine\_001\_010\_036}: E. Spiegazione non ancora disponibile.\par\medskip
\noindent\textbf{quiz\_pagine\_001\_010\_040}: B. Spiegazione non ancora disponibile.\par\medskip
\noindent\textbf{quiz\_pagine\_011\_020\_002}: B. Su $[0,+\infty)$ \textbackslash{}`e strettamente decrescente, dunque iniettiva.\par\medskip
\noindent\textbf{quiz\_pagine\_011\_020\_010}: D. $(f^{-1})\prime(2)=1/f\prime(0)=1/(2+3)=1/5$.\par\medskip
\noindent\textbf{quiz\_pagine\_011\_020\_011}: D. $e^{x^4}=1+x^4+o(x^4)$ e $\sqrt{1+x^4}=1+x^4/2+o(x^4)$; la differenza \textbackslash{}`e positiva.\par\medskip
\noindent\textbf{quiz\_pagine\_011\_020\_027}: E. $w=(5+i)(3+i)/10=(14+8i)/10$, parte reale $7/5$. Nel testo \textbackslash{}`e scritto $\operatorname{Re}(z)$ anzich\textbackslash{}'e $\operatorname{Re}(w)$.\par\medskip
\noindent\textbf{quiz\_pagine\_011\_020\_028}: E. $f\prime(x)=5x^4+3e^x>0$; $f(0)=1>0$ e $f(x)\to-\infty$ a sinistra.\par\medskip
\noindent\textbf{quiz\_pagine\_011\_020\_031}: B. $|w|=\sqrt{26}/\sqrt{10}=\sqrt{65}/5$.\par\medskip
\noindent\textbf{quiz\_pagine\_021\_030\_006}: E. Il radicando deve essere strettamente positivo, $1-x^2>0$, e $x+4>0$.\par\medskip
\noindent\textbf{quiz\_pagine\_021\_030\_007}: E. $49x^{48}+4>0$, quindi esiste un solo zero reale, positivo. Le radici non sono in generale equidistanti su una circonferenza.\par\medskip
\noindent\textbf{quiz\_pagine\_021\_030\_010}: D. $|z|z=\sqrt{13}(2+3i)$ non \textbackslash{}`e reale.\par\medskip
\noindent\textbf{quiz\_pagine\_021\_030\_011}: E. La base \textbackslash{}`e $e^{i\pi/3}$; elevando alla nona si ottiene $e^{i3\pi}=-1$.\par\medskip
\noindent\textbf{quiz\_pagine\_021\_030\_014}: A. $g\prime=-\sin(f)f\prime$: i tre zeri di $f$ e i due zeri distinti di $f\prime$ forniscono almeno cinque punti critici.\par\medskip
\noindent\textbf{quiz\_pagine\_021\_030\_016}: non verificata / nessuna alternativa corretta. I termini dominanti sono $\sqrt{x}$ al numeratore e $-4\sqrt{x}$ al denominatore: il limite risulterebbe $-1/4$. La lettura dell\textbackslash{}'espressione nel PDF richiede verifica.\par\medskip
\noindent\textbf{quiz\_pagine\_021\_030\_017}: C. $f(0)=3$, $f\prime(0)=6$, quindi $(f^{-1})\prime(3)=1/6$.\par\medskip
\noindent\textbf{quiz\_pagine\_021\_030\_018}: C. Vicino a $x=1$ il termine $|\cos(\pi x)|^3$ ha derivata nulla nel punto, mentre $2|x|$ ha derivata $2$.\par\medskip
\noindent\textbf{quiz\_pagine\_021\_030\_020}: D. Con $u=2-3x^2$, $du=-6x\,dx$, perci\textbackslash{}`o l\textbackslash{}'integrale \textbackslash{}`e $-u^6/36+C$.\par\medskip
\noindent\textbf{quiz\_pagine\_021\_030\_024}: D. Le radici hanno argomenti $-\pi/6$, $\pi/2$, $7\pi/6$: nessuna coppia \textbackslash{}`e coniugata.\par\medskip
\noindent\textbf{quiz\_pagine\_021\_030\_026}: E. La costante \textbackslash{}`e $\pi/2$ per $x>0$ e $-\pi/2$ per $x<0$, quindi non \textbackslash{}`e unica su $\mathbb R\setminus\{0\}$.\par\medskip
\noindent\textbf{quiz\_pagine\_021\_030\_030}: E. Poich\textbackslash{}'e $-1\le\cos x\le1$, si ha $1/2\le2^{\cos x}\le2$.\par\medskip
\noindent\textbf{quiz\_pagine\_021\_030\_033}: B. $\int_1^{+\infty}3x^{-3}dx=[-3/(2x^2)]_1^{+\infty}=3/2$.\par\medskip
\noindent\textbf{quiz\_pagine\_021\_030\_034}: D. L\textbackslash{}'integrale \textbackslash{}`e $-\log(\cos\pi/4)=\frac12\log2$; (d) vale il doppio.\par\medskip
\noindent\textbf{quiz\_pagine\_021\_030\_035}: A. $\sqrt[3]{1+3x}=1+x+o(x)$, quindi il rapporto tende a $1/7$.\par\medskip
\noindent\textbf{quiz\_pagine\_031\_040\_001}: E. Spiegazione non ancora disponibile.\par\medskip
\noindent\textbf{quiz\_pagine\_031\_040\_011}: D. Spiegazione non ancora disponibile.\par\medskip
\noindent\textbf{quiz\_pagine\_031\_040\_027}: A. Spiegazione non ancora disponibile.\par\medskip
\noindent\textbf{quiz\_pagine\_031\_040\_028}: A. Spiegazione non ancora disponibile.\par\medskip
\noindent\textbf{quiz\_pagine\_031\_040\_031}: A. Spiegazione non ancora disponibile.\par\medskip
\noindent\textbf{quiz\_pagine\_031\_040\_032}: B. Spiegazione non ancora disponibile.\par\medskip
\noindent\textbf{quiz\_pagine\_031\_040\_033}: B. Spiegazione non ancora disponibile.\par\medskip
\noindent\textbf{quiz\_pagine\_031\_040\_034}: B. Spiegazione non ancora disponibile.\par\medskip
\noindent\textbf{quiz\_pagine\_031\_040\_035}: A. Spiegazione non ancora disponibile.\par\medskip
\noindent\textbf{quiz\_pagine\_031\_040\_036}: A. Spiegazione non ancora disponibile.\par\medskip
\noindent\textbf{quiz\_pagine\_031\_040\_037}: B. Spiegazione non ancora disponibile.\par\medskip
\noindent\textbf{quiz\_pagine\_041\_050\_001}: B. Risposta ufficiale indicata nella tabella della decima simulazione.\par\medskip
\noindent\textbf{quiz\_pagine\_041\_050\_004}: E. Risposta ufficiale indicata nella tabella della decima simulazione.\par\medskip
\noindent\textbf{quiz\_pagine\_041\_050\_016}: A. Risposta ufficiale indicata nella tabella della decima simulazione.\par\medskip
\noindent\textbf{quiz\_pagine\_041\_050\_021}: non verificata / nessuna alternativa corretta. Soluzione non riportata nel PDF; da verificare.\par\medskip
\noindent\textbf{quiz\_pagine\_041\_050\_022}: non verificata / nessuna alternativa corretta. Soluzione non riportata nel PDF; da verificare.\par\medskip
\noindent\textbf{quiz\_pagine\_041\_050\_023}: non verificata / nessuna alternativa corretta. Soluzione non riportata nel PDF; da verificare.\par\medskip
\noindent\textbf{quiz\_pagine\_041\_050\_024}: non verificata / nessuna alternativa corretta. Soluzione non riportata nel PDF; da verificare.\par\medskip
\noindent\textbf{quiz\_pagine\_041\_050\_025}: non verificata / nessuna alternativa corretta. Soluzione non riportata nel PDF; da verificare.\par\medskip
\noindent\textbf{quiz\_pagine\_041\_050\_026}: non verificata / nessuna alternativa corretta. Soluzione non riportata nel PDF; da verificare.\par\medskip
\noindent\textbf{quiz\_pagine\_041\_050\_030}: non verificata / nessuna alternativa corretta. Formula verificata visivamente sulla pagina 8 del PDF.\par\medskip
\noindent\textbf{quiz\_pagine\_041\_050\_031}: non verificata / nessuna alternativa corretta. Soluzione non riportata nel PDF; da verificare.\par\medskip
\noindent\textbf{quiz\_pagine\_051\_060\_004}: C. Una funzione strettamente crescente è iniettiva e la sua inversa sull'immagine è strettamente crescente.\par\medskip
\noindent\textbf{quiz\_pagine\_051\_060\_008}: A. La composizione di due funzioni dispari è dispari.\par\medskip
\noindent\textbf{quiz\_pagine\_061\_070\_001}: E. Le prime quattro funzioni non sono pari; la risposta ufficiale è E.\par\medskip
\noindent\textbf{quiz\_pagine\_061\_070\_021}: B. Il dominio comune è x>1; su tale dominio vale la proprietà della radice del quoziente.\par\medskip
\noindent\textbf{quiz\_pagine\_071\_080\_020}: non verificata / nessuna alternativa corretta. Dato $f(0)=-2$, ogni minimo globale, se esiste, è necessariamente $\le-2$. Senza continuità l'esistenza del minimo non è garantita.\par\medskip
\noindent\textbf{quiz\_pagine\_091\_100\_001}: B. Ponendo $y=\log x+2$, si ottiene $x=e^{y-2}$; quindi $f^{-1}(x)=e^{x-2}$.\par\medskip
\noindent\textbf{quiz\_pagine\_091\_100\_018}: D. $y=\sqrt[7]{x+7}$ implica $y^7=x+7$, dunque $x=y^7-7$.\par\medskip
\noindent\textbf{quiz\_pagine\_101\_110\_001}: E. Ponendo $y=e^x+2$ si ottiene $x=\log(y-2)$, per $y>2$.\par\medskip
\noindent\textbf{quiz\_pagine\_101\_110\_018}: D. La composizione $g\circ f$ è $e^{4\log x}=x^4$, ma il suo dominio è $(0,+\infty)$, non $\mathbb R$.\par\medskip
\noindent\textbf{quiz\_pagine\_111\_120\_021}: E. $\operatorname{dom}f=(-\infty,-3)\cup[1,+\infty)$, $\operatorname{dom}g=[1,+\infty)$. Le funzioni coincidono sul dominio di $g$. Nota: anche b) e c) sono letteralmente vere; il quesito originale non ha risposta unica.\par\medskip
\noindent\textbf{quiz\_pagine\_131\_140\_007}: A. Il logaritmo e il seno sono dispari; $k\cos x$ e\textbackslash{}'e pari e impedisce la disparita\textbackslash{}'a se $k\neq0$.\par\medskip
\noindent\textbf{quiz\_pagine\_131\_140\_011}: A. Per $x<-1$ vale $-3x-1$, per $-1\le x\le0$ vale $1-x$, per $x>0$ vale $1+3x$. La funzione decresce fino a zero e poi cresce; su $(-2,a)$ e\textbackslash{}'e iniettiva se $a\le0$.\par\medskip
\noindent\textbf{quiz\_pagine\_131\_140\_016}: C. La stretta monotonia garantisce l\textbackslash{}'iniettivita\textbackslash{}'a e quindi l\textbackslash{}'esistenza dell\textbackslash{}'inversa sull\textbackslash{}'immagine $f(A)$, non necessariamente su tutto $\mathbb R$.\par\medskip
\noindent\textbf{quiz\_pagine\_131\_140\_017}: B. Occorre $x\neq0$ e $f(x)\neq0$; quest\textbackslash{}'ultima condizione equivale a $|x|\neq1$.\par\medskip
\noindent\textbf{quiz\_pagine\_131\_140\_022}: E. La funzione e\textbackslash{}'e strettamente crescente su $[3,4)$ e decrescente su $(-\infty,3]$; la risposta B e\textbackslash{}'e formulata su un intervallo che include punti fuori dal dominio. La convessita\textbackslash{}'a globale non vale: per $x<3$, $f=\log(4-x)$ e $f\prime\prime<0$. Nessuna opzione e\textbackslash{}'e corretta cosi\textbackslash{}'e scritta: quesito ambiguo.\par\medskip
\noindent\textbf{quiz\_pagine\_151\_160\_005}: C. Una funzione affine $ax+b$ ha in generale inversa $(x-b)/a$, diversa da $ax+b$; per $a=1,b=0$ oppure $a=-1$ coincide.\par\medskip
\noindent\textbf{quiz\_pagine\_151\_160\_020}: B. Se $a>0$ il punto $0$ appartiene a $(-1,a)$ e $f$ assume il minimo $0$; se $a=0$, $0$ e\textbackslash{}'e escluso.\par\medskip
\noindent\textbf{quiz\_pagine\_171\_180\_026}: E. Per $f$ servono $x-1>0$ e $x+1>0$, cioe' $x>1$. Su questo intervallo vale la proprieta' dei logaritmi.\par\medskip
\noindent\textbf{quiz\_pagine\_181\_190\_006}: C. Gli argomenti dei moduli sono simmetrici rispetto al punto medio $-10$. Anche E e' vera: la funzione e' derivabile fuori dai due punti angolosi. Il quesito presenta due risposte corrette.\par\medskip
\noindent\textbf{quiz\_pagine\_211\_220\_032}: B. Il logaritmo richiede argomento positivo: $4-x^2>0$ e $x\neq1$.\par\medskip
\noindent\textbf{quiz\_pagine\_211\_220\_033}: E. Per $x<-1$ il radicando e' $-3x-1$, per $x\geq-1$ e' $1-x$: il dominio e' $(-\infty,1]$ la funzione e' strettamente decrescente su tutto il dominio, quindi iniettiva.\par\medskip
\noindent\textbf{quiz\_pagine\_231\_240\_002}: A. $x^2+2\geq2>1$, dunque $\ln(x^2+2)>0$ per ogni $x$ e non ci sono zeri.\par\medskip
\noindent\textbf{quiz\_pagine\_231\_240\_014}: C. Il denominatore richiede $\ln(x^2-1)>0$, cioe' $x^2-1>1$ e $|x|>\sqrt2$.\par\medskip
\noindent\textbf{quiz\_pagine\_231\_240\_023}: B. Per $x_1<x_2$ vale $f(x_1)<f(x_2)$, quindi valori uguali non possono provenire da argomenti distinti. La derivata puo' anche non esistere.\par\medskip
\noindent\textbf{quiz\_pagine\_241\_250\_005}: A. $2\ln\sqrt{x^2+1}=\ln(x^2+1)$; sommando $-\ln(x^2+1)$ si ottiene zero.\par\medskip
\noindent\textbf{quiz\_pagine\_241\_250\_009}: D. $f(x/3)$ dilata orizzontalmente di fattore 3; $+2$ trasla verso l'alto di 2.\par\medskip
\noindent\textbf{quiz\_pagine\_241\_250\_011}: E. Per $x\geq5$, $f(x)=-x-5$ e' strettamente decrescente, dunque invertibile sulla restrizione. In $x=-5$ le derivate laterali sono 0 e -1.\par\medskip
\noindent\textbf{quiz\_pagine\_241\_250\_012}: A. $f=-8$ per $x<-4$, $f=2x$ per $-4\leq x\leq4$, $f=8$ per $x>4$: l'immagine e' $[-8,8]$, non solo nonnegativa.\par\medskip
\noindent\textbf{quiz\_pagine\_251\_260\_009}: C. Per $a\ne0$, $f^{-1}(x)=(x-b)/a$. L'uguaglianza richiede $a^2=1$ e $b(a+1)=0$: identita' oppure riflessione $f(x)=-x+b$.\par\medskip
\noindent\textbf{quiz\_pagine\_291\_300\_037}: E. La prima e' definita soltanto per $x>1$; in tale intervallo la proprieta' dei logaritmi rende le due espressioni uguali.\par\medskip
\noindent\textbf{quiz\_pagine\_321\_330\_002}: E. Il radicando richiede $|x|\leq1$, ma il logaritmo impone argomento strettamente positivo: $|x|<1$.\par\medskip
\noindent\textbf{quiz\_pagine\_321\_330\_036}: E. $e^{2x}-1=0$ solo per $x=0$. La singolarit\textbackslash{}`a \textbackslash{}`e eliminabile, ma il punto resta escluso dal dominio della formula.\par\medskip
\noindent\textbf{quiz\_pagine\_321\_330\_038}: C. Il dominio \textbackslash{}`e $[-1,1]$ e l\textbackslash{}textquotesingle immagine effettiva \textbackslash{}`e $[\sqrt2,2]$ (estremi inclusi). La chiave ufficiale C riporta erroneamente intervallo aperto: nessuna alternativa \textbackslash{}`e rigorosamente vera.\par\medskip
\noindent\textbf{quiz\_pagine\_321\_330\_040}: C. Per scrivere la differenza dei logaritmi occorrono simultaneamente $3-x>0$ e $3+x>0$, cio\textbackslash{}`e $-3<x<3$.\par\medskip
\noindent\textbf{quiz\_pagine\_331\_340\_004}: E. Essendo $-1\leq\cos x\leq1$, risulta $1/2\leq2^{\cos x}\leq2$.\par\medskip
\noindent\textbf{quiz\_pagine\_331\_340\_011}: E. La composizione $g(g(x))=\cos(\cos x)$ \textbackslash{}`e definita per ogni reale; le altre richiedono una radice quadrata.\par\medskip
\noindent\textbf{quiz\_pagine\_341\_350\_001}: E. $g$ e definita su tutto $\mathbb R$, quindi anche $g\circ g$. Tutte le altre espressioni coinvolgono $\log x$ o il logaritmo di un valore che puo essere nullo.\par\medskip
\noindent\textbf{quiz\_pagine\_341\_350\_008}: E. $\arctan x$ e definita su tutto $\mathbb R$. Il valore assoluto conserva il dominio ma introduce una cuspide in zero.\par\medskip
\noindent\textbf{quiz\_pagine\_341\_350\_016}: B. $g$ e definita solo quando $x+1>0$ e $x-1>0$, cioe $x>1$; qui vale la proprieta dei logaritmi. $f$ e definita anche per $x<-1$, ma $g$ no.\par\medskip
\noindent\textbf{quiz\_pagine\_341\_350\_019}: E. $|t|=t$ se e solo se $t\geq0$. Quindi $\log|x|\geq0$ equivale a $|x|\geq1$.\par\medskip
\noindent\textbf{quiz\_pagine\_341\_350\_031}: A. Per definizione, una funzione a valori reali ha immagine contenuta in $\mathbb R$. La continuita su un intervallo aperto non garantisce estremi assoluti, zeri o derivabilita.\par\medskip
\noindent\textbf{quiz\_pagine\_351\_360\_001}: E. $g$ e definita su tutto $\mathbb R$, quindi anche $g\circ g$. Tutte le altre espressioni coinvolgono $\log x$ o il logaritmo di un valore che puo essere nullo.\par\medskip
\noindent\textbf{quiz\_pagine\_351\_360\_008}: E. $\arctan x$ e definita su tutto $\mathbb R$. Il valore assoluto conserva il dominio ma introduce una cuspide in zero.\par\medskip
\noindent\textbf{quiz\_pagine\_351\_360\_016}: B. $g$ e definita solo quando $x+1>0$ e $x-1>0$, cioe $x>1$; qui vale la proprieta dei logaritmi. $f$ e definita anche per $x<-1$, ma $g$ no.\par\medskip
\noindent\textbf{quiz\_pagine\_351\_360\_019}: E. $|t|=t$ se e solo se $t\geq0$. Quindi $\log|x|\geq0$ equivale a $|x|\geq1$.\par\medskip
\noindent\textbf{quiz\_pagine\_351\_360\_031}: A. Per definizione, una funzione a valori reali ha immagine contenuta in $\mathbb R$. La continuita su un intervallo aperto non garantisce estremi assoluti, zeri o derivabilita.\par\medskip
\noindent\textbf{quiz\_pagine\_361\_370\_004}: C. Una funzione strettamente crescente e iniettiva e ammette inversa sulla sua immagine; tale inversa e strettamente crescente. Non e necessariamente suriettiva su $\mathbb R$.\par\medskip
\noindent\textbf{quiz\_pagine\_361\_370\_008}: A. $f(g(-x))=f(-g(x))=-f(g(x))$: la composizione di due funzioni dispari e dispari.\par\medskip
\noindent\textbf{quiz\_pagine\_361\_370\_021}: E. Le prime tre funzioni non sono pari (la seconda e la terza sono dispari); la quarta contiene una parte pari e una dispari. Nessuna e pari.\par\medskip
\noindent\textbf{quiz\_pagine\_371\_380\_011}: B. La seconda funzione richiede $x+1\geq0$ e $x-1>0$, dunque $x>1$. Sul dominio comune vale la proprieta della radice del quoziente.\par\medskip
\noindent\textbf{quiz\_pagine\_391\_400\_013}: B. Da $y=\log x+2$ si ricava $x=e^{y-2}$, dunque $f^{-1}(y)=e^{y-2}$.\par\medskip
\noindent\textbf{quiz\_pagine\_411\_420\_001}: E. Ponendo $y=e^x+2$, si ricava $x=\log(y-2)$, con $y>2$.\par\medskip
\noindent\textbf{quiz\_pagine\_411\_420\_018}: D. $g(f(x))=e^{4\log x}=x^4$ soltanto per $x>0$: il dominio non \textbackslash{}`e $\mathbb R$.\par\medskip
\noindent\textbf{quiz\_pagine\_421\_430\_004}: E. $\operatorname{dom}f=[1,+\infty)$ e $\operatorname{dom}g=(-\infty,-3)\cup[1,+\infty)$. Le due espressioni coincidono su $[1,+\infty)$. Anche B \textbackslash{}`e vera come affermazione non massimale: ambiguit\textbackslash{}`a nella fonte.\par\medskip
\noindent\textbf{quiz\_pagine\_441\_450\_003}: B. $\tan(-x)=-\tan x$ e $\cos^2$ \textbackslash{}`e pari, quindi $f$ \textbackslash{}`e pari e non dispari.\par\medskip
\noindent\textbf{quiz\_pagine\_441\_450\_027}: A. Il logaritmo del rapporto \textbackslash{}`e dispari, come $\sin x$; $k\cos x$ \textbackslash{}`e pari e deve scomparire per avere disparit\textbackslash{}`a.\par\medskip
\noindent\textbf{quiz\_pagine\_501\_510\_001}: B. Il numeratore ha il segno di $4-x^2$, mentre il denominatore $(x-1)^2$ \textbackslash{}`e positivo per $x\ne1$. Il logaritmo richiede $-2<x<2$, con $x\ne1$.\par\medskip
\noindent\textbf{quiz\_pagine\_501\_510\_002}: E. Il radicando \textbackslash{}`e $-3x-1$ per $x<-1$ e $1-x$ per $x\ge-1$. Nel dominio $(-\infty,1]$ la funzione \textbackslash{}`e strettamente decrescente, quindi iniettiva. Il testo scrive $f:\mathbb R\to\mathbb R$, ma la formula non \textbackslash{}`e definita per $x>1$: incongruenza nell\textbackslash{}'indicazione del dominio.\par\medskip
\noindent\textbf{quiz\_pagine\_501\_510\_021}: D. Per definire il logaritmo occorre $\sqrt{x-2}>3$, cio\textbackslash{}`e $x>11$.\par\medskip
\noindent\textbf{quiz\_pagine\_501\_510\_022}: B. Su $[4,+\infty)$ si ha $\log(x-3)\ge0$ e la funzione \textbackslash{}`e crescente, quindi iniettiva. Negli intervalli che attraversano $x=4$ il valore assoluto rompe la monotonia.\par\medskip
\noindent\textbf{quiz\_pagine\_501\_510\_033}: D. La continuit\textbackslash{}`a garantisce tutti i valori intermedi e la decrescenza individua il massimo in $0$ e il minimo in $3$.\par\medskip
\noindent\textbf{quiz\_pagine\_511\_520\_001}: C. Serve $x>1$ e $\log_{1/2}(x-1)\ge-2$. Essendo la base minore di $1$, si ottiene $x-1\le4$, quindi $1<x\le5$.\par\medskip
\noindent\textbf{quiz\_pagine\_511\_520\_022}: B. Per $x\to-\infty$, $e^x\to0$ e $\cos x$ oscilla fra $-1$ e $1$: per continuit\textbackslash{}`a si trovano infiniti zeri distinti. Non \textbackslash{}`e monotona n'e periodica.\par\medskip
\noindent\textbf{quiz\_pagine\_511\_520\_032}: B. Il valore $1$ \textbackslash{}`e assunto da $\sqrt{x}$ soltanto in $x=1$, dove la funzione vale invece $\alpha$. Per coprire $[0,2]$ occorre e basta $\alpha=1$.\par\medskip
\subsubsection*{1.5 -- Insiemi numerici}
\noindent\textbf{quiz\_pagine\_001\_010\_002}: A. Spiegazione non ancora disponibile.\par\medskip
\noindent\textbf{quiz\_pagine\_001\_010\_018}: D. Spiegazione non ancora disponibile.\par\medskip
\noindent\textbf{quiz\_pagine\_001\_010\_019}: E. Spiegazione non ancora disponibile.\par\medskip
\noindent\textbf{quiz\_pagine\_001\_010\_021}: C. Spiegazione non ancora disponibile.\par\medskip
\noindent\textbf{quiz\_pagine\_011\_020\_001}: C. $\cosh x>4$ equivale a $|x|>\operatorname{arcosh}4$: il dominio \textbackslash{}`e unione di due intervalli.\par\medskip
\noindent\textbf{quiz\_pagine\_011\_020\_003}: A. Per $z=x\in\mathbb R$, $(x-1)^2+1<9$, cio\textbackslash{}`e $1-\sqrt8<x<1+\sqrt8$.\par\medskip
\noindent\textbf{quiz\_pagine\_011\_020\_009}: B. Poich\textbackslash{}'e $f(0)>0$, vicino a zero $|f|=f$ e quindi $g\prime(0)=-4$.\par\medskip
\noindent\textbf{quiz\_pagine\_011\_020\_013}: D. $f\prime(x)=5x^4+\cosh x>0$, quindi $f$ \textbackslash{}`e strettamente crescente e il valore inferiore a $f(10)$ \textbackslash{}`e assunto prima di $10$.\par\medskip
\noindent\textbf{quiz\_pagine\_011\_020\_014}: A. $g\prime=3f^2f\prime$ si annulla in $0$, in $1$ e in un punto interno per Rolle.\par\medskip
\noindent\textbf{quiz\_pagine\_011\_020\_036}: E. $\log(1+x^2)\sim x^2$ e $x\sin x\sim x^2$; per $\alpha>2$ il limite \textbackslash{}`e $1$, per $\alpha=2$ \textbackslash{}`e $2$, altrimenti non \textbackslash{}`e $0$.\par\medskip
\noindent\textbf{quiz\_pagine\_021\_030\_002}: C. $f\prime$ ha segno di $-6x(x+1)-1$; \textbackslash{}`e positivo tra le radici $(-3\pm\sqrt3)/6$. Attenzione: l\textbackslash{}'intervallo attraversa $x=-1$? No, entrambe le radici sono maggiori di $-1$.\par\medskip
\noindent\textbf{quiz\_pagine\_021\_030\_005}: C. Ogni complesso non nullo ha esattamente quattro radici quarte distinte. La formula degli argomenti si intende modulo $2\pi$.\par\medskip
\noindent\textbf{quiz\_pagine\_021\_030\_013}: E. Moltiplicando per $1-i$ si ottiene $z=(1-i)^2/2=-i$.\par\medskip
\noindent\textbf{quiz\_pagine\_021\_030\_015}: C. Applicando Rolle a due punti con lo stesso valore, oppure usando un estremo interno quando $g$ non \textbackslash{}`e costante, si trova $g\prime(x_0)=0$.\par\medskip
\noindent\textbf{quiz\_pagine\_021\_030\_027}: B. Per Rolle applicato a $f-g$ esiste $c$ con $f\prime(c)=g\prime(c)$; basta prendere $c_1=c_2=c$.\par\medskip
\noindent\textbf{quiz\_pagine\_021\_030\_031}: C. Per definizione $[f(x)^3-f(0)^3]/(x-0)=f(x)^3/x$.\par\medskip
\noindent\textbf{quiz\_pagine\_021\_030\_032}: E. $(2-3i)^2=-5-12i$ e $i^7=-i$, quindi $(-i)(-5-12i)=-12+5i$. L\textbackslash{}'opzione corretta \textbackslash{}`e (e), non (a).\par\medskip
\noindent\textbf{quiz\_pagine\_031\_040\_012}: C. Spiegazione non ancora disponibile.\par\medskip
\noindent\textbf{quiz\_pagine\_031\_040\_014}: D. Spiegazione non ancora disponibile.\par\medskip
\noindent\textbf{quiz\_pagine\_031\_040\_015}: D. Spiegazione non ancora disponibile.\par\medskip
\noindent\textbf{quiz\_pagine\_031\_040\_016}: E. Spiegazione non ancora disponibile.\par\medskip
\noindent\textbf{quiz\_pagine\_031\_040\_017}: A. Spiegazione non ancora disponibile.\par\medskip
\noindent\textbf{quiz\_pagine\_031\_040\_018}: C. Spiegazione non ancora disponibile.\par\medskip
\noindent\textbf{quiz\_pagine\_031\_040\_029}: D. Spiegazione non ancora disponibile.\par\medskip
\noindent\textbf{quiz\_pagine\_031\_040\_030}: B. Spiegazione non ancora disponibile.\par\medskip
\noindent\textbf{quiz\_pagine\_041\_050\_005}: A. Risposta ufficiale indicata nella tabella della decima simulazione.\par\medskip
\noindent\textbf{quiz\_pagine\_061\_070\_004}: D. Poiché z coniugato(z)=|z|\textsuperscript{2}=1, il luogo è la circonferenza unitaria.\par\medskip
\noindent\textbf{quiz\_pagine\_061\_070\_025}: A. La parte immaginaria di (a+ib)(3+i) è a+3b; si annulla per a=-3b.\par\medskip
\noindent\textbf{quiz\_pagine\_071\_080\_017}: E. Le radici cubiche di $-27$ sono $-3$ e $3/2\pm(3\sqrt3/2)i$; le ultime due sono coniugate.\par\medskip
\noindent\textbf{quiz\_pagine\_081\_090\_013}: A. $\overline z=-1-2i$, quindi $\overline z^{\,2}+3=4i$ e il suo modulo è $4$.\par\medskip
\noindent\textbf{quiz\_pagine\_081\_090\_014}: C. $\frac{1-2i}{3+4i}=\frac{(1-2i)(3-4i)}{25}=\frac{-5-10i}{25}=-\frac15-\frac25i$.\par\medskip
\noindent\textbf{quiz\_pagine\_081\_090\_017}: E. Fra le sei radici complesse, le due reali sono $z=\pm\sqrt[6]{11}$.\par\medskip
\noindent\textbf{quiz\_pagine\_091\_100\_004}: B. Il massimo è $5$; $1+1/(3n)>1$ e tende a $1$, che non appartiene ad $A$. Quindi non c'è minimo.\par\medskip
\noindent\textbf{quiz\_pagine\_091\_100\_005}: D. Per esempio $A=(0,1)$ e $B=[0,1]$ sono distinti, con $A\subset B$, ma hanno entrambi $\inf=0$ e $\sup=1$.\par\medskip
\noindent\textbf{quiz\_pagine\_091\_100\_026}: E. Ogni minorante di $B$ è minorante di $A$; quindi, per insiemi non vuoti e inferiormente limitati, $\inf(A)\geq\inf(B)$.\par\medskip
\noindent\textbf{quiz\_pagine\_101\_110\_013}: D. Le radici di $i$ sono $\pm(1+i)/\sqrt2$, cioè opposte e non coniugate. Entrambe hanno modulo $1$.\par\medskip
\noindent\textbf{quiz\_pagine\_111\_120\_002}: D. Il modulo di $-16i$ e' $16$ e un argomento e' $-\pi/2$. Le radici quarte hanno modulo $2$ e argomenti $-\pi/8+k\pi/2$.\par\medskip
\noindent\textbf{quiz\_pagine\_111\_120\_019}: D. $z^2=3+4i$, quindi $z^2+3i=3+7i$ e il modulo e' $\sqrt{3^2+7^2}=\sqrt{58}$.\par\medskip
\noindent\textbf{quiz\_pagine\_121\_130\_013}: E. La parte immaginaria di $z^2$ e\textbackslash{}'e $2(k-3)(2-k)$, nulla esattamente per $k=3$ oppure $k=2$.\par\medskip
\noindent\textbf{quiz\_pagine\_131\_140\_005}: A. Se $z=iy$, $w=i(y-1/y)$ e\textbackslash{}'e reale solo per $y=\pm1$. Attenzione: anche ogni $z$ reale non nullo rende $w$ reale; il testo chiede i valori immaginari puri.\par\medskip
\noindent\textbf{quiz\_pagine\_151\_160\_006}: A. $\arctan(-n)\to-\pi/2$ senza raggiungerlo; $5$ e\textbackslash{}'e il massimo. Anche (b) e\textbackslash{}'e matematicamente vera, perche\textbackslash{}'e il massimo coincide con il supremo.\par\medskip
\noindent\textbf{quiz\_pagine\_151\_160\_007}: B. $\cosh x\geq1$ e non e\textbackslash{}'e limitato superiormente; aggiungendo $-9$ il minimo diventa $-9$.\par\medskip
\noindent\textbf{quiz\_pagine\_151\_160\_019}: D. $0\leq n/(n+5)<1$ e $5\in A$, dunque $A$ e\textbackslash{}'e limitato, con massimo $5$.\par\medskip
\noindent\textbf{quiz\_pagine\_171\_180\_010}: B. $|w|=\sqrt2/5$. Le alternative (b) e (d) risultano identiche nella trascrizione del PDF: risposta non univoca nel testo originale.\par\medskip
\noindent\textbf{quiz\_pagine\_171\_180\_021}: C. $(2+i)^2-(2+i)=(3+4i)-(2+i)=1+3i$.\par\medskip
\noindent\textbf{quiz\_pagine\_181\_190\_008}: A. Per la formula di Eulero $e^{it}=\cos t+i\sin t$; per $t=\pi/4$ le due componenti coincidono. Per $t=k\pi$ il valore puo' essere $-1$.\par\medskip
\noindent\textbf{quiz\_pagine\_211\_220\_034}: C. La condizione $|\bar z-9/z|=0$ equivale a $|z|=3$, ma esistono anche le tre radici di $z^3=8i$, di modulo $2$. Le uniche soluzioni reali sono $z=\pm3$ sul cerchio, quindi C e' vera.\par\medskip
\noindent\textbf{quiz\_pagine\_231\_240\_007}: C. Occorre $x^2-4\geq0$ e $x^2-4\leq1$, ossia $4\leq x^2\leq5$.\par\medskip
\noindent\textbf{quiz\_pagine\_231\_240\_009}: B. Da $-(2x+3)<x-3<2x+3$ segue $x>0$. Nel PDF la parentesi finale $+\infty]$ e' impropria; l'insieme corretto e' $(0,+\infty)$.\par\medskip
\noindent\textbf{quiz\_pagine\_241\_250\_002}: B. Il dominio e' $(-1,0)\cup(4,+\infty)$. La disuguaglianza equivale a $(x+1)^2<x^2-4x$, ossia $6x+1<0$: intersezione $(-1,-1/6)$.\par\medskip
\noindent\textbf{quiz\_pagine\_241\_250\_015}: D. Moltiplicando per $3+i$: $w=(14+8i)/10=7/5+4i/5$.\par\medskip
\noindent\textbf{quiz\_pagine\_251\_260\_005}: B. $|w|=|5+i|/|3-i|=\sqrt{26}/\sqrt{10}=\sqrt{65}/5$.\par\medskip
\noindent\textbf{quiz\_pagine\_271\_280\_003}: B. Per coefficienti reali anche $2-i$ e' radice: con $0$ si hanno almeno tre radici distinte.\par\medskip
\noindent\textbf{quiz\_pagine\_281\_290\_009}: D. Le potenze di $i$ sono periodiche di periodo 4: $i^{k+4}=i^k$.\par\medskip
\noindent\textbf{quiz\_pagine\_281\_290\_026}: C. $\overline z=1-2i$; moltiplicando per $3-4i$ si ottiene $(-5-10i)/25=-1/5-(2/5)i$.\par\medskip
\noindent\textbf{quiz\_pagine\_281\_290\_028}: A. $\overline z=-1-2i$, $\overline z^2=-3+4i$ e $|\overline z^2+3|=|4i|=4$.\par\medskip
\noindent\textbf{quiz\_pagine\_291\_300\_002}: A. $z^2=35-12i$ e $\overline z=6+i$, dunque $w=29-13i$.\par\medskip
\noindent\textbf{quiz\_pagine\_291\_300\_018}: D. $w=(1+3i)(4+2i)/20=-1/10+7i/10$, percio' $\Im w=7/10$.\par\medskip
\noindent\textbf{quiz\_pagine\_301\_310\_023}: A. Il disco aperto di centro $1-i$ e raggio $3$ interseca l'asse reale in un intervallo aperto non vuoto, con infinite soluzioni reali.\par\medskip
\noindent\textbf{quiz\_pagine\_311\_320\_022}: E. Moltiplicando per $3+i$ si ottiene $w=(14+8i)/10=7/5+4i/5$. La domanda scrive $\operatorname{Re}(z)$ anziche $\operatorname{Re}(w)$: refuso originale.\par\medskip
\noindent\textbf{quiz\_pagine\_311\_320\_026}: B. $|w|=\sqrt{5^2+1^2}/\sqrt{3^2+(-1)^2}=\sqrt{26/10}=\sqrt{65}/5$.\par\medskip
\noindent\textbf{quiz\_pagine\_321\_330\_001}: C. Ogni $z\neq0$ ha esattamente quattro radici quarte distinte. L\textbackslash{}textquotesingle uguaglianza degli argomenti in A richiede la convenzione modulo $2\pi$.\par\medskip
\noindent\textbf{quiz\_pagine\_321\_330\_003}: E. La derivata $49x^{48}+4$ \textbackslash{}`e positiva: esiste una sola radice reale, positiva. Le altre $48$ sono in coppie coniugate, ma non formano necessariamente un poligono regolare su una circonferenza.\par\medskip
\noindent\textbf{quiz\_pagine\_321\_330\_006}: D. $|z|z=\sqrt{13}(2+3i)$ ha parte immaginaria non nulla; le altre propriet\textbackslash{}`a sono vere.\par\medskip
\noindent\textbf{quiz\_pagine\_321\_330\_007}: E. La base \textbackslash{}`e $e^{i\pi/3}$; la nona potenza \textbackslash{}`e $e^{3\pi i}=-1$.\par\medskip
\noindent\textbf{quiz\_pagine\_321\_330\_009}: E. Moltiplicando per $1-i$ si ottiene $z=(1-i)^2/2=-i$.\par\medskip
\noindent\textbf{quiz\_pagine\_321\_330\_020}: D. Le radici hanno argomenti $-\pi/6,\pi/2,7\pi/6$: non c\textbackslash{}textquotesingle \textbackslash{}`e una coppia di radici coniugate.\par\medskip
\noindent\textbf{quiz\_pagine\_321\_330\_023}: E. Il modulo di un reale $k$ \textbackslash{}`e $|k|$, anche se $k$ \textbackslash{}`e negativo.\par\medskip
\noindent\textbf{quiz\_pagine\_321\_330\_025}: C. $|z^{-2}|=|z|^{-2}=(\sqrt{3^2+4^2})^{-2}=1/25$.\par\medskip
\noindent\textbf{quiz\_pagine\_321\_330\_026}: A. Scrivendo $z=1+iy$ si ha $1+y^2=16$, quindi $y=\pm\sqrt{15}$: due soluzioni.\par\medskip
\noindent\textbf{quiz\_pagine\_321\_330\_033}: D. $|z|=2$ e l\textbackslash{}textquotesingle argomento principale \textbackslash{}`e $\pi/6$.\par\medskip
\noindent\textbf{quiz\_pagine\_321\_330\_041}: C. $|z|^2$ \textbackslash{}`e reale positivo, quindi $z$ \textbackslash{}`e immaginario puro positivo: $z=it$, $t^3=8$, $t=2$.\par\medskip
\noindent\textbf{quiz\_pagine\_331\_340\_006}: E. $i^7=-i$ e $(2-3i)^2=-5-12i$; dunque $z=-12+5i$.\par\medskip
\noindent\textbf{quiz\_pagine\_341\_350\_009}: A. $u=(a+ib)(1-i)/2=(a+b)/2+i(b-a)/2$ e reale se $a=b$.\par\medskip
\noindent\textbf{quiz\_pagine\_341\_350\_013}: E. $i^3=-i$, quindi $z=(1+i)/(5-i)=(4+6i)/26=2/13+3i/13$.\par\medskip
\noindent\textbf{quiz\_pagine\_341\_350\_020}: A. $u=(a+ib)(3-i)/10=(3a+b)/10+i(3b-a)/10$. La parte immaginaria e zero se $a=3b$.\par\medskip
\noindent\textbf{quiz\_pagine\_351\_360\_009}: A. $u=(a+ib)(1-i)/2=(a+b)/2+i(b-a)/2$ e reale se $a=b$.\par\medskip
\noindent\textbf{quiz\_pagine\_351\_360\_013}: E. $i^3=-i$, quindi $z=(1+i)/(5-i)=(4+6i)/26=2/13+3i/13$.\par\medskip
\noindent\textbf{quiz\_pagine\_351\_360\_020}: A. $u=(a+ib)(3-i)/10=(3a+b)/10+i(3b-a)/10$. La parte immaginaria e zero se $a=3b$.\par\medskip
\noindent\textbf{quiz\_pagine\_361\_370\_024}: D. $z\overline z=|z|^2$ e reale: $|z|^2=1$ descrive la circonferenza unitaria, non il disco.\par\medskip
\noindent\textbf{quiz\_pagine\_371\_380\_015}: A. $(a+ib)(3+i)=(3a-b)+i(a+3b)$; la parte immaginaria si annulla per $a=-3b$.\par\medskip
\noindent\textbf{quiz\_pagine\_381\_390\_002}: E. Le radici sono $-3$ e $\frac32\pm\frac{3\sqrt3}{2}i$, quindi due sono coniugate. La forma esponenziale corretta include il modulo $3$ e angoli $(\pi+2k\pi)/3$.\par\medskip
\noindent\textbf{quiz\_pagine\_381\_390\_018}: A. $z^2=(-1+2i)^2=-3-4i$, dunque $z^2+3=-4i$ e $|-4i|=4$.\par\medskip
\noindent\textbf{quiz\_pagine\_381\_390\_019}: C. $\overline z/w=(1-2i)(3-4i)/25=(-5-10i)/25=-1/5-2i/5$.\par\medskip
\noindent\textbf{quiz\_pagine\_391\_400\_002}: E. Le sei radici sono $11^{1/6}e^{ik\pi/3}$, $k=0,\ldots,5$; quelle reali sono $\pm11^{1/6}$, quindi due.\par\medskip
\noindent\textbf{quiz\_pagine\_401\_410\_001}: C. Le radici quadrate di $i$ sono opposte: $z_2=-z_1$, dunque $z_1/z_2=-1\in\mathbb R$.\par\medskip
\noindent\textbf{quiz\_pagine\_401\_410\_017}: B. Poich\textbackslash{}`e i coefficienti sono reali, anche $2-i$ \textbackslash{}`e radice. Le radici distinte sono almeno tre: $0$, $2+i$, $2-i$.\par\medskip
\noindent\textbf{quiz\_pagine\_411\_420\_013}: D. Le due radici sono opposte, $z_2=-z_1$, non coniugate; entrambe hanno modulo $1$.\par\medskip
\noindent\textbf{quiz\_pagine\_411\_420\_022}: D. $-16i=16e^{-i\pi/2}$, perci\textbackslash{}`o le radici hanno modulo $2$ e argomenti $-\pi/8+k\pi/2$.\par\medskip
\noindent\textbf{quiz\_pagine\_411\_420\_037}: C. Il massimo \textbackslash{}`e $0$ (per $n=0$), il minimo \textbackslash{}`e $-3$ (elemento aggiunto); la successione tende a $-1$.\par\medskip
\noindent\textbf{quiz\_pagine\_421\_430\_002}: D. $z^2=3+4i$, dunque $z^2+3i=3+7i$ e il modulo vale $\sqrt{9+49}=\sqrt{58}$.\par\medskip
\noindent\textbf{quiz\_pagine\_421\_430\_008}: D. Un estremo inferiore finito garantisce soltanto limitatezza inferiore, non superiore: ad esempio $A=[0,+\infty)$.\par\medskip
\noindent\textbf{quiz\_pagine\_421\_430\_026}: A. Le quattro radici di $z^4=-4$ sono $\pm1\pm i$, di cui due con parte reale negativa. Le tre radici di $z^3=8$ sono $2$ e $-1\pm i\sqrt3$, di cui altre due negative: totale quattro.\par\medskip
\noindent\textbf{quiz\_pagine\_431\_440\_002}: E. $|z|=\sqrt{2+2}/\sqrt{3+1}=1$.\par\medskip
\noindent\textbf{quiz\_pagine\_431\_440\_009}: B. Moltiplicare per $-1$ ruota di $\pi$ l\textbackslash{}'argomento: B \textbackslash{}`e falsa. Le uguaglianze con $\operatorname{Arg}$ principale richiedono attenzione sui tagli del ramo (per esempio l\textbackslash{}'asse reale negativo).\par\medskip
\noindent\textbf{quiz\_pagine\_431\_440\_028}: E. La parte immaginaria di $z^2$ \textbackslash{}`e $2(k-3)(2-k)$, nulla per $k=2$ o $k=3$.\par\medskip
\noindent\textbf{quiz\_pagine\_441\_450\_007}: B. $i^5=i$; $iz=-2+4i$ e $i\overline z=2+4i$, perci\textbackslash{}`o le parti immaginarie coincidono. NOTA: nel PDF le opzioni C e D sono tipograficamente molto simili; verificare sul documento originale.\par\medskip
\noindent\textbf{quiz\_pagine\_441\_450\_025}: A. Per $z=iy$, $w=i(y-1/y)$ \textbackslash{}`e reale esattamente per $y=\pm1$. NOTA: anche ogni $z$ reale non nullo rende $w$ reale; A resta vera come condizione specifica sui puri immaginari.\par\medskip
\noindent\textbf{quiz\_pagine\_471\_480\_005}: D. Tutti i rapporti $n/(n+5)$ appartengono a $[0,1)$ e l'insieme contiene anche $5$. Quindi \textbackslash{}`e limitato e ha massimo $5$.\par\medskip
\noindent\textbf{quiz\_pagine\_471\_480\_013}: C. $|w|=\sqrt{1+9}/\sqrt{16+4}=\sqrt{10/20}=1/\sqrt2=\sqrt2/2$. Quindi la risposta corretta \textbackslash{}`e C.\par\medskip
\noindent\textbf{quiz\_pagine\_481\_490\_006}: E. L\textbackslash{}'estremo superiore pu\textbackslash{}`o non appartenere a $B$, per esempio $B=(0,3)$, che non ha massimo.\par\medskip
\noindent\textbf{quiz\_pagine\_481\_490\_007}: B. L\textbackslash{}'infimo \textbackslash{}`e il maggiore dei minoranti, ma non deve appartenere all\textbackslash{}'insieme. Per esempio $A=(5,6)$.\par\medskip
\noindent\textbf{quiz\_pagine\_491\_500\_005}: C. ATTENZIONE: la tabella ufficiale riporta E, ma $z^2=(2+i)^2=3+4i$, dunque $z^2-z=1+3i$: \textbackslash{}`e corretta C. Anche la spiegazione ufficiale contiene un errore di calcolo.\par\medskip
\noindent\textbf{quiz\_pagine\_501\_510\_003}: C. La condizione $\overline z=9/z$ equivale a $|z|=3$, che fornisce una circonferenza di soluzioni. Anche $z^3=8i$ fornisce tre radici di modulo $2$. Le uniche soluzioni reali sono $z=\pm3$.\par\medskip
\noindent\textbf{quiz\_pagine\_501\_510\_023}: A. $z^2=-3-4i$ e $\overline z=1+2i$; somma $-2-2i$, modulo $2\sqrt2$.\par\medskip
\noindent\textbf{quiz\_pagine\_511\_520\_002}: C. In $A$ la parte reale \textbackslash{}`e compresa tra $-1$ e $1$; in $B$ \textbackslash{}`e maggiore di $5$. I due insiemi sono disgiunti.\par\medskip
\noindent\textbf{quiz\_pagine\_511\_520\_023}: B. Per coefficienti reali anche $2-i$ \textbackslash{}`e radice. Le radici distinte $0$, $2+i$, $2-i$ impongono grado almeno $3$.\par\medskip
\subsection*{Capitolo 2}
\subsubsection*{2.1 -- Successioni reali e loro proprietà}
\noindent\textbf{quiz\_pagine\_031\_040\_013}: C. Spiegazione non ancora disponibile.\par\medskip
\noindent\textbf{quiz\_pagine\_031\_040\_020}: E. Spiegazione non ancora disponibile.\par\medskip
\noindent\textbf{quiz\_pagine\_031\_040\_021}: A. Spiegazione non ancora disponibile.\par\medskip
\noindent\textbf{quiz\_pagine\_031\_040\_022}: E. Spiegazione non ancora disponibile.\par\medskip
\noindent\textbf{quiz\_pagine\_031\_040\_039}: A. Spiegazione non ancora disponibile.\par\medskip
\noindent\textbf{quiz\_pagine\_031\_040\_040}: B. Spiegazione non ancora disponibile.\par\medskip
\noindent\textbf{quiz\_pagine\_031\_040\_042}: B. Spiegazione non ancora disponibile.\par\medskip
\noindent\textbf{quiz\_pagine\_141\_150\_002}: E. Un minorante reale $m$ esiste per ipotesi; scegliendo $k<\min\{m,0\}$ si ottiene $a_n>k$ per ogni $n$.\par\medskip
\noindent\textbf{quiz\_pagine\_141\_150\_016}: A. Essendo $a_n\to l\neq0$, la successione $1/a_n$ e\textbackslash{}'e limitata. Il prodotto con $b_n$ limitata e\textbackslash{}'e limitato.\par\medskip
\noindent\textbf{quiz\_pagine\_141\_150\_017}: C. La negazione della limitatezza e\textbackslash{}'e $\forall M>0\ \exists n:\ |a_n|>M$.\par\medskip
\noindent\textbf{quiz\_pagine\_181\_190\_026}: A. Per $a,b>0$ si ha $a^{\log b}=e^{(\log a)(\log b)}=b^{\log a}$. I due termini coincidono per ogni $n$, dunque $a_n=0$.\par\medskip
\noindent\textbf{quiz\_pagine\_191\_200\_003}: B. $i^5=i$; $iz=-2+4i$ e $i\overline z=2+4i$ hanno la stessa parte immaginaria. ATTENZIONE: anche D risulta vera letteralmente, poiche' $i^5=i$; la domanda originale e' ambigua.\par\medskip
\noindent\textbf{quiz\_pagine\_191\_200\_004}: D. Considerando gli argomenti principali delle tre radici cubiche, due risultano negativi per $-\pi<\theta<0$. Gli estremi richiedono una convenzione sugli argomenti.\par\medskip
\noindent\textbf{quiz\_pagine\_271\_280\_008}: A. Il numeratore e' $O(n)$, mentre il denominatore ha modulo asintotico $2^n$, per cui il rapporto tende a zero.\par\medskip
\noindent\textbf{quiz\_pagine\_311\_320\_002}: A. Per $n$ intero, $\cos(\pi n^2)=(-1)^n$, percio $a_n=(-1)^n n$: le sottosuccessioni pari e dispari divergono con segni opposti. Anche E descrive tale comportamento, ma la chiave ufficiale indica A.\par\medskip
\noindent\textbf{quiz\_pagine\_331\_340\_003}: E. Per $n=4k$ si ha $a_n=0$, per $n=4k+1$ si ha $a_n=n$ e per $n=4k+3$ si ha $a_n=-n$: non esiste limite. La B descrive pure il comportamento oscillante; la fonte privilegia E.\par\medskip
\noindent\textbf{quiz\_pagine\_331\_340\_021}: E. Per $n$ pari domina $3^n$ positivamente; per $n$ dispari domina $-3^n$. La successione oscilla con modulo crescente.\par\medskip
\noindent\textbf{quiz\_pagine\_401\_410\_010}: D. Ponendo $y=\sqrt[7]{x+7}$, si ottiene $x=y^7-7$, quindi $f^{-1}(x)=x^7-7$.\par\medskip
\noindent\textbf{quiz\_pagine\_461\_470\_002}: C. La negazione della limitatezza $\exists M>0\ \forall n:\ |a_n|\le M$ \textbackslash{}`e $\forall M>0\ \exists n:\ |a_n|>M$.\par\medskip
\noindent\textbf{quiz\_pagine\_501\_510\_027}: E. Una successione limitata inferiormente ammette un minorante reale $m$; scegliendo $k<\min\{m,0\}$ si ha $a_n>k$ per ogni $n$.\par\medskip
\subsubsection*{2.2 -- Limiti di successioni}
\noindent\textbf{quiz\_pagine\_001\_010\_006}: A. Spiegazione non ancora disponibile.\par\medskip
\noindent\textbf{quiz\_pagine\_001\_010\_010}: E. Spiegazione non ancora disponibile.\par\medskip
\noindent\textbf{quiz\_pagine\_001\_010\_025}: B. Spiegazione non ancora disponibile.\par\medskip
\noindent\textbf{quiz\_pagine\_011\_020\_007}: A. $\cos(\pi n^2)=(-1)^n$, dunque $a_n=(-1)^n n$ oscilla senza limite.\par\medskip
\noindent\textbf{quiz\_pagine\_021\_030\_019}: A. Entrambi gli addendi sono non negativi, dunque ciascuno tende a zero.\par\medskip
\noindent\textbf{quiz\_pagine\_021\_030\_029}: B. $\sin(\pi n/2)$ assume ciclicamente $1,0,-1,0$, quindi la successione diverge oscillando e non ha limite.\par\medskip
\noindent\textbf{quiz\_pagine\_031\_040\_019}: A. Spiegazione non ancora disponibile.\par\medskip
\noindent\textbf{quiz\_pagine\_031\_040\_023}: E. Spiegazione non ancora disponibile.\par\medskip
\noindent\textbf{quiz\_pagine\_031\_040\_024}: D. Spiegazione non ancora disponibile.\par\medskip
\noindent\textbf{quiz\_pagine\_031\_040\_025}: C. Spiegazione non ancora disponibile.\par\medskip
\noindent\textbf{quiz\_pagine\_031\_040\_038}: B. Spiegazione non ancora disponibile.\par\medskip
\noindent\textbf{quiz\_pagine\_031\_040\_041}: B. Spiegazione non ancora disponibile.\par\medskip
\noindent\textbf{quiz\_pagine\_031\_040\_043}: A. Spiegazione non ancora disponibile.\par\medskip
\noindent\textbf{quiz\_pagine\_041\_050\_012}: C. Il rapporto tra radice cubica e termine lineare tende a zero; il seno è limitato.\par\medskip
\noindent\textbf{quiz\_pagine\_041\_050\_027}: non verificata / nessuna alternativa corretta. Soluzione non riportata nel PDF; da verificare.\par\medskip
\noindent\textbf{quiz\_pagine\_041\_050\_028}: non verificata / nessuna alternativa corretta. Soluzione non riportata nel PDF; da verificare.\par\medskip
\noindent\textbf{quiz\_pagine\_041\_050\_029}: non verificata / nessuna alternativa corretta. Soluzione non riportata nel PDF; da verificare.\par\medskip
\noindent\textbf{quiz\_pagine\_071\_080\_007}: A. Poiché $\sin(\pi/6)=1/2$ e l'esponente è negativo e tende a $-\infty$, la successione diverge a $+\infty$.\par\medskip
\noindent\textbf{quiz\_pagine\_081\_090\_006}: C. $b_n=c_n-a_n$ è la differenza fra una successione convergente e una che tende a $+\infty$; dunque $b_n\to-\infty$.\par\medskip
\noindent\textbf{quiz\_pagine\_081\_090\_016}: C. $\sin(5/(1+n))\sim5/n$ e $(\log n)/n\to0$, quindi il prodotto tende a $0$.\par\medskip
\noindent\textbf{quiz\_pagine\_101\_110\_004}: B. Il logaritmo vale $\log(1+\sqrt n)/n^2\to0$, quindi il limite è $e^0=1$.\par\medskip
\noindent\textbf{quiz\_pagine\_101\_110\_015}: D. Per il teorema del confronto, da $0\leq a_n\leq b_n\to0$ segue $a_n\to0$.\par\medskip
\noindent\textbf{quiz\_pagine\_111\_120\_032}: C. La base tende a $1$ e l'esponente a $-\pi/2$, valore finito: il limite e' $1$.\par\medskip
\noindent\textbf{quiz\_pagine\_121\_130\_003}: C. Poiche\textbackslash{}'e $-1\leq\sin n\leq1$, risulta $3^n-1\leq a_n\leq3^n+1$ e $a_n\to+\infty$.\par\medskip
\noindent\textbf{quiz\_pagine\_121\_130\_020}: C. La differenza $a_n-b_n=(-1)^n n/(n^2+1)$ tende a $0$; pertanto $a_n$ e $b_n$ hanno lo stesso comportamento di convergenza.\par\medskip
\noindent\textbf{quiz\_pagine\_131\_140\_001}: E. Ponendo $c_n=(1+6/n)^{2n}$, il limite e\textbackslash{}'e $e^{12}$, non presente tra le alternative.\par\medskip
\noindent\textbf{quiz\_pagine\_131\_140\_002}: A. Gli esponenziali dominano i polinomi: $c_n\sim(9/10)^n\to0$.\par\medskip
\noindent\textbf{quiz\_pagine\_131\_140\_003}: D. Il modulo e\textbackslash{}'e $\log(n+1)/(n+1)\to0$; applicare la definizione di limite con $\varepsilon=0.002$.\par\medskip
\noindent\textbf{quiz\_pagine\_151\_160\_011}: C. $n\log(1+4/(n^2+4))\sim4/n\to0$, quindi il limite e\textbackslash{}'e $e^0=1$.\par\medskip
\noindent\textbf{quiz\_pagine\_151\_160\_024}: A. $\sqrt[n]{3^n-2^n}=3\bigl(1-(2/3)^n\bigr)^{1/n}\to3$.\par\medskip
\noindent\textbf{quiz\_pagine\_171\_180\_029}: D. $|a_n|\leq(\log n+1)/(n^2+4)\to0$.\par\medskip
\noindent\textbf{quiz\_pagine\_181\_190\_018}: D. Dividendo per $n!$, il numeratore diventa $n+1-2n/n!\sim n$; il denominatore $\sin n-2n\sim-2n$, quindi il limite e' $-1/2$.\par\medskip
\noindent\textbf{quiz\_pagine\_211\_220\_038}: D. Il numeratore cresce al piu' come $10\log n+1$, trascurabile rispetto a $\sqrt n$.\par\medskip
\noindent\textbf{quiz\_pagine\_221\_230\_014}: A. Il denominatore cresce in modulo come $2^n$, il numeratore come $n$; dunque $a_n\to0$.\par\medskip
\noindent\textbf{quiz\_pagine\_221\_230\_015}: B. Ogni successione monotona ammette limite esteso, finito oppure infinito; non e' indeterminata.\par\medskip
\noindent\textbf{quiz\_pagine\_251\_260\_019}: D. Per $n$ grande $a_n<0$ e sempre $b_n<0$, quindi $a_n/b_n>0$: il limite non puo' essere negativo.\par\medskip
\noindent\textbf{quiz\_pagine\_261\_270\_029}: D. La chiave originale indica D; tuttavia $a_n/b_n$ puo' non avere limite (ad esempio $a_n=-n$, $b_n=-n$ per $n$ pari e $b_n=-n^2$ per $n$ dispari). Pertanto nessuna delle alternative e' impossibile: quesito difettoso.\par\medskip
\noindent\textbf{quiz\_pagine\_281\_290\_003}: C. Il logaritmo cresce piu' lentamente di $n^2$, quindi il rapporto tende a zero.\par\medskip
\noindent\textbf{quiz\_pagine\_281\_290\_024}: E. La risposta ufficiale E indica una condizione sufficiente: limitata per infinitesima e' infinitesima. Tuttavia NON e' necessaria, nonostante "solo se": per $b_n=\sqrt n$, $c_n=1/n$, il prodotto tende a zero con $b_n$ illimitata. Quesito formalmente ambiguo.\par\medskip
\noindent\textbf{quiz\_pagine\_291\_300\_006}: A. La base tende a $1$ e l'esponente tende a $\pi/2$, quindi il limite e' $1$.\par\medskip
\noindent\textbf{quiz\_pagine\_291\_300\_010}: E. $\sqrt[n]{4^n+2^n}=4(1+2^{-n})^{1/n}\to4$.\par\medskip
\noindent\textbf{quiz\_pagine\_301\_310\_005}: B. Poiche' $\log n/n^{1/3}\to0$, anche $|a_n|\to0$, quindi per definizione di limite vale la seconda proposizione.\par\medskip
\noindent\textbf{quiz\_pagine\_321\_330\_015}: A. Dato che $0\leq|a_n|,|b_n|\leq|a_n|+|b_n|\to0$, entrambe le successioni tendono a zero.\par\medskip
\noindent\textbf{quiz\_pagine\_341\_350\_014}: E. $a_n-b_n=n/(n^2-3n+10)\to0$, indipendentemente da $b_n$.\par\medskip
\noindent\textbf{quiz\_pagine\_341\_350\_015}: A. $a_n=(a_nb_n)/b_n$ e rapporto tra un numero tendente a $-5$ e uno tendente a $-\infty$: tende a zero da valori positivi.\par\medskip
\noindent\textbf{quiz\_pagine\_341\_350\_027}: C. Il modulo e maggiorato da $\sqrt[3]{2n+1}/(3n-1)\to0$; per confronto $a_n\to0$.\par\medskip
\noindent\textbf{quiz\_pagine\_351\_360\_014}: E. $a_n-b_n=n/(n^2-3n+10)\to0$, indipendentemente da $b_n$.\par\medskip
\noindent\textbf{quiz\_pagine\_351\_360\_015}: A. $a_n=(a_nb_n)/b_n$ e rapporto tra un numero tendente a $-5$ e uno tendente a $-\infty$: tende a zero da valori positivi.\par\medskip
\noindent\textbf{quiz\_pagine\_351\_360\_027}: C. Il modulo e maggiorato da $\sqrt[3]{2n+1}/(3n-1)\to0$; per confronto $a_n\to0$.\par\medskip
\noindent\textbf{quiz\_pagine\_371\_380\_022}: A. La base e $1/2$ e l\textbackslash{}textquotesingle esponente e negativo con modulo almeno $\sqrt n$, quindi $a_n\geq2^{\sqrt n}\to+\infty$.\par\medskip
\noindent\textbf{quiz\_pagine\_381\_390\_005}: A. Qualora il minimo esista, e al piu $f(0)=-2$. Il testo non assume continuita e quindi il minimo potrebbe non essere raggiunto: rigorosamente si puo affermare solo $\inf f\leq-2$.\par\medskip
\noindent\textbf{quiz\_pagine\_391\_400\_001}: C. Poiche $\sin(5/(n+1))\sim5/(n+1)$ e $\log n/n\to0$, il prodotto tende a $0$.\par\medskip
\noindent\textbf{quiz\_pagine\_411\_420\_004}: B. Il logaritmo vale $\log(1+\sqrt n)/n^2\to0$, quindi il limite \textbackslash{}`e $e^0=1$.\par\medskip
\noindent\textbf{quiz\_pagine\_411\_420\_015}: D. Per il teorema del confronto, $0\leq a_n\leq b_n\to0$ implica $a_n\to0$.\par\medskip
\noindent\textbf{quiz\_pagine\_421\_430\_015}: C. La base tende a $1$ e l\textbackslash{}'esponente tende a $-\pi/2$, quindi il limite \textbackslash{}`e $1$.\par\medskip
\noindent\textbf{quiz\_pagine\_431\_440\_018}: C. $3n-1\leq a_n\leq3n+1$, quindi $a_n\to+\infty$.\par\medskip
\noindent\textbf{quiz\_pagine\_431\_440\_035}: C. $a_n-b_n=(-1)^nn/(n^2+1)\to0$, quindi le due successioni hanno lo stesso limite quando esiste.\par\medskip
\noindent\textbf{quiz\_pagine\_441\_450\_012}: B. $|a_n|\le3n$ ma il fattore oscillante non impone limitatezza. La non limitatezza richiede un argomento sulla distribuzione dei valori $n^2+1$ modulo $2\pi$; non \textbackslash{}`e una conseguenza del solo $|\sin|\le1$.\par\medskip
\noindent\textbf{quiz\_pagine\_441\_450\_021}: E. $c_n=(1+6/n)^{2n}\to e^{12}$, che non compare fra le prime quattro risposte.\par\medskip
\noindent\textbf{quiz\_pagine\_441\_450\_022}: A. Dominano gli esponenziali: $c_n\sim(9/10)^n\to0$.\par\medskip
\noindent\textbf{quiz\_pagine\_441\_450\_023}: D. $|c_n|=\log(n+1)/(n+1)\to0$, dunque la condizione epsilon vale.\par\medskip
\noindent\textbf{quiz\_pagine\_451\_460\_019}: E. Una successione limitata inferiormente ammette un minorante reale $m$; scegliendo $k<\min\{m,0\}$ si ha $a_n>k$ per ogni $n$.\par\medskip
\noindent\textbf{quiz\_pagine\_461\_470\_001}: A. Dato che $a_n\to l\neq0$, la successione $1/a_n$ \textbackslash{}`e definitivamente limitata; i primi termini sono in numero finito e non nulli. Il prodotto $b_n/a_n$ \textbackslash{}`e dunque limitato.\par\medskip
\noindent\textbf{quiz\_pagine\_461\_470\_021}: C. $n\log(1+4/(n^2+4))\sim4n/(n^2+4)\to0$, perci\textbackslash{}`o il limite vale $e^0=1$.\par\medskip
\noindent\textbf{quiz\_pagine\_471\_480\_010}: A. Si raccoglie $3^n$: $\sqrt[n]{3^n-2^n}=3\left(1-(2/3)^n\right)^{1/n}\to3$.\par\medskip
\noindent\textbf{quiz\_pagine\_491\_500\_013}: D. $|a_n|\le(\log n+1)/(n^2+4)\to0$; si applica il teorema del confronto.\par\medskip
\noindent\textbf{quiz\_pagine\_511\_520\_003}: C. Dividendo per $x$, numeratore e denominatore tendono rispettivamente a $5$ e $-1$; il rapporto tende a $-5$.\par\medskip
\noindent\textbf{quiz\_pagine\_511\_520\_004}: B. La definizione di limite finito all'infinito implica che, scegliendo $M$ sufficientemente grande, i valori $f(x)$ per $x<-M$ appartengono a un intervallo limitato intorno a $-5$.\par\medskip
\noindent\textbf{quiz\_pagine\_511\_520\_005}: E. $\cosh x=1+x^2/2+o(x^2)$, perci\textbackslash{}`o $(1-\cosh x)/(2x^2)\to-1/4$. Il valore assegnato in $0$ non conta.\par\medskip
\noindent\textbf{quiz\_pagine\_511\_520\_006}: D. $a_n=(1+2/(n-1))^n\to e^2$.\par\medskip
\noindent\textbf{quiz\_pagine\_511\_520\_007}: B. Ogni successione monotona ha limite nell'insieme dei reali estesi: pu\textbackslash{}`o convergere oppure divergere a $\pm\infty$, ma non oscillare senza limite.\par\medskip
\noindent\textbf{quiz\_pagine\_511\_520\_024}: E. Poich'e il limite \textbackslash{}`e $1>0$, esiste un intorno bucato di $2$ in cui $f(x)>0$. Nessuna condizione \textbackslash{}`e imposta al valore $f(2)$.\par\medskip
\noindent\textbf{quiz\_pagine\_511\_520\_025}: C. $9^x-1\sim x\log9$ e $(1+8x)^{1/4}-1\sim2x$: il rapporto tende a $\log9/2=\log3$.\par\medskip
\noindent\textbf{quiz\_pagine\_511\_520\_026}: C. Lungo $x=k\pi$ il prodotto \textbackslash{}`e $0$, mentre lungo $x=\pi/2+k\pi$ tende a $+\infty$. Il limite non esiste.\par\medskip
\noindent\textbf{quiz\_pagine\_511\_520\_027}: B. $f(x)=x+1+o(1)$ e $g(x)=2x+o(1)$; pertanto $f-g=-x+1+o(1)\to-\infty$.\par\medskip
\noindent\textbf{quiz\_pagine\_511\_520\_028}: A. Poich'e $|(-2)^n|=2^n$ domina $n$ e il numeratore \textbackslash{}`e $O(n)$, si ha $a_n\to0$ da entrambi i sottogruppi.\par\medskip
\subsection*{Capitolo 3}
\subsubsection*{3.1 -- Definizioni e prime proprietà. Limiti}
\noindent\textbf{quiz\_pagine\_001\_010\_022}: D. Spiegazione non ancora disponibile.\par\medskip
\noindent\textbf{quiz\_pagine\_001\_010\_026}: C. Spiegazione non ancora disponibile.\par\medskip
\noindent\textbf{quiz\_pagine\_011\_020\_004}: D. $x^3\log(1+1/(4x^3))\to1/4$ per $\log(1+t)\sim t$.\par\medskip
\noindent\textbf{quiz\_pagine\_011\_020\_005}: E. $x[(1+2/x)^{1/4}-1]\to 1/2$.\par\medskip
\noindent\textbf{quiz\_pagine\_031\_040\_003}: C. Spiegazione non ancora disponibile.\par\medskip
\noindent\textbf{quiz\_pagine\_031\_040\_026}: A. Spiegazione non ancora disponibile.\par\medskip
\noindent\textbf{quiz\_pagine\_041\_050\_020}: A. Risposta ufficiale indicata nella tabella della decima simulazione.\par\medskip
\noindent\textbf{quiz\_pagine\_041\_050\_033}: non verificata / nessuna alternativa corretta. Soluzione non riportata nel PDF; da verificare.\par\medskip
\noindent\textbf{quiz\_pagine\_051\_060\_014}: D. La definizione esprime che $f(x)$ tende a $5$ quando $x$ tende a $+\infty$.\par\medskip
\noindent\textbf{quiz\_pagine\_051\_060\_015}: A. La risposta ufficiale è A; la condizione con un solo $M$ fissato non equivale alla definizione completa di limite.\par\medskip
\noindent\textbf{quiz\_pagine\_061\_070\_003}: A. Il numeratore è equivalente a x\textsuperscript{2}/2 e il denominatore a 3x\textsuperscript{2}; limite 1/6.\par\medskip
\noindent\textbf{quiz\_pagine\_071\_080\_013}: A. Per i limiti notevoli $e^{2x}-1\sim2x$ e $\log(1+3x)\sim3x$, quindi il limite vale $2/3$.\par\medskip
\noindent\textbf{quiz\_pagine\_081\_090\_001}: E. Per $x\to0$, il rapporto è asintotico a $\frac{3x}{25x^2/2}=\frac{6}{25x}$: i limiti laterali hanno segni opposti.\par\medskip
\noindent\textbf{quiz\_pagine\_081\_090\_005}: B. Il denominatore è compreso fra $4$ e $5$ e il numeratore tende a $+\infty$; per confronto il rapporto tende a $+\infty$.\par\medskip
\noindent\textbf{quiz\_pagine\_081\_090\_026}: A. Per $x\to0$, $\log(1+7x)\sim7x$ e $e^{4x}-1\sim4x$, quindi il limite è $7/4$.\par\medskip
\noindent\textbf{quiz\_pagine\_081\_090\_027}: A. $(4^x+2^x)^{1/x}=4(1+2^{-x})^{1/x}\to4$, poiché $\log(1+2^{-x})/x\to0$.\par\medskip
\noindent\textbf{quiz\_pagine\_091\_100\_022}: A. La definizione di limite con $\varepsilon=1$ dà la condizione A, con intorno bucato. Il valore $f(1)$ non è determinato.\par\medskip
\noindent\textbf{quiz\_pagine\_101\_110\_005}: A. Poiché $x^2\to+\infty$ quando $x\to+\infty$, si ha $f(x^2)\to0$. Non è noto il segno.\par\medskip
\noindent\textbf{quiz\_pagine\_101\_110\_009}: A. Usando $2^x-1\sim x\log2$ e $\sin5x\sim5x$, il limite vale $\log2/5$.\par\medskip
\noindent\textbf{quiz\_pagine\_101\_110\_019}: A. Si scrive $x^2-9=(x-3)(x+3)$ e si usa $\sin u/u\to1$: il limite è $x+3\to6$.\par\medskip
\noindent\textbf{quiz\_pagine\_111\_120\_006}: E. Il numeratore e' equivalente a $4x$ e il denominatore a $9x^2/2$: il rapporto e' equivalente a $8/(9x)$, con limiti laterali opposti.\par\medskip
\noindent\textbf{quiz\_pagine\_111\_120\_014}: C. Il denominatore resta compreso fra $6$ e $7$, mentre il numeratore tende a $+\infty$.\par\medskip
\noindent\textbf{quiz\_pagine\_111\_120\_018}: B. Il numeratore e' compreso fra $-5$ e $-3$; dividendo per $x<0$ il rapporto tende a $+\infty$.\par\medskip
\noindent\textbf{quiz\_pagine\_111\_120\_029}: E. $\frac{x^2+3}{x-4}-x=\frac{4x+3}{x-4}\to4$, quindi il rapporto tende a $e^4$.\par\medskip
\noindent\textbf{quiz\_pagine\_111\_120\_030}: B. $1-\cos4x\sim8x^2$ e $\sinh^2 3x\sim9x^2$, dunque il limite e' $8/9$.\par\medskip
\noindent\textbf{quiz\_pagine\_121\_130\_004}: D. Per $x\to1^+$ o $x\to1^-$ nel dominio, $\log x^2=2\log x$ e $\log(1+\log x)\sim\log x$, quindi il rapporto tende a $2$.\par\medskip
\noindent\textbf{quiz\_pagine\_121\_130\_012}: E. Per $x\to0$, $\sinh(5x)\sim5x$ e $1-\cos(3x)\sim9x^2/2$; il limite vale $25/(9/2)=50/9$.\par\medskip
\noindent\textbf{quiz\_pagine\_121\_130\_016}: A. Una funzione positiva non puo\textbackslash{}'o avere limite negativo; il limite puo\textbackslash{}'o essere $0$, un numero positivo oppure $+\infty$.\par\medskip
\noindent\textbf{quiz\_pagine\_121\_130\_019}: D. Si scrive $\log(e^{2x}+x)=2x+\log(1+xe^{-2x})=2x+o(1)$; il limite e\textbackslash{}'e $6$.\par\medskip
\noindent\textbf{quiz\_pagine\_131\_140\_024}: E. Il logaritmo dell\textbackslash{}'espressione e\textbackslash{}'e $\log(1+4x^2)/\log x\to2$, percio\textbackslash{}'e il limite e\textbackslash{}'e $e^2$.\par\medskip
\noindent\textbf{quiz\_pagine\_141\_150\_001}: B. Poiche\textbackslash{}'e $\sin(1/x)\sim1/x$, il prodotto e\textbackslash{}'e asintotico a $3x^2-4x$, che tende a $+\infty$.\par\medskip
\noindent\textbf{quiz\_pagine\_151\_160\_001}: E. Per $x>0$, $\sqrt{x^4+x^6}=x^2\sqrt{1+x^2}=x^2+\frac12x^4+o(x^4)$.\par\medskip
\noindent\textbf{quiz\_pagine\_151\_160\_021}: E. $\cos(1/x)\to1$ mentre $e^{x/2}\to+\infty$, percio\textbackslash{}'o il prodotto tende a $-\infty$.\par\medskip
\noindent\textbf{quiz\_pagine\_151\_160\_023}: E. Per definizione $f(x)/(4x)\to1$. Il resto $f(x)-4x$ puo\textbackslash{}'o invece oscillare o divergere.\par\medskip
\noindent\textbf{quiz\_pagine\_161\_170\_002}: B. Per $t=1/x\to0^+$, $e^{3t}-\cos t=3t+O(t^2)$; il prodotto e\textbackslash{}'e asintotico a $3x$, quindi diverge a $+\infty$.\par\medskip
\noindent\textbf{quiz\_pagine\_161\_170\_006}: A. L'arcotangente e\textbackslash{}'e limitata e tende a $\pi/2$; pertanto $f(x)/x\to1$.\par\medskip
\noindent\textbf{quiz\_pagine\_161\_170\_007}: A. Per $x\to0^+$, $\arctan(x+x^{3/2})\sim x+x^{3/2}\sim x$ e $f(x)\sim2x$. Nota: $x^{3/2}$ reale richiede $x\geq0$, dunque il limite e\textbackslash{}'e destro.\par\medskip
\noindent\textbf{quiz\_pagine\_161\_170\_011}: B. Poiche\textbackslash{}'e $x=o(x^{2/3})$ per $x\to0$, $\arctan(x+x^{2/3})\sim x^{2/3}$.\par\medskip
\noindent\textbf{quiz\_pagine\_161\_170\_012}: E. Per $x\to+\infty$, $\cos(1/x)-e^{4/x}=-4/x+O(x^{-2})$; moltiplicando per la potenza indicata si ottiene il risultato.\par\medskip
\noindent\textbf{quiz\_pagine\_161\_170\_013}: E. Per $x\to+\infty$, $\cos(1/x)-e^{4/x}=-4/x+O(x^{-2})$; moltiplicando per la potenza indicata si ottiene il risultato.\par\medskip
\noindent\textbf{quiz\_pagine\_161\_170\_014}: A. Per $x\to+\infty$, $\cos(1/x)-e^{4/x}=-4/x+O(x^{-2})$; moltiplicando per la potenza indicata si ottiene il risultato.\par\medskip
\noindent\textbf{quiz\_pagine\_161\_170\_018}: E. Il logaritmo del termine e\textbackslash{}'e $5\log(n+3)/\log n\to5$, quindi il limite vale $e^5$.\par\medskip
\noindent\textbf{quiz\_pagine\_161\_170\_020}: B. Con $t=5x-4$, il denominatore e\textbackslash{}'e $t^2/25$ e $(1-\cos t)/t^2\to1/2$.\par\medskip
\noindent\textbf{quiz\_pagine\_171\_180\_008}: A. Ponendo $t=x-4/5$, $1-\cos(5t)\sim25t^2/2$.\par\medskip
\noindent\textbf{quiz\_pagine\_171\_180\_011}: B. Il numeratore e' $-\frac12\log x\to-\infty$, il denominatore resta tra $1$ e $2$.\par\medskip
\noindent\textbf{quiz\_pagine\_171\_180\_023}: B. $5^x-1\sim x\log5$ e $\sin4x\sim4x$.\par\medskip
\noindent\textbf{quiz\_pagine\_181\_190\_004}: D. Per continuita', $e^{f(x)}\to e^9>4$; in un intorno di $3$ si ha dunque $e^{f(x)}\geq4$.\par\medskip
\noindent\textbf{quiz\_pagine\_191\_200\_009}: A. $\sin(x/2+x^2/3)\sim x/2$ e $\tan^2x\sim x^2$; il prodotto tende a $1/2$.\par\medskip
\noindent\textbf{quiz\_pagine\_191\_200\_013}: B. Il numeratore appartiene a $[-5,-4]$, il denominatore tende a $-\infty$, quindi il rapporto tende a zero da valori positivi.\par\medskip
\noindent\textbf{quiz\_pagine\_211\_220\_035}: E. Razionalizzando si ottiene $1/(\sqrt{1+x^2}+1)\to1/2$.\par\medskip
\noindent\textbf{quiz\_pagine\_221\_230\_006}: B. Per $x\to-\infty$ il termine $e^x$ tende a zero e la funzione oscilla vicino a $\cos x$, attraversando lo zero infinite volte.\par\medskip
\noindent\textbf{quiz\_pagine\_221\_230\_020}: B. Per $x\neq1$ la radice assume tutti i valori di $[0,2]$ eccetto $1$, che deve essere assunto in $x=1$.\par\medskip
\noindent\textbf{quiz\_pagine\_221\_230\_021}: C. La continuita' in zero raccorda i due intervalli di monotonia e permette il confronto tra punti di segno opposto.\par\medskip
\noindent\textbf{quiz\_pagine\_221\_230\_025}: B. $-x^2+2x=1-(x-1)^2\leq1$; la funzione positiva ha massimo $3$ e tende a zero senza raggiungerlo.\par\medskip
\noindent\textbf{quiz\_pagine\_221\_230\_028}: A. Per una potenza con esponente reale variabile si richiede $\ln(5-x^2)>0$, ossia $5-x^2>1$.\par\medskip
\noindent\textbf{quiz\_pagine\_221\_230\_029}: C. $g(f(x))=g(x^2)=4x^2+1$.\par\medskip
\noindent\textbf{quiz\_pagine\_231\_240\_011}: C. Per $x\to-\infty$, $f(x)\to0$ (asintoto $y=0$); per $x\to-3$ il denominatore tende a zero con numeratore positivo (asintoto $x=-3$).\par\medskip
\noindent\textbf{quiz\_pagine\_231\_240\_020}: D. Fattorizzando: $(x-1)(x+4)/[(x-1)(x+5)]\to5/6$.\par\medskip
\noindent\textbf{quiz\_pagine\_241\_250\_013}: C. Per il teorema del confronto $|f(x)|\leq3^{-x}\to0$, quindi il limite e' zero.\par\medskip
\noindent\textbf{quiz\_pagine\_241\_250\_016}: A. $\log x^{-5/2}=-(5/2)\log x$, dunque il numeratore e' asintotico a $x$ e il denominatore resta tra 2 e 3.\par\medskip
\noindent\textbf{quiz\_pagine\_241\_250\_018}: E. $M$ indica la mantissa (parte frazionaria), limitata. Nell'ultima opzione il rapporto cresce come $x/(4\ln x)\to+\infty$, quindi NON e' un limite inesistente.\par\medskip
\noindent\textbf{quiz\_pagine\_241\_250\_019}: A. Il seno e' limitato e il logaritmo cresce come $2\ln x$, percio' il rapporto con $x-6$ tende a zero, non a $+\infty$.\par\medskip
\noindent\textbf{quiz\_pagine\_241\_250\_020}: B. Per $x\to0^+$, $\sin(\pi x/2)\to0$ e $\log(\sqrt[3]x+x+x^2)\to-\infty$, ma il prodotto tende a zero. Nota: la funzione reale non e' definita per $x<0$ sufficientemente vicino a zero; il limite e' destro.\par\medskip
\noindent\textbf{quiz\_pagine\_251\_260\_013}: D. $1\leq\sin x+2\leq3$, percio' $x(\sin x+2)\geq x\to+\infty$: la D e' falsa. Nell'originale il simbolo di limite inesistente e' riportato come simbolo speciale.\par\medskip
\noindent\textbf{quiz\_pagine\_261\_270\_005}: E. Le sottosuccessioni $f(3n)=0$ e $f(3n+1)=3^{3n+1}\to+\infty$ hanno comportamenti incompatibili.\par\medskip
\noindent\textbf{quiz\_pagine\_261\_270\_010}: C. $\log(x^{-5/2})$ e' definito solo per $x>0$; non ha senso il limite $x\to-\infty$.\par\medskip
\noindent\textbf{quiz\_pagine\_261\_270\_012}: E. Per $x\to0^+$ il denominatore $e^{1/x}$ domina $1/x$: il rapporto tende a $0$, non a $+\infty$.\par\medskip
\noindent\textbf{quiz\_pagine\_261\_270\_017}: B. $\sin(x^2-25)\sim(x-5)(x+5)$, quindi il limite vale $10$.\par\medskip
\noindent\textbf{quiz\_pagine\_261\_270\_018}: A. Il denominatore e' compreso tra $3$ e $4$, mentre il numeratore tende a $+\infty$.\par\medskip
\noindent\textbf{quiz\_pagine\_271\_280\_002}: B. Per $x\to-\infty$, $e^x\to0$; in ogni intervallo tra multipli successivi di $\pi$ la parte coseno cambia segno e vi sono infiniti zeri per continuita'.\par\medskip
\noindent\textbf{quiz\_pagine\_271\_280\_004}: E. Dalla definizione di limite con $\varepsilon=1$ esiste un intorno bucato di $2$ in cui $0<f(x)<2$. Il valore $f(2)$ non e' vincolato.\par\medskip
\noindent\textbf{quiz\_pagine\_271\_280\_012}: B. Il valore $1$ e' assunto da $\sqrt x$ soltanto per $x=1$; quindi deve essere $\alpha=1$.\par\medskip
\noindent\textbf{quiz\_pagine\_271\_280\_013}: C. La continuita' in zero raccorda i due intervalli di monotonia; per $x<0<y$, $f(x)\geq f(0)\geq f(y)$.\par\medskip
\noindent\textbf{quiz\_pagine\_291\_300\_003}: B. Il numeratore tende a $-2$ e il denominatore a $+\infty$: il rapporto tende a zero.\par\medskip
\noindent\textbf{quiz\_pagine\_291\_300\_009}: E. Ponendo $h=5(x-4/5)$ e usando $1-\cos h\sim h^2/2$ si trova $25/2$.\par\medskip
\noindent\textbf{quiz\_pagine\_291\_300\_015}: A. $\sqrt{1+4/x}=1+2/x+O(x^{-2})$ e $\cos(1/x)=1+O(x^{-2})$: il prodotto e' $2x+O(1)\to+\infty$.\par\medskip
\noindent\textbf{quiz\_pagine\_291\_300\_020}: E. Le sottosuccessioni nei pari e nei dispari tendono rispettivamente a $-\infty$ e $+\infty$, quindi il limite non esiste.\par\medskip
\noindent\textbf{quiz\_pagine\_291\_300\_022}: B. $\alpha x-4\sin x=(\alpha-4)x+\frac23x^3+o(x^3)$; per $\alpha<4$ il termine $(\alpha-4)/x^2$ porta a $-\infty$.\par\medskip
\noindent\textbf{quiz\_pagine\_291\_300\_026}: B. Il logaritmo e' $\log(1+\sin x)/x^2\sim1/x\to+\infty$.\par\medskip
\noindent\textbf{quiz\_pagine\_291\_300\_032}: B. $\sin(x^2-25)\sim(x-5)(x+5)$, quindi il limite e' $10$.\par\medskip
\noindent\textbf{quiz\_pagine\_291\_300\_036}: B. $5^x-1\sim x\log5$ e $\sin4x\sim4x$.\par\medskip
\noindent\textbf{quiz\_pagine\_301\_310\_002}: D. Dividendo numeratore e denominatore per $x^2$, tutti gli altri addendi tendono a zero; resta $-1$.\par\medskip
\noindent\textbf{quiz\_pagine\_301\_310\_006}: C. Usando $\sin(3x)\sim3x$ e $1-\cos(2x)\sim2x^2$, il rapporto tende a $9/2$.\par\medskip
\noindent\textbf{quiz\_pagine\_301\_310\_019}: C. $f(x)=x+o(x)$ e $g(x)=x^2+o(x^2)=o(x)$; pertanto la somma e' $x+o(x)$.\par\medskip
\noindent\textbf{quiz\_pagine\_301\_310\_024}: D. Il limite e' $\lim x^3\log(1+1/(4x^3))=1/4$ poiche' $\log(1+u)\sim u$.\par\medskip
\noindent\textbf{quiz\_pagine\_301\_310\_025}: E. Si scrive $x((1+2/x)^{1/4}-1)\sim x\cdot(2/(4x))=1/2$.\par\medskip
\noindent\textbf{quiz\_pagine\_311\_320\_031}: E. Poiche $x\sin x\sim x^2$ e $\log(1+x^2)\sim x^2$, il rapporto e asintotico a $x^{\alpha-2}+1$, mai zero per $x\to0^+$.\par\medskip
\noindent\textbf{quiz\_pagine\_321\_330\_012}: E. Il numeratore \textbackslash{}`e $\sqrt{x}+O(x)$, il denominatore $-2\sqrt{x}+O(x^2)$: il rapporto tende a $-1/2$.\par\medskip
\noindent\textbf{quiz\_pagine\_321\_330\_030}: D. Per continuit\textbackslash{}`a $9x+2^x\to1$ e $x-2^x\to-1$.\par\medskip
\noindent\textbf{quiz\_pagine\_321\_330\_031}: E. Dividendo per $x$ si ottiene $(1+1/\sqrt x)/(10+1/\sqrt x)\to1/10$.\par\medskip
\noindent\textbf{quiz\_pagine\_331\_340\_013}: C. Per permanenza del segno in forma non stretta, se $f(x)>7$ e il limite esiste, allora $l\geq7$.\par\medskip
\noindent\textbf{quiz\_pagine\_361\_370\_014}: D. E la definizione epsilon-soglia del limite finito $5$ per $x\to+\infty$.\par\medskip
\noindent\textbf{quiz\_pagine\_361\_370\_015}: A. RISPOSTA UFFICIALE A, MA UN SOLO VALORE FISSATO $M=10^{-5}$ NON DEFINISCE UN LIMITE. Servirebbe $\forall M>0\ \exists\epsilon$ con la stessa implicazione. La formulazione originale e insufficiente.\par\medskip
\noindent\textbf{quiz\_pagine\_381\_390\_011}: C. $b_n=c_n-a_n$ e una successione convergente meno una che tende a $+\infty$: $b_n\to-\infty$.\par\medskip
\noindent\textbf{quiz\_pagine\_391\_400\_011}: A. Per $x\to0$, $\log(1+7x)\sim7x$ e $e^{4x}-1\sim4x$, quindi $7/4$.\par\medskip
\noindent\textbf{quiz\_pagine\_391\_400\_020}: A. $\log(1+3x)=3x+O(x^2)$, percio il numeratore e $4x^3+6x-6x+O(x^2)=O(x^2)$ e il rapporto tende a $0$.\par\medskip
\noindent\textbf{quiz\_pagine\_401\_410\_005}: C. $\sinh(4x)\sim4x$ e $1-\cos(3x)\sim9x^2/2$; il rapporto tende a $32/9$.\par\medskip
\noindent\textbf{quiz\_pagine\_401\_410\_014}: A. Dalla definizione di limite con $\varepsilon=1$ segue $|f(x)-3|<1$ per $0<|x-1|<\delta$. Non si conosce $f(1)$.\par\medskip
\noindent\textbf{quiz\_pagine\_401\_410\_020}: B. Posto $t=\sqrt3\,x$, $e^t-1\sim t$ e $\sin t\sim t$, dunque il limite vale $1$.\par\medskip
\noindent\textbf{quiz\_pagine\_411\_420\_009}: A. Poich\textbackslash{}'e $2^x-1\sim x\log2$ e $\sin5x\sim5x$, il limite vale $\log2/5$.\par\medskip
\noindent\textbf{quiz\_pagine\_411\_420\_019}: A. $\sin(x^2-9)\sim(x-3)(x+3)$, dunque il limite \textbackslash{}`e $6$.\par\medskip
\noindent\textbf{quiz\_pagine\_411\_420\_026}: E. Il rapporto \textbackslash{}`e asintotico a $4x/((9/2)x^2)=8/(9x)$: i limiti laterali sono opposti e quello bilatero non esiste. La fonte usa il simbolo $\nexists$.\par\medskip
\noindent\textbf{quiz\_pagine\_411\_420\_034}: C. Il denominatore resta fra $6$ e $7$, mentre il numeratore tende a $+\infty$.\par\medskip
\noindent\textbf{quiz\_pagine\_421\_430\_001}: B. Il numeratore \textbackslash{}`e compreso fra $-5$ e $-3$, mentre $x\to0^-$: il rapporto tende a $+\infty$.\par\medskip
\noindent\textbf{quiz\_pagine\_421\_430\_012}: E. $\frac{x^2+3}{x-4}-x=\frac{4x+3}{x-4}\to4$; il limite vale $e^4$.\par\medskip
\noindent\textbf{quiz\_pagine\_421\_430\_013}: B. $1-\cos4x\sim8x^2$ e $\sinh^2(3x)\sim9x^2$, quindi $8/9$.\par\medskip
\noindent\textbf{quiz\_pagine\_421\_430\_027}: B. $1-\cos(7x^2)\sim(49/2)x^4$, dunque il logaritmo \textbackslash{}`e $4\log x+O(1)$ e il limite vale $4$.\par\medskip
\noindent\textbf{quiz\_pagine\_431\_440\_019}: D. $\log x^2=2\log x$ vicino a $1$ e $\log(1+\log x)\sim\log x$; rapporto $2$.\par\medskip
\noindent\textbf{quiz\_pagine\_431\_440\_027}: E. $\sinh^2(5x)\sim25x^2$, $1-\cos(3x)\sim9x^2/2$, rapporto $50/9$.\par\medskip
\noindent\textbf{quiz\_pagine\_431\_440\_031}: A. Un limite di funzione positiva pu\textbackslash{}`o essere zero, positivo finito oppure $+\infty$.\par\medskip
\noindent\textbf{quiz\_pagine\_431\_440\_034}: D. $\log(e^{2x}+x)=2x+\log(1+xe^{-2x})=2x+o(1)$; il prodotto tende a $6$.\par\medskip
\noindent\textbf{quiz\_pagine\_441\_450\_019}: B. Confronto: $f\le g$ e $f\to+\infty$ impongono $g\to+\infty$.\par\medskip
\noindent\textbf{quiz\_pagine\_461\_470\_024}: A. $\lim_{x\to0^-}\sin(\alpha x)/x=\alpha$, mentre $f(0)=1/2$. La continuit\textbackslash{}`a richiede $\alpha=1/2$.\par\medskip
\noindent\textbf{quiz\_pagine\_471\_480\_007}: E. L'esponenziale $e^{x/2}\to+\infty$, mentre il coseno \textbackslash{}`e limitato; il prodotto con $3x^2>0$ tende a $-\infty$. Nel testo originale il fattore appare come $3x^2$.\par\medskip
\noindent\textbf{quiz\_pagine\_471\_480\_014}: B. $e^{3/x}=1+3/x+O(1/x^2)$ e $\cos(1/x)=1+O(1/x^2)$, quindi il prodotto \textbackslash{}`e $3x+O(1)\to+\infty$.\par\medskip
\noindent\textbf{quiz\_pagine\_471\_480\_024}: E. Poich'e $\cos(1/x)-e^{4/x}=-4/x+O(1/x^2)$, moltiplicando per $x^3$ si ottiene $-4x^{2}+O(x^{1})\to-\infty$.\par\medskip
\noindent\textbf{quiz\_pagine\_471\_480\_025}: E. Poich'e $\cos(1/x)-e^{4/x}=-4/x+O(1/x^2)$, moltiplicando per $x^2$ si ottiene $-4x^{1}+O(x^{0})\to-\infty$.\par\medskip
\noindent\textbf{quiz\_pagine\_481\_490\_015}: D. Le derivate di $\sinh x$ alternano $\cosh x$ e $\sinh x$; quelle dispari in zero valgono $1$, quelle pari valgono $0$.\par\medskip
\noindent\textbf{quiz\_pagine\_491\_500\_008}: D. La definizione $\varepsilon$--$\delta$ vale anche per punti isolati del dominio; la formulazione con limiti laterali no.\par\medskip
\noindent\textbf{quiz\_pagine\_491\_500\_010}: E. Entrambi i logaritmi di $f$ esistono soltanto per $x>1$, dove la propriet\textbackslash{}`a $\log a-\log b=\log(a/b)$ vale.\par\medskip
\noindent\textbf{quiz\_pagine\_501\_510\_005}: D. $f=o(g)$ significa esattamente $f/g\to0$ dove il rapporto \textbackslash{}`e definito. Il prodotto \textbackslash{}`e soltanto $o(x^{10})$, non necessariamente $o(x^{11})$.\par\medskip
\noindent\textbf{quiz\_pagine\_501\_510\_007}: D. Il numeratore \textbackslash{}`e $10\log n+O(1)$, il denominatore \textbackslash{}`e equivalente a $\sqrt n$; il rapporto tende a zero.\par\medskip
\noindent\textbf{quiz\_pagine\_501\_510\_024}: E. Gli esponenziali decadono e i termini trigonometrici restano limitati: il rapporto \textbackslash{}`e equivalente a $-2x/(3x)=-2/3$.\par\medskip
\noindent\textbf{quiz\_pagine\_501\_510\_025}: C. $f(x)\sim x^2$ e $g(x)\sim x$, dunque $f/g\sim x\to0$. Le opzioni (a) e (d) sono duplicate nel testo originale.\par\medskip
\noindent\textbf{quiz\_pagine\_501\_510\_026}: E. L\textbackslash{}'ampiezza cresce senza limite e $\sin x$ oscilla: lungo sottosuccessioni il prodotto diverge a $+\infty$ e a $-\infty$. Pertanto il limite non esiste; la (d) descrive impropriamente due limiti diversi.\par\medskip
\subsubsection*{3.2 -- Continuità di una funzione}
\noindent\textbf{quiz\_pagine\_001\_010\_027}: B. Spiegazione non ancora disponibile.\par\medskip
\noindent\textbf{quiz\_pagine\_011\_020\_022}: B. Per Lagrange esiste $c$ con $f\prime(c)=(0-6)/(6-4)=-3$.\par\medskip
\noindent\textbf{quiz\_pagine\_031\_040\_010}: C. Spiegazione non ancora disponibile.\par\medskip
\noindent\textbf{quiz\_pagine\_041\_050\_013}: E. Risposta ufficiale indicata nella tabella della decima simulazione.\par\medskip
\noindent\textbf{quiz\_pagine\_041\_050\_036}: non verificata / nessuna alternativa corretta. Soluzione non riportata nel PDF; da verificare.\par\medskip
\noindent\textbf{quiz\_pagine\_041\_050\_039}: non verificata / nessuna alternativa corretta. Soluzione non riportata nel PDF; da verificare.\par\medskip
\noindent\textbf{quiz\_pagine\_071\_080\_011}: E. La tabella ufficiale indica E. Nota: anche A è una formulazione vera del teorema degli zeri, ammettendo lo zero agli estremi; il quesito non ha risposta univoca.\par\medskip
\noindent\textbf{quiz\_pagine\_101\_110\_016}: C. La funzione è continua su $[a,10]$ soltanto se $1<a\leq10$. Fra le opzioni la C è l\textbackslash{}textquoteright unica che impone anche un intervallo non degenere. La condizione $a=10$ sarebbe pure ammissibile, ma non è proposta isolatamente.\par\medskip
\noindent\textbf{quiz\_pagine\_121\_130\_001}: D. La continuita\textbackslash{}'a su $[a,e]$ richiede $a>0$. Inoltre $f(e)=1$ e $f(a)<0$ se $a<1$: dunque $0<a<1$.\par\medskip
\noindent\textbf{quiz\_pagine\_121\_130\_005}: D. Su $[0,1]$ l\textbackslash{}'immagine e\textbackslash{}'e $[0,1]$, mentre su $[1,8]$ e\textbackslash{}'e $[0,e^8-e]$; poiche\textbackslash{}'e $e^8-e>1$, le immagini coincidono.\par\medskip
\noindent\textbf{quiz\_pagine\_121\_130\_006}: A. Per il teorema dei valori intermedi, ogni valore strettamente compreso tra $1$ e $10$ e\textbackslash{}'e assunto in $(0,5)$. Tra le opzioni solo $2$ appartiene a tale intervallo.\par\medskip
\noindent\textbf{quiz\_pagine\_141\_150\_008}: E. La funzione puo\textbackslash{}'e essere costante. Inoltre l\textbackslash{}'immagine di $[0,3]$ e\textbackslash{}'e un intervallo chiuso, non necessariamente aperto.\par\medskip
\noindent\textbf{quiz\_pagine\_141\_150\_020}: B. Per $x\neq0$ l\textbackslash{}'immagine e\textbackslash{}'e $[-1,0)\cup(0,1]$. Per non lasciare il valore $0$ mancante, deve essere $a=0$.\par\medskip
\noindent\textbf{quiz\_pagine\_151\_160\_014}: A. Il limite sinistro vale $\alpha$, mentre $f(0)=1/2$.\par\medskip
\noindent\textbf{quiz\_pagine\_171\_180\_009}: E. Per continuita' su $[a,7]$ occorre $a>5$; $f(7)=\log2>0$ e $f(a)<0$ se $5<a<6$. Per $a=6$ si ha gia' uno zero all'estremo, ma non il cambio di segno stretto richiesto dalla soluzione ufficiale.\par\medskip
\noindent\textbf{quiz\_pagine\_171\_180\_024}: D. La (d) e' la definizione epsilon-delta di continuita' relativa al dominio. La (c) riporta nel PDF $x\to0^-$ al secondo limite, non $x\to x_0^-$.\par\medskip
\noindent\textbf{quiz\_pagine\_181\_190\_017}: C. $2^x-3^{-x}=x(\log2+\log3)+o(x)$; $\log(2^x-x^2\sin(1/x))=x\log2+o(x)$, quindi il rapporto tende a $\log6/\log2$.\par\medskip
\noindent\textbf{quiz\_pagine\_201\_210\_013}: B. $h(x)=f(x)-x$ e' continua, $h(0)\ge0$ e $h(10)\le0$: per il teorema degli zeri esiste un punto fisso.\par\medskip
\noindent\textbf{quiz\_pagine\_211\_220\_044}: B. Teorema degli zeri: continuita' sull'intervallo chiuso e valori di segno opposto agli estremi.\par\medskip
\noindent\textbf{quiz\_pagine\_221\_230\_009}: C. Il numeratore e' equivalente a $x^{1/3}$ e il denominatore a $x^{1/7}$; il rapporto tende a zero (con radici reali).\par\medskip
\noindent\textbf{quiz\_pagine\_221\_230\_010}: E. Per permanenza del segno, dato il limite positivo, $f(x)>0$ in un intorno bucato di $2$.\par\medskip
\noindent\textbf{quiz\_pagine\_221\_230\_011}: C. $9^x-1\sim x\log9$ e $(1+8x)^{1/4}-1\sim2x$; il rapporto vale $\log9/2=\log3$.\par\medskip
\noindent\textbf{quiz\_pagine\_221\_230\_012}: C. Sulla successione $x=k\pi$ vale zero, mentre per $x=\pi/2+k\pi$ diverge a $+\infty$.\par\medskip
\noindent\textbf{quiz\_pagine\_221\_230\_013}: B. $f=x+1+o(1)$ e $g=2x+o(1)$, percio' $f-g=-x+1+o(1)\to-\infty$.\par\medskip
\noindent\textbf{quiz\_pagine\_221\_230\_031}: B. $e^x-1\sim x$ e $x^2+2x\sim2x$, quindi il rapporto tende a $2$.\par\medskip
\noindent\textbf{quiz\_pagine\_231\_240\_008}: B. Una funzione continua su un compatto assume massimo e minimo assoluti (Weierstrass).\par\medskip
\noindent\textbf{quiz\_pagine\_231\_240\_016}: D. In $x=-1$ i limiti delle due espressioni sono $2+k$ ed $e^2$; imponendo $2+k=e^2$ segue $k=e^2-2$.\par\medskip
\noindent\textbf{quiz\_pagine\_231\_240\_018}: B. Su $[5,10]$ vale il polinomio continuo e $f(5)=-4<0$, $f(10)=21>0$. Su $[-3,1]$ la funzione e' discontinua in $0$; negli altri intervalli non vi e' cambio stretto di segno agli estremi.\par\medskip
\noindent\textbf{quiz\_pagine\_241\_250\_001}: A. Per il teorema degli zeri esiste $c\in(-1,1)$ con $f(c)=0$. Nel testo originale la variabile quantificata e' $x$ mentre la funzione usa $c$: refuso conservato.\par\medskip
\noindent\textbf{quiz\_pagine\_241\_250\_004}: C. $1/x$ non e' definita in $x=0$ e dunque non e' continua su tutto $[-1,1]$.\par\medskip
\noindent\textbf{quiz\_pagine\_251\_260\_010}: E. $4-e^{-2x}=0$ se $x=-\ln4/2=-\ln2\in[-1,0]$.\par\medskip
\noindent\textbf{quiz\_pagine\_271\_280\_005}: C. $9^x-1\sim x\log9$ e $(1+8x)^{1/4}-1\sim2x$; il rapporto tende a $\log9/2=\log3$.\par\medskip
\noindent\textbf{quiz\_pagine\_271\_280\_006}: C. Per $x=n\pi$ il prodotto vale zero; per $x=\pi/2+n\pi$ tende a $+\infty$.\par\medskip
\noindent\textbf{quiz\_pagine\_271\_280\_015}: D. $x^4+x^2=x^2(1+x^2)$ e $1+x^2\to1$.\par\medskip
\noindent\textbf{quiz\_pagine\_271\_280\_016}: C. Posto $t=\sqrt[3]x$, il numeratore e' $e^{2t}-1-2t=O(t^2)=O(|x|^{2/3})$; il denominatore e' asintotico a $\sqrt[7]x$. Il rapporto tende a zero.\par\medskip
\noindent\textbf{quiz\_pagine\_291\_300\_013}: A. La risposta ufficiale e' A, ma anche B e' vera: per $a\in[-32,32]$ l'immagine resta $[-32,32]$. Quesito con piu' risposte corrette.\par\medskip
\noindent\textbf{quiz\_pagine\_291\_300\_025}: D. Il limite sinistro in zero e' $10-1=9$; il valore destro e' $\beta$. Serve $\beta=9$, assente dalle alternative numeriche.\par\medskip
\noindent\textbf{quiz\_pagine\_301\_310\_008}: E. La definizione di limite impone $f(x)<0$ in un intorno bucato di $2$; non vincola il valore $f(2)$.\par\medskip
\noindent\textbf{quiz\_pagine\_301\_310\_011}: A. Il dominio e' $\mathbb R\setminus\{0\}$ e $f(-x)=\sin(\log|-x|)=f(x)$.\par\medskip
\noindent\textbf{quiz\_pagine\_301\_310\_020}: B. Il teorema degli zeri garantisce almeno una radice interna quando la funzione e' continua su tutto l'intervallo chiuso e i valori agli estremi hanno segno opposto.\par\medskip
\noindent\textbf{quiz\_pagine\_301\_310\_021}: C. Occorre $\cosh x>4$, ossia $x<-\operatorname{arcosh}4$ oppure $x>\operatorname{arcosh}4$: unione di due intervalli disgiunti.\par\medskip
\noindent\textbf{quiz\_pagine\_301\_310\_022}: B. La funzione e' strettamente decrescente per $x>0$, quindi e' iniettiva su $[0,+\infty)$; su intervalli simmetrici non lo e'.\par\medskip
\noindent\textbf{quiz\_pagine\_321\_330\_042}: C. Il prodotto negativo implica segni opposti agli estremi: per il teorema degli zeri esiste almeno un punto con $f(x)=0$.\par\medskip
\noindent\textbf{quiz\_pagine\_341\_350\_005}: B. $2^x-1\sim x\log2$ e $\sin3x\sim3x$, quindi il limite vale $\log2/3$.\par\medskip
\noindent\textbf{quiz\_pagine\_341\_350\_007}: C. Da $f(x)/x^2\to1$ segue $f(x)/x=(f(x)/x^2)x\to0$.\par\medskip
\noindent\textbf{quiz\_pagine\_341\_350\_033}: C. $\sqrt[4]{x^3+x^4}=x(1+1/x)^{1/4}\sim x$, dunque $f(x)\sim3x$.\par\medskip
\noindent\textbf{quiz\_pagine\_351\_360\_005}: B. $2^x-1\sim x\log2$ e $\sin3x\sim3x$, quindi il limite vale $\log2/3$.\par\medskip
\noindent\textbf{quiz\_pagine\_351\_360\_007}: C. Da $f(x)/x^2\to1$ segue $f(x)/x=(f(x)/x^2)x\to0$.\par\medskip
\noindent\textbf{quiz\_pagine\_351\_360\_033}: C. $\sqrt[4]{x^3+x^4}=x(1+1/x)^{1/4}\sim x$, dunque $f(x)\sim3x$.\par\medskip
\noindent\textbf{quiz\_pagine\_411\_420\_016}: C. Per applicare Weierstrass l\textbackslash{}'intervallo deve essere compatto non vuoto e non contenere $1$: dunque $1<a\leq10$. L\textbackslash{}'opzione C esclude $a=10$ (intervallo degenere); l\textbackslash{}'opzione E include $a>10$ (intervallo non valido).\par\medskip
\noindent\textbf{quiz\_pagine\_491\_500\_002}: D. Il teorema di prolungamento della derivata garantisce $f\prime(2)=\lim_{x\to2}f\prime(x)=5$, usando la continuit\textbackslash{}`a in $2$.\par\medskip
\noindent\textbf{quiz\_pagine\_491\_500\_003}: B. Ponendo $t=x+1$, $f=t-2t^2+O(t^3)$; $f(-1)=0$ e $f\prime(-1)=1$, quindi $y=x+1$.\par\medskip
\noindent\textbf{quiz\_pagine\_491\_500\_011}: E. $f\prime(0)=3+1=4$ e $f^{-1}(1)=0$; dunque $(f^{-1})\prime(1)=1/4$.\par\medskip
\noindent\textbf{quiz\_pagine\_491\_500\_014}: D. ATTENZIONE: la tabella ufficiale indica A, ma Rolle applicato a $h=f-g$ assicura soltanto $f\prime(c)=g\prime(c)$, non che entrambe siano zero. Controesempio: $f=g=x$, con derivate sempre uguali a $1$. D \textbackslash{}`e corretta.\par\medskip
\noindent\textbf{quiz\_pagine\_491\_500\_015}: E. In $(0,1)$, $(2-x)>0$, $(x-5)<0$ e $\arctan x>0$, quindi $f(x)<0$ e l' integrale \textbackslash{}`e negativo. Rolle assicura A e C; Lagrange su $[0,1]$ assicura D.\par\medskip
\noindent\textbf{quiz\_pagine\_491\_500\_016}: A. La continuit\textbackslash{}`a in $0$ impone $a=1$; in quel caso le derivate laterali sono entrambe zero. La funzione non \textbackslash{}`e globalmente limitata a sinistra.\par\medskip
\noindent\textbf{quiz\_pagine\_501\_510\_013}: B. Il teorema degli zeri richiede continuit\textbackslash{}`a sull\textbackslash{}'intervallo chiuso e valori agli estremi di segno opposto.\par\medskip
\subsubsection*{3.3 -- La continuità uniforme}
\noindent\textbf{quiz\_pagine\_261\_270\_025}: B. $h(x)=(a-c)x+(b-d)+o(1)$; l'asintoto e' obliquo esattamente quando $a-c\neq0$.\par\medskip
\noindent\textbf{quiz\_pagine\_271\_280\_007}: B. $f(x)=x+1+o(1)$ e $g(x)=2x+o(1)$: la differenza e' $-x+1+o(1)\to-\infty$.\par\medskip
\noindent\textbf{quiz\_pagine\_301\_310\_003}: D. Per $x\to-\infty$, $\log(3+4/x^2)=\log3+o(1)$; dunque $f(x)=2x+\log3+o(1)$.\par\medskip
\noindent\textbf{quiz\_pagine\_311\_320\_001}: B. Per $x\to+\infty$, $\log(3+2e^x)=x+\log2+o(1)$; dunque $f(x)=7x+\log2+o(1)$.\par\medskip
\noindent\textbf{quiz\_pagine\_341\_350\_026}: B. Per $x\to+\infty$, $\log(e^{-x}+3)\to\log3$, $\sin x/x\to0$, $(1+1/x)^x\to e$. Resta $3x+e+\log3$.\par\medskip
\noindent\textbf{quiz\_pagine\_351\_360\_026}: B. Per $x\to+\infty$, $\log(e^{-x}+3)\to\log3$, $\sin x/x\to0$, $(1+1/x)^x\to e$. Resta $3x+e+\log3$.\par\medskip
\noindent\textbf{quiz\_pagine\_431\_440\_006}: B. La frazione \textbackslash{}`e $x+1+o(1)$; per $x\to-\infty$, $\sqrt{x^2+6x}=-x-3+o(1)$. Quindi $f(x)=-2+o(1)$.\par\medskip
\noindent\textbf{quiz\_pagine\_461\_470\_003}: D. Dalla transitivit\textbackslash{}`a di $o$ piccolo, $f=o(g)$ e $g=o(x^5)$ implicano $f=o(x^5)$.\par\medskip
\noindent\textbf{quiz\_pagine\_461\_470\_004}: A. Poich\textbackslash{}'e $g\to0$, si ha $g^2=o(g)$. Pertanto $f=(f+g^2)-g^2=o(g)$.\par\medskip
\noindent\textbf{quiz\_pagine\_461\_470\_011}: E. Per $x>0$, $\sqrt{x^4+x^6}=x^2\sqrt{1+x^2}=x^2+\frac12x^4+o(x^4)$; il rapporto tende a $1/2$.\par\medskip
\noindent\textbf{quiz\_pagine\_481\_490\_010}: E. $f(x^2)=f(0)+f\prime(0)x^2+o(x^2)$, assumendo la regolarit\textbackslash{}`a necessaria. La risposta D differisce per il fattore $1/2$.\par\medskip
\noindent\textbf{quiz\_pagine\_481\_490\_013}: D. $\sin(x^6)=x^6+o(x^6)$ e $\log(2+u)=\log2+u/2+o(u)$.\par\medskip
\noindent\textbf{quiz\_pagine\_481\_490\_023}: D. Per $x>0$, $\sqrt{x^2+x^3}=x\sqrt{1+x}=x+\frac12x^2+O(x^3)$. Sottraendo $x+x^3$ rimane $\frac12x^2+O(x^3)$.\par\medskip
\subsubsection*{3.4 -- Confronto tra funzioni: infinitesimi e infiniti}
\noindent\textbf{quiz\_pagine\_001\_010\_012}: D. Spiegazione non ancora disponibile.\par\medskip
\noindent\textbf{quiz\_pagine\_001\_010\_017}: E. Spiegazione non ancora disponibile.\par\medskip
\noindent\textbf{quiz\_pagine\_001\_010\_033}: B. Spiegazione non ancora disponibile.\par\medskip
\noindent\textbf{quiz\_pagine\_001\_010\_039}: C. Spiegazione non ancora disponibile.\par\medskip
\noindent\textbf{quiz\_pagine\_011\_020\_016}: non verificata / nessuna alternativa corretta. Dall\textbackslash{}'ipotesi segue solo $f(x)/x^3\to0$; nessuna delle opzioni qui trascritte lo esprime. La formula dell\textbackslash{}'opzione (c) nel PDF va ricontrollata.\par\medskip
\noindent\textbf{quiz\_pagine\_021\_030\_036}: C. $f(x)/x=[f(x)/x^2]x\to0$.\par\medskip
\noindent\textbf{quiz\_pagine\_031\_040\_009}: D. Spiegazione non ancora disponibile.\par\medskip
\noindent\textbf{quiz\_pagine\_041\_050\_018}: C. Risposta ufficiale indicata nella tabella della decima simulazione.\par\medskip
\noindent\textbf{quiz\_pagine\_051\_060\_010}: C. Lo sviluppo $e^x-\sin x=1+O(x^2)$ mostra che il termine costante deve annullarsi.\par\medskip
\noindent\textbf{quiz\_pagine\_071\_080\_004}: D. Da $f(x)/e^x\to1$ segue $\log f(x)=x+\log(f(x)/e^x)=x+o(1)$, quindi $\log f(x)\sim x$.\par\medskip
\noindent\textbf{quiz\_pagine\_071\_080\_012}: E. Per $x\to0^+$, $\sqrt[4]{x^3+x^4}=x^{3/4}(1+x)^{1/4}\sim x^{3/4}$ e $2x=o(x^{3/4})$.\par\medskip
\noindent\textbf{quiz\_pagine\_071\_080\_019}: non verificata / nessuna alternativa corretta. L'equivalenza $f(x)/x\to1$ per $x\to+\infty$ implica, ponendo $t=1/x$, che $f(1/t)/(1/t)\to1$ per $t\to0^+$; l'affermazione con limite bilatero richiede un'ulteriore ipotesi.\par\medskip
\noindent\textbf{quiz\_pagine\_081\_090\_023}: non verificata / nessuna alternativa corretta. Poiché $f(x)/x=[f(x)/x^2]\,x\to0$, segue $f=o(x)$. Anche l\textbackslash{}textquotesingle alternativa D è vera, perché per $x\to+\infty$ vale $f(1/x)=o(1/x^2)$ e quindi anche $o(x^2)$: pertanto anche D è matematicamente vera.\par\medskip
\noindent\textbf{quiz\_pagine\_091\_100\_008}: A. $\log(1+3x)=3x-\frac92x^2+o(x^2)$, dunque il numeratore è $9x^2+4x^3+o(x^2)$ e il quoziente tende a $0$.\par\medskip
\noindent\textbf{quiz\_pagine\_091\_100\_010}: A. $\sqrt[3]{x^3+x^4}=x(1+x)^{1/3}=x+\frac13x^2+o(x^2)$; quindi $f(x)=\frac23x^2+o(x^2)$.\par\medskip
\noindent\textbf{quiz\_pagine\_091\_100\_013}: C. $\sinh^2(4x)\sim16x^2$ e $1-\cos(3x)\sim9x^2/2$, dunque il limite è $32/9$.\par\medskip
\noindent\textbf{quiz\_pagine\_091\_100\_028}: B. Per $x\to0$, $e^{\sqrt3x}-1\sim\sqrt3x$ e $\sin(\sqrt3x)\sim\sqrt3x$; il rapporto tende a $1$.\par\medskip
\noindent\textbf{quiz\_pagine\_101\_110\_002}: B. Per $x\to0$ la funzione tende a $2$, dunque $f(x)\sim2$.\par\medskip
\noindent\textbf{quiz\_pagine\_111\_120\_004}: E. L'equivalenza $f(x)/(3x)\to1$ non implica differenza infinitesima ne' equivalenza degli esponenziali. L'opzione d) non specifica il limite laterale necessario.\par\medskip
\noindent\textbf{quiz\_pagine\_111\_120\_010}: B. Da $\log(1+3x)=3x-\frac92x^2+9x^3+o(x^3)$ segue $f(x)=-18x^3+o(x^3)$.\par\medskip
\noindent\textbf{quiz\_pagine\_111\_120\_031}: E. Il radicando e' equivalente a $3x$ per $x\to0$; essendo la radice settima reale e continua, $f(x)\sim\sqrt[7]{3x}$.\par\medskip
\noindent\textbf{quiz\_pagine\_111\_120\_038}: A. Dividendo per $x>0$ si ha $1+1/x<f(x)/x<1+2/x$, quindi $f(x)/x\to1$.\par\medskip
\noindent\textbf{quiz\_pagine\_121\_130\_002}: D. Poiche\textbackslash{}'e $\sqrt[3]{x^3+x^4}=x^{4/3}(1+1/x)^{1/3}\sim x^{4/3}$ e $x=o(x^{4/3})$, si ha $f(x)\sim-x^{4/3}$.\par\medskip
\noindent\textbf{quiz\_pagine\_121\_130\_018}: A. Dividendo per $x^2$ per $x\neq0$ si ottiene $|f(x)/x^2|<\varepsilon$ in un intorno di $0$: e\textbackslash{}'e la definizione di $f=o(x^2)$.\par\medskip
\noindent\textbf{quiz\_pagine\_131\_140\_006}: E. $f(x)=\frac{2+\sin x}{x^2}(1+o(1))\to0$; ammette asintoto orizzontale $y=0$, ma $x^2f(x)$ oscilla tra $1$ e $3$.\par\medskip
\noindent\textbf{quiz\_pagine\_131\_140\_020}: D. $\log(1+x)=x-x^2/2+O(x^3)$ e $\sin^2\sqrt{x}=x-x^2/3+O(x^3)$. Il prodotto e\textbackslash{}'e $x^2-5x^3/6+O(x^4)$; sottraendo $x^2$ rimane $-5x^3/6$.\par\medskip
\noindent\textbf{quiz\_pagine\_131\_140\_025}: C. $f(x)\sim x^2$ e $g(x)\sim x$; dunque $f(x)/g(x)\sim x\to0$.\par\medskip
\noindent\textbf{quiz\_pagine\_141\_150\_007}: E. Dal prodotto $f/h^3=(f/g)(g/h^3)$ si ottiene $0\cdot1=0$.\par\medskip
\noindent\textbf{quiz\_pagine\_141\_150\_010}: E. Per $x<0$, $\sqrt{4x^2+3x^3}=|x|\sqrt{4+3x}\sim-2x$. Dunque $f(x)\sim x$.\par\medskip
\noindent\textbf{quiz\_pagine\_141\_150\_018}: D. $f/x^5=(f/g)(g/x^5)\to0$. Non e\textbackslash{}'e invece garantito l\textbackslash{}'ordine $x^{11}$ del prodotto.\par\medskip
\noindent\textbf{quiz\_pagine\_141\_150\_019}: A. Da $f=o(g)-g^2$ e $g^2=o(g)$ segue $f=o(g)$.\par\medskip
\noindent\textbf{quiz\_pagine\_171\_180\_005}: A. $\log(e^{2x}+x)=2x+\log(1+xe^{-2x})=2x+o(1)$.\par\medskip
\noindent\textbf{quiz\_pagine\_171\_180\_014}: D. Per $x>0$, $\sqrt{x^2+x^3}=x\sqrt{1+x}=x+x^2/2+O(x^3)$.\par\medskip
\noindent\textbf{quiz\_pagine\_171\_180\_028}: A. Per $x>0$, $f=x+x\sqrt{1+x^2}\sim2x$. Per $x<0$, $f=x-x\sqrt{1+x^2}=-x^3/2+O(x^5)$.\par\medskip
\noindent\textbf{quiz\_pagine\_171\_180\_033}: D. $f-g=(9\sin x+1)/(x^2+1)\to0$ e $g(x)\sim3x$, percio' $f/g\to1$.\par\medskip
\noindent\textbf{quiz\_pagine\_181\_190\_005}: A. Poiche' $\sinh x\sim e^x/2$, il rapporto $(e^x/x)/\sinh x\sim2/x\to0$.\par\medskip
\noindent\textbf{quiz\_pagine\_181\_190\_012}: D. $\sqrt[7]{x+x^2}=x^{1/7}(1+x)^{1/7}=x^{1/7}+\frac17x^{8/7}+o(x^{8/7})$ per $x\to0^+$. Poiche' $7/8<8/7$, domina $-x^{7/8}$. Il testo indica $x\to0$, ma la potenza $x^{7/8}$ e' reale solo per $x\geq0$.\par\medskip
\noindent\textbf{quiz\_pagine\_181\_190\_013}: E. Per $x\to0^+$, $\sqrt[7]{x+x^2}=x^{1/7}+\frac17x^{8/7}+o(x^{8/7})$. Sommando $-x^{8/7}$ si ottiene $-\frac87x^{8/7}+o(x^{8/7})$.\par\medskip
\noindent\textbf{quiz\_pagine\_211\_220\_036}: D. Per definizione $f=o(g)$ significa $f/g\to0$ dove il rapporto e' definito.\par\medskip
\noindent\textbf{quiz\_pagine\_211\_220\_042}: C. Usando $\sin^2u=u^2-u^4/3+o(u^4)$ e $u=x-x^2$, resta $(2/3)x^4+o(x^4)$.\par\medskip
\noindent\textbf{quiz\_pagine\_211\_220\_043}: B. Il polinomio di Taylor mostra $f(0)=-2$, $f'(0)=0$, $f''(0)=4$, $f'''(0)=-18$.\par\medskip
\noindent\textbf{quiz\_pagine\_211\_220\_046}: A. $\log\cosh x=x^2/2-x^4/12+O(x^6)$; quindi $f(x)=-3x^3-x^4/3+O(x^6)\sim-3x^3$.\par\medskip
\noindent\textbf{quiz\_pagine\_221\_230\_024}: D. Poiche' $x^4=o(x^2)$, si ha $x^4+x^2=x^2(1+x^2)\sim x^2$.\par\medskip
\noindent\textbf{quiz\_pagine\_241\_250\_021}: A. Per $\alpha=-1$ si annulla il termine lineare: $\sqrt[4]{x^4+x^3}-x\to1/4$, dunque non vale la P.P. lineare per ogni $\alpha$.\par\medskip
\noindent\textbf{quiz\_pagine\_241\_250\_024}: E. $\sin x-\tan x=(x-x^3/6)-(x+x^3/3)+o(x^3)=-x^3/2+o(x^3)$, dunque la parte principale non e' zero.\par\medskip
\noindent\textbf{quiz\_pagine\_281\_290\_022}: B. $\sqrt[3]{x^3+x^6}=x(1+x^3)^{1/3}=x+\frac13x^4+o(x^4)$. Sommando $-x+x^4+x^6$ si ha $\frac43x^4+o(x^4)$.\par\medskip
\noindent\textbf{quiz\_pagine\_281\_290\_027}: E. Nel numeratore domina $\sqrt x\log^2x$ e nel denominatore $\sqrt x\log x$, dunque il rapporto e' asintotico a $\log x\to-\infty$.\par\medskip
\noindent\textbf{quiz\_pagine\_291\_300\_012}: D. $\sin(x+x^{6/5})=x+x^{6/5}+o(x)$, dunque $f(x)\sim5x$ (nel dominio reale della potenza).\par\medskip
\noindent\textbf{quiz\_pagine\_291\_300\_017}: E. $f(x)\log x\sim(\log x)^2=o(x^2)$. L'equivalenza non implica che $f-\log x$ tenda a zero.\par\medskip
\noindent\textbf{quiz\_pagine\_291\_300\_023}: E. Ponendo $t=1/x\to+\infty$, da $g(t)/t\to1$ segue $xg(1/x)\to1$.\par\medskip
\noindent\textbf{quiz\_pagine\_291\_300\_027}: D. $1-\cos^2(1/x)=\sin^2(1/x)\sim x^{-2}$, mentre $(1+x^5)^2\sim x^{10}$: prodotto $\sim x^8$.\par\medskip
\noindent\textbf{quiz\_pagine\_291\_300\_029}: B. $x^{8/7}$ domina $x^{2/7}$ e $x^{1/7}$, dunque $f(x)\sim x^{8/7}$.\par\medskip
\noindent\textbf{quiz\_pagine\_331\_340\_009}: A. $(1+3x)^{1/3}-1\sim x$ e $7x+20x^4\sim7x$: limite $1/7$.\par\medskip
\noindent\textbf{quiz\_pagine\_331\_340\_010}: C. Poich\textbackslash{}`e $f(x)/x^2\to0$, si ha $f(x)/x=x[f(x)/x^2]\to0$.\par\medskip
\noindent\textbf{quiz\_pagine\_361\_370\_006}: B. RISPOSTA UFFICIALE B, MA IL TESTO NON IMPONE CHE I LIMITI DI $f$ E $g$ SIANO ZERO. Il teorema di de l\textbackslash{}textquotesingle Hopital richiede una forma indeterminata $0/0$ (o infinito/infinito). Controesempio: $f(x)=1+2x$, $g(x)=1+x$, $x_0=0$: il rapporto delle derivate vale $2$, quello delle funzioni tende a $1$. Nessuna alternativa e garantita come scritta.\par\medskip
\noindent\textbf{quiz\_pagine\_361\_370\_010}: C. $e^x-\sin x=1+O(x^2)$; per $\lambda=1$ si ha $f(x)=O(x^2)=o(x)$.\par\medskip
\noindent\textbf{quiz\_pagine\_361\_370\_023}: A. $x-\log(1+x)=x^2/2+o(x^2)$ e $x\sin3x=3x^2+o(x^2)$; il limite e $1/6$.\par\medskip
\noindent\textbf{quiz\_pagine\_371\_380\_019}: D. $f(x)=e^x(1+o(1))$ implica $\log f(x)=x+\log(1+o(1))=x+o(1)$, quindi $\log f(x)\sim x$.\par\medskip
\noindent\textbf{quiz\_pagine\_371\_380\_027}: E. Per $x\to0^+$, $\sqrt[4]{x^3+x^4}=x^{3/4}(1+x)^{1/4}\sim x^{3/4}$, che domina $2x$.\par\medskip
\noindent\textbf{quiz\_pagine\_371\_380\_028}: A. $e^{2x}-1\sim2x$ e $\log(1+3x)\sim3x$, dunque il limite e $2/3$.\par\medskip
\noindent\textbf{quiz\_pagine\_381\_390\_004}: D. Sostituendo $t=1/x$ si ha $t\to+\infty$ per $x\to0^+$, dunque $f(1/x)\sim1/x$ da destra. Il testo scrive $x\to0$ senza specificare il lato: la conclusione bilaterale non e garantita, quindi la risposta ufficiale D richiede $x\to0^+$.\par\medskip
\noindent\textbf{quiz\_pagine\_381\_390\_006}: E. $\sin3x\sim3x$ e $\cosh5x-1\sim25x^2/2$, quindi il rapporto e equivalente a $6/(25x)$: i limiti laterali hanno segni opposti.\par\medskip
\noindent\textbf{quiz\_pagine\_381\_390\_010}: B. Il denominatore resta tra $4$ e $5$ mentre $\log(5x)\to+\infty$; il rapporto tende a $+\infty$.\par\medskip
\noindent\textbf{quiz\_pagine\_391\_400\_008}: A. $f(x)/x=(f(x)/x^2)x\to0$. Non segue un ordine di annullamento maggiore di due.\par\medskip
\noindent\textbf{quiz\_pagine\_391\_400\_012}: A. $(4^x+2^x)^{1/x}=4(1+2^{-x})^{1/x}\to4$.\par\medskip
\noindent\textbf{quiz\_pagine\_411\_420\_002}: B. Per $x\to0$ la funzione tende a $2$, quindi $f(x)\sim2$. Il commento nella fonte che cita $4$ come parte principale \textbackslash{}`e errato.\par\medskip
\noindent\textbf{quiz\_pagine\_411\_420\_005}: A. Se $x\to+\infty$, allora $x^2\to+\infty$ e per ipotesi $f(x^2)\to0$; il segno non \textbackslash{}`e determinato.\par\medskip
\noindent\textbf{quiz\_pagine\_411\_420\_024}: E. $f(x)/(3x)\to1$ non implica $f(x)-3x\to0$, n\textbackslash{}'e l\textbackslash{}'equivalenza degli esponenziali. Le altre conclusioni non seguono.\par\medskip
\noindent\textbf{quiz\_pagine\_411\_420\_030}: B. $\log(1+3x)=3x-9x^2/2+9x^3+o(x^3)$; i termini lineari e quadratici si cancellano e rimane $-18x^3$.\par\medskip
\noindent\textbf{quiz\_pagine\_421\_430\_014}: E. All\textbackslash{}'interno della radice il termine dominante \textbackslash{}`e $3x$; la radice settima \textbackslash{}`e definita anche per $x<0$.\par\medskip
\noindent\textbf{quiz\_pagine\_421\_430\_021}: A. Dividendo per $x>0$, $1+1/x<f(x)/x<1+2/x$; per confronto $f(x)/x\to1$.\par\medskip
\noindent\textbf{quiz\_pagine\_431\_440\_013}: E. $x-\sin x\sim x^3/6$, dunque il limite \textbackslash{}`e $0$ per $\alpha<3$, $1/6$ per $\alpha=3$ e $+\infty$ per $\alpha>3$.\par\medskip
\noindent\textbf{quiz\_pagine\_431\_440\_017}: D. $\sqrt[3]{x^3+x^4}\sim x^{4/3}$ domina $2x$, quindi $f(x)\sim-x^{4/3}$.\par\medskip
\noindent\textbf{quiz\_pagine\_431\_440\_033}: A. \textbackslash{}`E esattamente la definizione di $f(x)=o(x^2)$ per $x\to0$ (per $x\ne0$).\par\medskip
\noindent\textbf{quiz\_pagine\_441\_450\_008}: C. $f\sim x^2$, $g\sim-x^3$; $f^2/g\sim-x\to0$. Le opzioni B ed E non valgono per limite bilatero.\par\medskip
\noindent\textbf{quiz\_pagine\_441\_450\_013}: D. Numeratore $x(-x^2/2)+x^3=x^3/2+O(x^5)$, denominatore $x^3-x^4+o(x^4)$; rapporto $\to1/2$.\par\medskip
\noindent\textbf{quiz\_pagine\_451\_460\_016}: E. Il logaritmo del limite \textbackslash{}`e $\lim\log(1+4x^2)/\log x=2$, quindi il limite vale $e^2$.\par\medskip
\noindent\textbf{quiz\_pagine\_451\_460\_017}: C. $f(x)\sim x^2$ e $g(x)\sim x$; dunque $f(x)/g(x)\sim x\to0$.\par\medskip
\noindent\textbf{quiz\_pagine\_451\_460\_018}: B. Poich\textbackslash{}`e $\sin(1/x)\sim1/x$, il prodotto \textbackslash{}`e asintotico a $3x^2-4x\to+\infty$.\par\medskip
\noindent\textbf{quiz\_pagine\_451\_460\_024}: A. Da $f/g\to0$ e $g/h^3\to1$ segue $f/h^3\to0$. ATTENZIONE: anche E afferma esattamente la definizione di A, quindi due risposte risultano equivalenti nel testo originale.\par\medskip
\noindent\textbf{quiz\_pagine\_471\_480\_009}: E. La condizione $f(x)/x\to4$ equivale a $f(x)/(4x)\to1$, cio\textbackslash{}`e $f(x)\sim4x$. Non implica un asintoto.\par\medskip
\noindent\textbf{quiz\_pagine\_471\_480\_018}: A. La funzione $\arctan(x+x^{3/2})$ tende a $\pi/2$, dunque $f(x)=x+O(1)\sim x$.\par\medskip
\noindent\textbf{quiz\_pagine\_471\_480\_019}: A. Per $x\to0^+$, $\arctan(x+x^{3/2})=x+x^{3/2}+O(x^3)$, dunque $f(x)\sim2x$. ATTENZIONE: per $x<0$ la potenza $x^{3/2}$ non \textbackslash{}`e reale; il limite va inteso destro.\par\medskip
\noindent\textbf{quiz\_pagine\_471\_480\_023}: B. Poich'e $x=o(x^{2/3})$ e $\arctan u\sim u$ per $u\to0$, si ha $f(x)\sim x^{2/3}$.\par\medskip
\noindent\textbf{quiz\_pagine\_481\_490\_001}: A. Ponendo $t=1/x$, si ha $\cos t-e^{4t}=-4t+O(t^2)$, dunque il limite vale $-4$.\par\medskip
\noindent\textbf{quiz\_pagine\_481\_490\_003}: E. Il logaritmo della potenza \textbackslash{}`e $5\log(n+3)/\log n\to5$; per continuit\textbackslash{}`a della funzione esponenziale il limite \textbackslash{}`e $e^5$.\par\medskip
\noindent\textbf{quiz\_pagine\_481\_490\_005}: B. Con $t=5x-4=5(x-4/5)$, il limite \textbackslash{}`e $25\lim_{t\to0}(1-\cos t)/t^2=25/2$.\par\medskip
\noindent\textbf{quiz\_pagine\_481\_490\_014}: A. $\log(e^{2x}+x)=2x+\log(1+xe^{-2x})=2x+o(1)$; quindi il limite vale $6$.\par\medskip
\noindent\textbf{quiz\_pagine\_481\_490\_017}: A. Come nel quesito 20 della 28a simulazione, $1-\cos t\sim t^2/2$ con $t=5(x-4/5)$, dunque $25/2$.\par\medskip
\noindent\textbf{quiz\_pagine\_481\_490\_020}: B. $\log(x^{-1/2})=-\frac12\log x\to-\infty$, mentre $1\le2-\cos^2x\le2$. Il quoziente tende a $-\infty$.\par\medskip
\noindent\textbf{quiz\_pagine\_491\_500\_007}: B. $5^x-1\sim x\log5$ e $\sin4x\sim4x$.\par\medskip
\noindent\textbf{quiz\_pagine\_491\_500\_012}: A. $\sqrt{x^2+x^4}=|x|\sqrt{1+x^2}$. Per $x>0$ si ha $2x+O(x^3)$; per $x<0$, $x-x(1+x^2/2+O(x^4))=-x^3/2+O(x^5)$.\par\medskip
\noindent\textbf{quiz\_pagine\_491\_500\_017}: D. $f-g=(9\sin x+1)/(x^2+1)\to0$; poich\textbackslash{}'e $g(x)\to+\infty$, segue $f/g\to1$.\par\medskip
\noindent\textbf{quiz\_pagine\_501\_510\_004}: E. Per $x>0$, $\sqrt{x^4+x^6}=x^2\sqrt{1+x^2}$; razionalizzando si ottiene $1/(\sqrt{1+x^2}+1)\to1/2$.\par\medskip
\noindent\textbf{quiz\_pagine\_501\_510\_011}: C. Usando $\sin^2t=t^2-t^4/3+o(t^4)$ con $t=x-x^2$, si trova $x^2-2x^3+(1-1/3)x^4+o(x^4)$; i termini fino al cubico si cancellano.\par\medskip
\noindent\textbf{quiz\_pagine\_501\_510\_012}: B. Confrontando con Taylor si ha $f(0)=-2$, $f\prime(0)=0$, $f\prime\prime(0)=4$, $f\prime\prime\prime(0)=-18$. Inoltre $x=0$ \textbackslash{}`e un minimo locale.\par\medskip
\noindent\textbf{quiz\_pagine\_501\_510\_015}: A. $\log\cosh x=x^2/2-x^4/12+O(x^6)$, quindi $4\log\cosh x-2x^2-3x^3=-3x^3-x^4/3+O(x^6)$; parte principale $-3x^3$.\par\medskip
\noindent\textbf{quiz\_pagine\_501\_510\_031}: B. $\cos(x^3)=1-x^6/2+O(x^{12})$; quindi $e^{\cos(x^3)}=e(1-x^6/2+O(x^{12}))$. Le alternative (d) ed (e) risultano identiche nel PDF.\par\medskip
\noindent\textbf{quiz\_pagine\_501\_510\_032}: E. Per $x$ vicino a $2$ e diverso da $2$, $f(x)-f(2)=-3(x-2)^6+o((x-2)^6)<0$: massimo locale stretto.\par\medskip
\noindent\textbf{quiz\_pagine\_501\_510\_035}: A. Per $x>0$, $\sqrt{4x^2+2x^3}=x\sqrt{4+2x}\sim2x$; quindi $f(x)\sim5x$.\par\medskip
\noindent\textbf{quiz\_pagine\_511\_520\_011}: B. Per $a>0$, $\log(1+2x^a)\sim2x^a$, quindi $f(x)\sim2x^{3a}$: $3a=5$ implica $a=5/3$.\par\medskip
\noindent\textbf{quiz\_pagine\_511\_520\_014}: B. $e^{1-\cos x}-1=x^2/2+O(x^4)$ e $x\sin(2x)=2x^2+O(x^4)$: $f(x)\sim-3x^2/2$.\par\medskip
\noindent\textbf{quiz\_pagine\_511\_520\_035}: D. $x^4+x^2=x^2(1+x^2)$ e il rapporto con $x^2$ tende a $1$.\par\medskip
\noindent\textbf{quiz\_pagine\_511\_520\_036}: C. $\sqrt[3]{8x}=2\sqrt[3]{x}$; lo sviluppo dell'esponenziale cancella il termine lineare, lasciando $O(|x|^{2/3})$. Il denominatore \textbackslash{}`e asintotico a $\sqrt[7]{x}$, dunque il rapporto tende a $0$.\par\medskip
\subsection*{Capitolo 4}
\subsubsection*{4.1 -- Derivate}
\noindent\textbf{quiz\_pagine\_001\_010\_001}: E. Spiegazione non ancora disponibile.\par\medskip
\noindent\textbf{quiz\_pagine\_001\_010\_005}: E. Spiegazione non ancora disponibile.\par\medskip
\noindent\textbf{quiz\_pagine\_001\_010\_011}: C. Spiegazione non ancora disponibile.\par\medskip
\noindent\textbf{quiz\_pagine\_011\_020\_008}: D. $f(x)=15\sin^2x$, $f\prime(x)=30\sin x\cos x$, che si annulla in $\pi/2$.\par\medskip
\noindent\textbf{quiz\_pagine\_021\_030\_009}: C. Per $a^{u(x)}$, la derivata \textbackslash{}`e $a^{u(x)}\log(a)u\prime(x)$.\par\medskip
\noindent\textbf{quiz\_pagine\_021\_030\_023}: B. Una primitiva \textbackslash{}`e $\log|f(x)|+C$ dove $f\ne0$; moltiplicare il logaritmo per $2$ ne raddoppia la derivata.\par\medskip
\noindent\textbf{quiz\_pagine\_031\_040\_002}: D. Spiegazione non ancora disponibile.\par\medskip
\noindent\textbf{quiz\_pagine\_031\_040\_005}: E. Spiegazione non ancora disponibile.\par\medskip
\noindent\textbf{quiz\_pagine\_031\_040\_006}: A. Spiegazione non ancora disponibile.\par\medskip
\noindent\textbf{quiz\_pagine\_041\_050\_007}: E. Risposta ufficiale indicata nella tabella della decima simulazione.\par\medskip
\noindent\textbf{quiz\_pagine\_041\_050\_015}: D. Risposta ufficiale indicata nella tabella della decima simulazione.\par\medskip
\noindent\textbf{quiz\_pagine\_041\_050\_032}: non verificata / nessuna alternativa corretta. Soluzione non riportata nel PDF; da verificare.\par\medskip
\noindent\textbf{quiz\_pagine\_051\_060\_016}: D. Il rapporto incrementale in zero è $\sin(1/x)$ e non ha limite.\par\medskip
\noindent\textbf{quiz\_pagine\_061\_070\_002}: D. Imponendo continuità e uguaglianza delle derivate in x=1 si ottiene alfa=4/5, beta=3/5.\par\medskip
\noindent\textbf{quiz\_pagine\_061\_070\_010}: D. Espandendo $e^u=1+u+\frac{u^2}{2}+o(x^2)$, con $u=2x+x^2$, si ottiene $1+2x+3x^2$. Le risposte A e B sono duplicate nel PDF originale.\par\medskip
\noindent\textbf{quiz\_pagine\_061\_070\_014}: A. Si usa log(1+u)=u-u\textsuperscript{2}/2+u\textsuperscript{3}/3+o(x\textsuperscript{3}), con u=x+x\textsuperscript{2}.\par\medskip
\noindent\textbf{quiz\_pagine\_061\_070\_022}: E. Sviluppando sqrt(1-x) e poi l'esponenziale si ottiene e(1-x/2-x\textsuperscript{3}/48+o(x\textsuperscript{3})); soluzione calcolata, tabella ufficiale non presente nell'estratto.\par\medskip
\noindent\textbf{quiz\_pagine\_071\_080\_001}: E. Poiché $g=f^{-2}$, si ha $g'(0)=-2f'(0)/f(0)^3=-6/(-8)=3/4$; la risposta ufficiale è E.\par\medskip
\noindent\textbf{quiz\_pagine\_071\_080\_008}: D. In $0$ entrambi i rami valgono $1$ e le derivate laterali sono entrambe $0$; altrove i rami sono regolari.\par\medskip
\noindent\textbf{quiz\_pagine\_071\_080\_010}: E. Per derivazione composta, $f'(x)=(6x+6x^2)/(3x^2+2x^3)=6(1+x)/(3x+2x^2)$, dove definita.\par\medskip
\noindent\textbf{quiz\_pagine\_071\_080\_016}: C. Per la derivata dell'inversa, se $f'(0)\ne0$ allora $m=1/f'(0)$: con $-1<f'(0)<0$ si ha $m<-1$. Se $f'(0)=0$, la tangente può essere verticale.\par\medskip
\noindent\textbf{quiz\_pagine\_081\_090\_002}: D. $f(e)=e$ e $f\prime(e)=2$, quindi $y=e+2(x-e)=2x-e$.\par\medskip
\noindent\textbf{quiz\_pagine\_081\_090\_008}: E. Per la regola della catena, $(g\circ f)\prime(-2)=g\prime(f(-2))f\prime(-2)=(-7)(-4)=28$. Le alternative A e D sono duplicate nel PDF.\par\medskip
\noindent\textbf{quiz\_pagine\_081\_090\_010}: A. La derivata di $a^x$ è $a^x\log a$ per $a>0$.\par\medskip
\noindent\textbf{quiz\_pagine\_081\_090\_011}: C. $g\prime=f\prime\prime\prime$ esiste, quindi $g$ è derivabile e perciò continua. Non è garantita la continuità di $f\prime\prime\prime$.\par\medskip
\noindent\textbf{quiz\_pagine\_081\_090\_015}: B. Per $x\neq0$, la derivata del primo addendo è $2/(4+x^2)$ e quella del secondo è $-2/(4+x^2)$; sottraendo si ottiene $4/(4+x^2)$.\par\medskip
\noindent\textbf{quiz\_pagine\_081\_090\_021}: D. Per $a=0$, i valori laterali in $0$ sono entrambi $0$ e le derivate laterali sono entrambe $1$; altrove i rami sono regolari.\par\medskip
\noindent\textbf{quiz\_pagine\_081\_090\_024}: non verificata / nessuna alternativa corretta. Per $\sin x<0$, $f\prime(x)=(-\cos x)/(-\sin x)=\cos x/\sin x=1/\tan x$. Le alternative A ed E sono equivalenti, quindi non esiste una sola risposta corretta.\par\medskip
\noindent\textbf{quiz\_pagine\_091\_100\_003}: E. La regola della catena dà $f\prime(x)=\frac13(e^{4x}\cos^8x)^{-2/3}e^{4x}(4\cos^8x-8\cos^7x\sin x)$, equivalente alla E, con prolungamento nei punti $\cos x=0$.\par\medskip
\noindent\textbf{quiz\_pagine\_091\_100\_016}: non verificata / nessuna alternativa corretta. Le derivate di ordine dispari di $\cos x$ si annullano in $0$: A è vera; anche D è vera perché descrive un sottoinsieme degli ordini dispari. Il testo ufficiale indica A, ma il quesito ha più risposte vere.\par\medskip
\noindent\textbf{quiz\_pagine\_091\_100\_017}: E. $f\prime(x)=-5e^{-5x}/(3-e^{-5x})^2=-5e^{5x}/(3e^{5x}-1)^2$.\par\medskip
\noindent\textbf{quiz\_pagine\_091\_100\_023}: A. Il teorema di prolungamento della derivata: continuità in $2$ e limite finito di $f\prime$ implicano derivabilità in $2$ con $f\prime(2)=4$.\par\medskip
\noindent\textbf{quiz\_pagine\_101\_110\_006}: E. Con la regola del quoziente, $f\prime(x)=\frac{-e^{1/x}/x- (e^{1/x}-1)}{x^2}=\frac{x-(x+1)e^{1/x}}{x^3}$.\par\medskip
\noindent\textbf{quiz\_pagine\_101\_110\_012}: E. Le derivate di ordine dispari del coseno sono proporzionali al seno, e quindi si annullano in $0$. Si intende $f^n$ come derivata di ordine $n$.\par\medskip
\noindent\textbf{quiz\_pagine\_101\_110\_014}: A. La continuità in $0$ richiede $a=\cos0=1$. Le derivate laterali sono entrambe nulle, dunque la funzione è derivabile anche in $0$.\par\medskip
\noindent\textbf{quiz\_pagine\_111\_120\_001}: C. Per la derivata della funzione inversa, $m=1/f'(0)$. Poiche' $f'(0)<-1$, risulta $-1<m<0$.\par\medskip
\noindent\textbf{quiz\_pagine\_111\_120\_007}: A. Si leggono $f(-2\pi)=\pi$ e $f'(-2\pi)=-2$; la tangente e' $y=\pi-2(x+2\pi)=-2x-3\pi$.\par\medskip
\noindent\textbf{quiz\_pagine\_111\_120\_012}: D. La derivata e' $f'(x)=3x^2\sin x+x^3\cos x$ e dunque $f'(0)=0$.\par\medskip
\noindent\textbf{quiz\_pagine\_111\_120\_027}: E. Si ha $f(x)=\log(\cos2x)+3x$ dove definita, quindi $f'(x)=3-2\tan2x$. L'opzione e) manca del fattore $2$.\par\medskip
\noindent\textbf{quiz\_pagine\_111\_120\_037}: A. La regola del prodotto da' $h'(0)=f'(0)g(0)+f(0)g'(0)=g(0)-3g'(0)$.\par\medskip
\noindent\textbf{quiz\_pagine\_111\_120\_039}: C. Per $a=2$ la funzione e' $(x-4)^2$, regolare; per $a>2$ la seconda derivata tende a $0$ nel punto $4$. Per $a<2$ la derivabilita' seconda in $4$ fallisce.\par\medskip
\noindent\textbf{quiz\_pagine\_121\_130\_007}: B. Nel dominio $\cos x>0$ vale $f(x)=4x+\log(\cos x)$; derivando si ottiene $4-\tan x$.\par\medskip
\noindent\textbf{quiz\_pagine\_131\_140\_008}: B. Il termine $2|x|\sin x$ ha rapporto incrementale $2|x|\sin x/x\to0$; la derivata in zero e\textbackslash{}'e $1$.\par\medskip
\noindent\textbf{quiz\_pagine\_131\_140\_010}: B. $\log|f(x)|=\log|x|+\arctan(1/x)$, quindi $f\prime/f=1/x-1/(x^2+1)=(x^2-x+1)/(x(x^2+1))$.\par\medskip
\noindent\textbf{quiz\_pagine\_131\_140\_012}: C. Per la regola della catena $g\prime(1)=e^{f(1)}f\prime(1)=-2e^3$.\par\medskip
\noindent\textbf{quiz\_pagine\_131\_140\_013}: D. La derivabilita\textbackslash{}'a assicura la tangente, orizzontale se $f\prime(x_1)=0$. Su domini non intervallari, $f\prime>0$ non garantisce monotonia globale.\par\medskip
\noindent\textbf{quiz\_pagine\_131\_140\_019}: A. Il rapporto incrementale e\textbackslash{}'e $((1+e^x)^2-4)/x=(e^x-1)(e^x+3)/x$; il suo limite per $x\to0$ vale $4$.\par\medskip
\noindent\textbf{quiz\_pagine\_141\_150\_003}: B. Scrivendo $f(x)=\exp(4\cos x\log(3\cos x))$, la derivata dell\textbackslash{}'esponente contiene il fattore $\sin x$ e si annulla in $0$.\par\medskip
\noindent\textbf{quiz\_pagine\_141\_150\_004}: A. I due rami tendono a $0$; le derivate laterali tendono entrambe a $2$. Anche la derivata e\textbackslash{}'e continua in $0$.\par\medskip
\noindent\textbf{quiz\_pagine\_141\_150\_005}: A. Si ha $f(0)=1$ e $f\prime(0)=1$; dunque $(f^{-1})\prime(1)=1/f\prime(0)=1$.\par\medskip
\noindent\textbf{quiz\_pagine\_141\_150\_023}: B. $f\prime=e^{\sinh x}(\cosh x\cos x-\sin x)$. Il termine tra parentesi cambia segno tra $-\pi/2$ e $\pi/2$, e ancora tra $\pi/2$ e $3\pi/2$: almeno due zeri.\par\medskip
\noindent\textbf{quiz\_pagine\_141\_150\_024}: C. La regola del quoziente da\textbackslash{}'e $g\prime(x)=(xf\prime(x)-f(x))/x^2$ e quindi $g\prime(1)=2-f(1)$.\par\medskip
\noindent\textbf{quiz\_pagine\_151\_160\_010}: D. Le derivate del seno sono periodiche di periodo $4$; all\textbackslash{}'argomento $\pi/2$ quelle di ordine multiplo di $4$ valgono $1$. L\textbackslash{}'alternativa (b) e\textbackslash{}'e formalmente vera se si ammettono derivate di ordine negativo, che pero\textbackslash{}'o non sono definite nel contesto.\par\medskip
\noindent\textbf{quiz\_pagine\_151\_160\_013}: E. Regola della catena: $(\log u)\prime=u\prime/u$ con $u=h^3+1$.\par\medskip
\noindent\textbf{quiz\_pagine\_151\_160\_026}: D. $f\prime(x)=3+1/(1+x^2)$, quindi $f\prime(0)=4$ e $(f^{-1})\prime(1)=1/4$.\par\medskip
\noindent\textbf{quiz\_pagine\_161\_170\_004}: D. Per la regola della catena $f'(x)=e^{\sin(1/(x+1))}\cos(1/(x+1))\cdot[-1/(x+1)^2]$.\par\medskip
\noindent\textbf{quiz\_pagine\_161\_170\_015}: A. Derivando il logaritmo si ottiene $(4x+9x^2)/(2x^2+3x^3)=(4+9x)/(x(2+3x))$.\par\medskip
\noindent\textbf{quiz\_pagine\_171\_180\_006}: D. Le derivate di ordine dispari di $\sinh x$ sono $\cosh x$, che in $0$ vale $1$; quelle pari sono $\sinh x$, che in $0$ vale $0$.\par\medskip
\noindent\textbf{quiz\_pagine\_171\_180\_015}: C. $f(0)=2$, $f'(0)=4$; la derivata dell'inversa e' $1/f'(0)=1/4$.\par\medskip
\noindent\textbf{quiz\_pagine\_171\_180\_018}: D. Per il teorema del valor medio applicato tra $2$ e $x$, il rapporto incrementale tende a $5$; per continuita' di $f$ in $2$ si conclude $f'(2)=5$.\par\medskip
\noindent\textbf{quiz\_pagine\_171\_180\_019}: B. Si ha $f(-1)=0$ e $f'(-1)=1$ (il termine cubico e' di ordine almeno $3$), quindi $y=x+1$.\par\medskip
\noindent\textbf{quiz\_pagine\_171\_180\_027}: E. $f'(0)=4$ e $f^{-1}(1)=0$, quindi $(f^{-1})'(1)=1/4$.\par\medskip
\noindent\textbf{quiz\_pagine\_171\_180\_032}: A. La continuita' in $0$ richiede $a=1$; allora entrambe le derivate laterali valgono $0$.\par\medskip
\noindent\textbf{quiz\_pagine\_181\_190\_009}: D. Il rapporto incrementale $[f(x)-f(0)]/x=f(x)/x\to0$, quindi $f'(0)=0$.\par\medskip
\noindent\textbf{quiz\_pagine\_181\_190\_014}: C. $f(0)=e^2$ e $f'(x)=2e^xe^{2e^x}$, quindi $f'(0)=2e^2$.\par\medskip
\noindent\textbf{quiz\_pagine\_181\_190\_019}: B. Vicino a $x=1$ i due fattori si comportano ciascuno come $\sqrt{|x-1|}$, quindi il prodotto e' proporzionale a $|x-1|$ (punto angoloso). Vicino a $x=2$ si ha $C\sqrt{|x-2|}$ (cuspide).\par\medskip
\noindent\textbf{quiz\_pagine\_181\_190\_020}: non verificata / nessuna alternativa corretta. Per $x<0$, $|x|=-x$ e $|x-3|=3-x$, quindi $f(x)=x\sqrt{3-x}+\sin(4x)$: la derivata sinistra e' $\sqrt3+4$, non $-4-\sqrt3$. Il PDF indica A ma nessuna alternativa coincide con il calcolo; controllare il segno originale della funzione.\par\medskip
\noindent\textbf{quiz\_pagine\_191\_200\_002}: A. La funzione $f(x)=x^3$ e' strettamente crescente e derivabile, ma $f^{-1}(y)=\sqrt[3]{y}$ non e' derivabile in $0$.\par\medskip
\noindent\textbf{quiz\_pagine\_191\_200\_007}: D. $f'(x)=11x^{10}+7>0$ e $f(-1)=4>0$: $f(x)>0$ su $[-1,1]$.\par\medskip
\noindent\textbf{quiz\_pagine\_191\_200\_017}: A. Con la regola della catena, $f'=\frac13(e^{2x}\sin^4x)^{-2/3}e^{2x}(2\sin^4x+4\sin^3x\cos x)$, che si riduce alla A dove definita, con estensione per continuita' negli zeri.\par\medskip
\noindent\textbf{quiz\_pagine\_201\_210\_001}: B. Per la derivata dell'inversa, $(f^{-1})'(f(2))=1/f'(2)=7$.\par\medskip
\noindent\textbf{quiz\_pagine\_201\_210\_015}: A. Il numeratore tende a $\log1=0$ e il denominatore tende a $4$, quindi il limite e' zero. NOTA: nelle spiegazioni originali si usa $x\to-2^+$ e $x\sqrt{x+2}$; la traccia stampata invece indica $x\to2^+$.\par\medskip
\noindent\textbf{quiz\_pagine\_201\_210\_020}: A. Per $h=f-g$, $h(0)=h(5)=0$. Rolle fornisce $h'(c)=0$, ossia $f'(c)=g'(c)$.\par\medskip
\noindent\textbf{quiz\_pagine\_201\_210\_025}: A. Per $x>0$, $f'(x)=\pi^x\log\pi-\pi x^{\pi-1}$; in $x=\pi$ si ottiene $\pi^\pi(\log\pi-1)$. ATTENZIONE: $x^\pi$ non e' definita come funzione reale su tutto $\mathbb R$ secondo la convenzione standard, nonostante il testo menzioni $\mathbb R$.\par\medskip
\noindent\textbf{quiz\_pagine\_211\_220\_039}: D. La derivata di $\log|u|$ e' $u'/u$; con $u=\sin2x$ si ottiene $2\cos2x/\sin2x$.\par\medskip
\noindent\textbf{quiz\_pagine\_211\_220\_040}: B. $h'(0)=f'(0)/f(0)=2$. Le altre condizioni richiederebbero ipotesi ulteriori sulla positivita' di $f$ in tutto $[0,4]$.\par\medskip
\noindent\textbf{quiz\_pagine\_221\_230\_016}: D. Si deriva $2^{u(x)}$ con $u=x\sin x$: $f'=2^u\log2\,u'$ e $u'=\sin x+x\cos x$.\par\medskip
\noindent\textbf{quiz\_pagine\_221\_230\_017}: D. Per definizione, il rapporto incrementale di $h=fg$ e' $(h(x)-h(x_0))/(x-x_0)$.\par\medskip
\noindent\textbf{quiz\_pagine\_221\_230\_018}: A. $g(x_0)=0$ e il rapporto incrementale di $g$ e' $f(x)$, che tende a $f(x_0)=0$ per continuita'.\par\medskip
\noindent\textbf{quiz\_pagine\_221\_230\_019}: C. La derivata dell'inversa in $f(2)=1$ vale $1/f'(2)=1/3$.\par\medskip
\noindent\textbf{quiz\_pagine\_221\_230\_027}: non verificata / nessuna alternativa corretta. La continuita' in $1$ impone $b=1/2+a$; la derivabilita' impone $3/2+a=1$, cioe' $a=-1/2$, $b=0$. La risposta C esprime solo la continuita', non entrambe: nessuna opzione e' sufficiente. Si segnala l'incongruenza.\par\medskip
\noindent\textbf{quiz\_pagine\_231\_240\_001}: D. $f(1)=1$, $f\prime(1)=2$, $f\prime\prime(1)=-1$, quindi $T_2(x)=1+2(x-1)-\frac12(x-1)^2$.\par\medskip
\noindent\textbf{quiz\_pagine\_231\_240\_003}: B. $f\prime=e^x(x^2+2x)$ e $f\prime\prime=e^x(x^2+4x+2)$.\par\medskip
\noindent\textbf{quiz\_pagine\_231\_240\_004}: D. $f\prime(x_0)=3e^{x_0^2+x_0}(2x_0+1)$, quindi $f\prime(x_0)x_0/f(x_0)=2x_0^2+x_0$.\par\medskip
\noindent\textbf{quiz\_pagine\_231\_240\_006}: A. $f(x)=3x^{-1/2}$ per $x>0$, dunque $f\prime(x)=-\frac32x^{-3/2}$ e $f\prime(1)=-3/2$.\par\medskip
\noindent\textbf{quiz\_pagine\_231\_240\_010}: C. La derivabilita' implica continuita': $f(x)-f(x_0)=\frac{f(x)-f(x_0)}{x-x_0}(x-x_0)\to0$.\par\medskip
\noindent\textbf{quiz\_pagine\_231\_240\_015}: C. $e^x=1+x+x^2/2+x^3/6+o(x^3)$; sottraendo $x^2$ si ottiene il coefficiente $-1/2$.\par\medskip
\noindent\textbf{quiz\_pagine\_231\_240\_017}: C. $f\prime(x)=e^x(1+x)$, percio' $xf\prime(x)=(1+x)xe^x=(1+x)f(x)$.\par\medskip
\noindent\textbf{quiz\_pagine\_231\_240\_024}: A. $f^{-1}(y)=y^2-3$ per $y\geq0$, quindi $(f^{-1})\prime(2)=2\cdot2=4$.\par\medskip
\noindent\textbf{quiz\_pagine\_241\_250\_010}: D. Per $x>0$, $f(x)=x^{4x^2}$ e $f\prime(x)=x^{4x^2}(8x\ln x+4x)=4x f(x)(2\ln x+1)$.\par\medskip
\noindent\textbf{quiz\_pagine\_241\_250\_025}: C. Con prodotto e catena: $7\sin^6x\cos^{12}x-11\sin^8x\cos^{10}x=\sin^6x\cos^{10}x(7-18\sin^2x)$.\par\medskip
\noindent\textbf{quiz\_pagine\_241\_250\_026}: B. Per $x>0$, $f=e^{\sin x\ln x}$; la derivata e' $f(\cos x\ln x+\sin x/x)$.\par\medskip
\noindent\textbf{quiz\_pagine\_241\_250\_029}: B. $C^k(I)$ significa che esistono le derivate fino all'ordine $k$ e sono continue su $I$ (in particolare la derivata $k$-esima).\par\medskip
\noindent\textbf{quiz\_pagine\_281\_290\_006}: E. E' la definizione di derivabilita' in un punto, con $l=f'(x_0)$.\par\medskip
\noindent\textbf{quiz\_pagine\_281\_290\_025}: B. Per $x\neq0$, la prima derivata e' $2/(4+x^2)$ e la seconda (con il segno meno) e' ancora $2/(4+x^2)$: somma $4/(4+x^2)$.\par\medskip
\noindent\textbf{quiz\_pagine\_291\_300\_001}: E. La derivazione della radice cubica e del prodotto da' $\frac43(e^{4x})^{1/3}(\cos x)^{5/3}(\cos x-2\sin x)$, sul dominio reale della funzione.\par\medskip
\noindent\textbf{quiz\_pagine\_291\_300\_005}: E. $h'(0)=f'(0)/(2\sqrt{f(0)})=3/4$.\par\medskip
\noindent\textbf{quiz\_pagine\_291\_300\_011}: C. La regola della catena fornisce $-\sin x/\cos x=-\tan x$.\par\medskip
\noindent\textbf{quiz\_pagine\_291\_300\_014}: B. $(f^{-1})'(f(2))=1/f'(2)=7$.\par\medskip
\noindent\textbf{quiz\_pagine\_291\_300\_019}: E. Le derivate pari di $\sinh x$ sono $\sinh x$ e si annullano in zero; quelle dispari sono $\cosh x$. Anche B puo' essere letta come vera se si ammettono indici negativi nell'equivalenza, ma le derivate hanno ordine $n\geq0$.\par\medskip
\noindent\textbf{quiz\_pagine\_291\_300\_024}: A. Derivazione logaritmica: $f'=x^{x^4}(4x^3\log x+x^3)=x^{x^4+3}(1+\log x^4)$.\par\medskip
\noindent\textbf{quiz\_pagine\_291\_300\_028}: A. Derivando $\log f=5\sin x\log(\cos x)$ si ottiene $f'=5(\cos x)^{5\sin x-1}(\cos^2x\log(\cos x)-\sin^2x)$.\par\medskip
\noindent\textbf{quiz\_pagine\_291\_300\_030}: E. $f(x)=x^{3x^2}$ e $f'(x)=x^{3x^2}(6x\log x+3x)=3x^{3x^2+1}(2\log x+1)$.\par\medskip
\noindent\textbf{quiz\_pagine\_291\_300\_031}: E. Per $x>0$, $f(x)=\sin x-x$; per $x<0$, $f(x)=-\sin x+x$. Entrambe le derivate laterali in zero sono nulle.\par\medskip
\noindent\textbf{quiz\_pagine\_291\_300\_033}: C. Regola del prodotto e $\cos^2x=1-\sin^2x$: $f'=\sin^6x\cos^{10}x(7\cos^2x-11\sin^2x)$.\par\medskip
\noindent\textbf{quiz\_pagine\_291\_300\_034}: B. $x^{\sin x}=\exp(\sin x\log x)$; la derivata e' $x^{\sin x}(\cos x\log x+\sin x/x)$.\par\medskip
\noindent\textbf{quiz\_pagine\_291\_300\_035}: D. Per il teorema di Lagrange, i rapporti incrementali vicino a 2 coincidono con valori di $f'$ in punti intermedi; tendono quindi a $5$.\par\medskip
\noindent\textbf{quiz\_pagine\_291\_300\_038}: E. $f'(0)=3+1=4$ e $f(0)=1$, dunque $(f^{-1})'(1)=1/4$.\par\medskip
\noindent\textbf{quiz\_pagine\_301\_310\_004}: A. La derivata della funzione inversa in $5=f(4)$ vale $1/f\prime(4)=8$.\par\medskip
\noindent\textbf{quiz\_pagine\_301\_310\_015}: D. Per $x\ne0$, $f\prime(x)=\frac23 x^{-1/3}$, che diverge in modulo vicino a zero. La funzione e' continua in zero ma non derivabile.\par\medskip
\noindent\textbf{quiz\_pagine\_311\_320\_003}: D. Per derivazione logaritmica $f\prime/f=5\cos x\log(3\sin x)+5\cos x$; a $\pi/2$ entrambi i termini si annullano.\par\medskip
\noindent\textbf{quiz\_pagine\_311\_320\_004}: B. Poiche $f(0)=3>0$, $f$ rimane positiva vicino a zero: $|f|=f$ e $g\prime(0)=-4$.\par\medskip
\noindent\textbf{quiz\_pagine\_311\_320\_005}: D. $f\prime(0)=2+3=5$, dunque $(f^{-1})\prime(2)=1/5$.\par\medskip
\noindent\textbf{quiz\_pagine\_311\_320\_006}: D. $e^{x^4}=1+x^4+O(x^8)$, $\sqrt{1+x^4}=1+x^4/2+O(x^8)$: la differenza e $x^4/2+O(x^8)\geq0$.\par\medskip
\noindent\textbf{quiz\_pagine\_311\_320\_007}: non verificata / nessuna alternativa corretta. NESSUNA ALTERNATIVA CORRETTA, come segnala la fonte: $f(2)=2>0$, $f\prime(2)=-3<0$, $f\prime\prime(2)=-10<0$: positiva, decrescente e concava.\par\medskip
\noindent\textbf{quiz\_pagine\_311\_320\_010}: B. $\sqrt{1+3x^3}=1+3x^3/2+o(x^3)$, dunque $f(x)=e(1+3x^3/2)+o(x^3)$.\par\medskip
\noindent\textbf{quiz\_pagine\_311\_320\_011}: C. Dalla definizione $f(x)/x^3\to0$; dividendo per $x^3+x^2=x^3(1+1/x)$ si ottiene ancora zero.\par\medskip
\noindent\textbf{quiz\_pagine\_311\_320\_016}: C. $1/\cos x=1+x^2/2+o(x^2)$ e $\sin x=x+o(x^2)$, quindi il prodotto e $1+x+x^2/2+o(x^2)$.\par\medskip
\noindent\textbf{quiz\_pagine\_311\_320\_018}: D. $1/(1+3x)=1-3x+9x^2+o(x^2)$; moltiplicando per $3x^2$ e sottraendo $3x^2+3x^4$ resta $-9x^3+24x^4+o(x^4)$.\par\medskip
\noindent\textbf{quiz\_pagine\_311\_320\_019}: B. La derivata dell\textbackslash{}'inversa nel punto 7 vale $1/f\prime(2)=7$, poiche $f(2)=7$.\par\medskip
\noindent\textbf{quiz\_pagine\_311\_320\_020}: A. $\sin^2x=x^2+o(x^2)$ e $\log(5-\sin^2x)=\log5-x^2/5+o(x^2)$, quindi zero e massimo locale.\par\medskip
\noindent\textbf{quiz\_pagine\_311\_320\_021}: E. Poiche $\sinh(\log(1+x))=((1+x)-(1+x)^{-1})/2=x-x^2/2+x^3/2+o(x^3)$. Le alternative A e D riportano anche una costante iniziale errata, conservata dalla fonte.\par\medskip
\noindent\textbf{quiz\_pagine\_311\_320\_024}: C. Lo sviluppo e centrato in 7, non in zero (refuso originale): $f(7)=-3$ e $f\prime(7)=2$, quindi la tangente e $y=-3+2(x-7)$.\par\medskip
\noindent\textbf{quiz\_pagine\_311\_320\_025}: A. $\cosh x=1+x^2/2+O(x^4)$; $\log(\cosh x)=x^2/2+O(x^4)=x^2/2+o(x^3)$.\par\medskip
\noindent\textbf{quiz\_pagine\_311\_320\_027}: A. $\sinh x=x+x^3/6+o(x^3)$; $\log(1+u)=u-u^2/2+u^3/3+o(u^3)$, dunque $x-x^2/2+(1/6+1/3)x^3+o(x^3)$.\par\medskip
\noindent\textbf{quiz\_pagine\_311\_320\_028}: D. $2e^{-x}-x^2=2-2x+o(x^2)$; moltiplicando per $(1+x^2)^{-1}=1-x^2+o(x^2)$ si ottiene $2-2x-2x^2+o(x^2)$. Il polinomio e senza resto: D.\par\medskip
\noindent\textbf{quiz\_pagine\_311\_320\_029}: E. Il termine dominante rispetto al valore $f(2)=5$ e $-4(x-2)^6$: zero derivata e massimo locale, concavita negativa ai due lati. Non e un flesso.\par\medskip
\noindent\textbf{quiz\_pagine\_311\_320\_030}: C. $1/(1-x^2)=1+x^2+x^4+\cdots$; il prodotto con $1+2x^2$ da $1+3x^2+3x^4+\cdots$.\par\medskip
\noindent\textbf{quiz\_pagine\_311\_320\_032}: B. $f(1)=e^{-1/3}$, $f\prime(1)=-2e^{-1/3}/3$, $f\prime\prime(1)=-2e^{-1/3}/9$; Taylor: $e^{-1/3}[1-2(x-1)/3-(x-1)^2/9]$.\par\medskip
\noindent\textbf{quiz\_pagine\_311\_320\_033}: C. $f\prime(x)=-e^{-3x^2}(6x^2+6x+1)/[4(x+1)^2]$, positivo tra le radici $-1/2\pm1/(2\sqrt3)$, entrambe maggiori di $-1$.\par\medskip
\noindent\textbf{quiz\_pagine\_311\_320\_035}: C. $\cosh x=1+x^2/2+x^4/24+o(x^4)$, quindi $e^{\cosh x}=e[1+x^2/2+(1/24+1/8)x^4]+o(x^4)=e[1+x^2/2+x^4/6]+o(x^4)$.\par\medskip
\noindent\textbf{quiz\_pagine\_321\_330\_005}: C. La derivata di $a^{u(x)}$ \textbackslash{}`e $a^u\log(a)u\prime$.\par\medskip
\noindent\textbf{quiz\_pagine\_321\_330\_013}: C. $f(0)=3$ e $f\prime(0)=6$, quindi $(f^{-1})\prime(3)=1/6$.\par\medskip
\noindent\textbf{quiz\_pagine\_321\_330\_014}: C. Vicino a $x=1$ il termine $|\cos\pi x|^3$ ha derivata nulla in $1$, mentre $2|x|$ ha derivata $2$.\par\medskip
\noindent\textbf{quiz\_pagine\_321\_330\_027}: C. $f\prime(x)=(x\cos x-\sin x)/x^2$; in $\pi/2$ vale $-4/\pi^2$.\par\medskip
\noindent\textbf{quiz\_pagine\_321\_330\_037}: B. La regola del quoziente porta a $e^{2x}[(1+2x)(e^{2x}-1)-2xe^{2x}]/(e^{2x}-1)^2=e^{2x}(e^{2x}-2x-1)/(e^{2x}-1)^2$.\par\medskip
\noindent\textbf{quiz\_pagine\_331\_340\_005}: C. Il rapporto incrementale di $h(x)=[f(x)]^3$ in zero \textbackslash{}`e $[h(x)-h(0)]/x=f(x)^3/x$.\par\medskip
\noindent\textbf{quiz\_pagine\_331\_340\_015}: E. Per la regola della catena $f\prime=e^{\sin(1/(x+1))}\cos(1/(x+1))[-1/(x+1)^2]$.\par\medskip
\noindent\textbf{quiz\_pagine\_331\_340\_016}: A. $f(0)=1$, $f\prime(0)=2$, quindi $(f^{-1})\prime(1)=1/f\prime(0)=1/2$.\par\medskip
\noindent\textbf{quiz\_pagine\_341\_350\_004}: C. Per la regola della catena $f\prime=(1/2)(1+\sin^2x)^{-1/2}\cdot2\sin x\cos x$.\par\medskip
\noindent\textbf{quiz\_pagine\_341\_350\_006}: A. La formula della derivata inversa da $(f^{-1})\prime(1)=1/f\prime(0)=1/(1+3)=1/4$.\par\medskip
\noindent\textbf{quiz\_pagine\_341\_350\_011}: B. $g(1)=\sin(f(1))=\sin\pi=0$; il rapporto incrementale e $[g(x)-g(1)]/(x-1)$.\par\medskip
\noindent\textbf{quiz\_pagine\_341\_350\_022}: E. $g\prime(x)=-3f\prime(x)/f(x)^4$, quindi $g\prime(0)=-9/16$. ATTENZIONE: il risultato matematico e $-9/16$, dunque la risposta corretta e E; l\textbackslash{}textquotesingle alternativa A non e corretta.\par\medskip
\noindent\textbf{quiz\_pagine\_341\_350\_028}: E. Per $c=1$ i due rami si raccordano con valore $0$, ma le derivate laterali sono $-1$ e $1$: il punto e un minimo assoluto, non derivabile.\par\medskip
\noindent\textbf{quiz\_pagine\_341\_350\_030}: D. La regola del prodotto da $3x^2\log(\cos x)-x^3\tan x$. L\textbackslash{}textquotesingle alternativa B e equivalente sul dominio reale, dove $\cos x>0$: il quesito ha due alternative corrette.\par\medskip
\noindent\textbf{quiz\_pagine\_351\_360\_004}: C. Per la regola della catena $f\prime=(1/2)(1+\sin^2x)^{-1/2}\cdot2\sin x\cos x$.\par\medskip
\noindent\textbf{quiz\_pagine\_351\_360\_006}: A. La formula della derivata inversa da $(f^{-1})\prime(1)=1/f\prime(0)=1/(1+3)=1/4$.\par\medskip
\noindent\textbf{quiz\_pagine\_351\_360\_011}: B. $g(1)=\sin(f(1))=\sin\pi=0$; il rapporto incrementale e $[g(x)-g(1)]/(x-1)$.\par\medskip
\noindent\textbf{quiz\_pagine\_351\_360\_022}: E. $g\prime(x)=-3f\prime(x)/f(x)^4$, quindi $g\prime(0)=-9/16$. ATTENZIONE: il risultato matematico e $-9/16$, dunque la risposta corretta e E; l\textbackslash{}textquotesingle alternativa A non e corretta.\par\medskip
\noindent\textbf{quiz\_pagine\_351\_360\_028}: E. Per $c=1$ i due rami si raccordano con valore $0$, ma le derivate laterali sono $-1$ e $1$: il punto e un minimo assoluto, non derivabile.\par\medskip
\noindent\textbf{quiz\_pagine\_351\_360\_030}: D. La regola del prodotto da $3x^2\log(\cos x)-x^3\tan x$. L\textbackslash{}textquotesingle alternativa B e equivalente sul dominio reale, dove $\cos x>0$: il quesito ha due alternative corrette.\par\medskip
\noindent\textbf{quiz\_pagine\_361\_370\_016}: D. $|x\sin(1/x)|\leq|x|$ prova la continuita in zero; il rapporto incrementale $\sin(1/x)$ non ha limite, quindi non e derivabile in zero.\par\medskip
\noindent\textbf{quiz\_pagine\_361\_370\_022}: D. Continuita in $1$: $2\alpha-\beta=1$. Uguaglianza delle derivate: $2\alpha=4\beta-\alpha$, ossia $3\alpha=4\beta$. Risolvendo $\alpha=4/5$, $\beta=3/5$.\par\medskip
\noindent\textbf{quiz\_pagine\_371\_380\_016}: E. $g\prime=-2f\prime/f^3$ e $g\prime(0)=-6/(-8)=3/4$.\par\medskip
\noindent\textbf{quiz\_pagine\_371\_380\_023}: D. In $0$ i valori laterali coincidono e valgono $1$; le derivate laterali sono entrambe nulle, perche $e^0-1=0$ e $-\sin0=0$.\par\medskip
\noindent\textbf{quiz\_pagine\_371\_380\_025}: E. $f\prime=(6x+6x^2)/(3x^2+2x^3)=6(1+x)/(x(3+2x))$ sul dominio del logaritmo.\par\medskip
\noindent\textbf{quiz\_pagine\_381\_390\_001}: C. Se $-1<f\prime(0)<0$, per il teorema della derivata dell\textbackslash{}textquotesingle inversa $m=1/f\prime(0)<-1$. Attenzione: la sola stretta decrescenza non esclude $f\prime(0)=0$, nel qual caso la derivata dell\textbackslash{}textquotesingle inversa non e finita: manca una ipotesi nel testo.\par\medskip
\noindent\textbf{quiz\_pagine\_381\_390\_013}: E. La regola della catena da $(g\circ f)\prime(-2)=g\prime(f(-2))f\prime(-2)=(-7)(-4)=28$. Il testo originale omette l\textbackslash{}textquotesingle argomento della derivata, che va inteso $-2$.\par\medskip
\noindent\textbf{quiz\_pagine\_381\_390\_015}: A. Per ogni base $a>0$, $(a^x)\prime=a^x\log a$.\par\medskip
\noindent\textbf{quiz\_pagine\_381\_390\_016}: C. Poiche $f$ e derivabile tre volte, $f\prime\prime$ e derivabile una volta e quindi continua. La continuita di $f\prime\prime\prime$ non e garantita.\par\medskip
\noindent\textbf{quiz\_pagine\_381\_390\_020}: B. Per $x\neq0$, la prima derivata e $2/(4+x^2)$; la seconda e $-2/(x^2+4)$, e sottraendola si ottiene $4/(4+x^2)$.\par\medskip
\noindent\textbf{quiz\_pagine\_391\_400\_006}: D. In $0$ la continuita richiede $a=0$; le derivate laterali valgono entrambe $1$, quindi la funzione e derivabile e continua ovunque.\par\medskip
\noindent\textbf{quiz\_pagine\_391\_400\_009}: E. Per la regola della catena $f\prime(x)=(-\cos x)/(-\sin x)=\cos x/\sin x=1/\tan x$. ATTENZIONE: nel PDF le opzioni a), d), e) sembrano equivalenti (a contiene una notazione tipografica anomala); la risposta ufficiale e E.\par\medskip
\noindent\textbf{quiz\_pagine\_391\_400\_015}: E. Derivando $e^{4x/3}|\cos x|^{8/3}$ si ottiene $\frac43\sqrt[3]{e^{4x}\cos^5x}(\cos x-2\sin x)$, anche negli zeri del coseno per continuita della derivata.\par\medskip
\noindent\textbf{quiz\_pagine\_401\_410\_007}: C. Il teorema di Rolle garantisce almeno un punto $c\in(0,3)$ con $f\prime(c)=0$.\par\medskip
\noindent\textbf{quiz\_pagine\_401\_410\_009}: E. Per la regola della catena $f\prime(x)=-5e^{-5x}/(3-e^{-5x})^2=-5e^{5x}/(3e^{5x}-1)^2$.\par\medskip
\noindent\textbf{quiz\_pagine\_401\_410\_015}: A. Per il teorema di prolungamento della derivata, la continuit\textbackslash{}`a in $2$ e il limite finito di $f\prime(x)$ implicano $f\prime(2)=4$.\par\medskip
\noindent\textbf{quiz\_pagine\_411\_420\_006}: E. Per la regola del quoziente $f\prime(x)=\{x[-e^{1/x}/x^2]-(e^{1/x}-1)\}/x^2=[x-(x+1)e^{1/x}]/x^3$.\par\medskip
\noindent\textbf{quiz\_pagine\_411\_420\_012}: E. Le derivate di ordine dispari di $\cos x$ sono proporzionali a $\sin x$, quindi si annullano in $0$. L\textbackslash{}'opzione A descrive anch\textbackslash{}'essa tutti gli ordini naturali dispari, pur usando $k\in\mathbb Z$: ambiguit\textbackslash{}`a nella fonte.\par\medskip
\noindent\textbf{quiz\_pagine\_411\_420\_014}: A. Per $a=1$ i limiti coincidono in $0$ e le derivate laterali valgono entrambe $0$; altrove i rami sono regolari.\par\medskip
\noindent\textbf{quiz\_pagine\_411\_420\_021}: C. $m=(f^{-1})\prime(f(0))=1/f\prime(0)$; essendo $f\prime(0)<-1$, risulta $-1<m<0$. Anche B ed E sono vere, quindi il quesito ha pi\textbackslash{}`u risposte valide.\par\medskip
\noindent\textbf{quiz\_pagine\_411\_420\_032}: D. Poich\textbackslash{}'e $|x^3|\sin x=O(|x|^4)$, il rapporto incrementale in $0$ tende a $0$: la derivata esiste ed \textbackslash{}`e nulla.\par\medskip
\noindent\textbf{quiz\_pagine\_421\_430\_010}: E. La derivata \textbackslash{}`e $3-2\tan(2x)$; l\textbackslash{}'opzione E manca del fattore $2$.\par\medskip
\noindent\textbf{quiz\_pagine\_421\_430\_020}: A. La regola del prodotto fornisce $h\prime(0)=f\prime(0)g(0)+f(0)g\prime(0)=g(0)-3g\prime(0)$.\par\medskip
\noindent\textbf{quiz\_pagine\_421\_430\_022}: C. Per $a=2$ \textbackslash{}`e un polinomio; per $a>2$ la derivata seconda tende a zero nel punto $4$. Per $1<a<2$ la seconda derivata non esiste in $4$.\par\medskip
\noindent\textbf{quiz\_pagine\_431\_440\_004}: D. $f(0)=-1$ e $f\prime(0)=2$; la tangente \textbackslash{}`e $y=2x-1$.\par\medskip
\noindent\textbf{quiz\_pagine\_431\_440\_012}: E. Derivando $g\prime=-\sin x\,f\prime(\cos x)$ con prodotto e catena si ottiene E.\par\medskip
\noindent\textbf{quiz\_pagine\_431\_440\_022}: B. Nel dominio, $f(x)=4x+\log(\cos x)$; quindi $f\prime(x)=4-\tan x$.\par\medskip
\noindent\textbf{quiz\_pagine\_441\_450\_005}: D. $f\prime(0)=1-6=-5$; la derivata dell\textbackslash{}'inversa nel valore $-2$ \textbackslash{}`e $1/f\prime(0)=-1/5$.\par\medskip
\noindent\textbf{quiz\_pagine\_441\_450\_011}: A. Regola del prodotto: $f\prime=\frac{\sin\log(x+1)}{2\sqrt{x+1}}+\frac{\cos\log(x+1)}{\sqrt{x+1}}$.\par\medskip
\noindent\textbf{quiz\_pagine\_441\_450\_017}: A. L\textbackslash{}'inversa \textbackslash{}`e continua e crescente sulla sua immagine ma pu\textbackslash{}`o non essere derivabile (es. $f(x)=x^3$). ATTENZIONE: E \textbackslash{}`e anch\textbackslash{}'essa falsa nella forma $\le0$? Pu\textbackslash{}`o valere $f\prime=0$ (es. $x^3$), quindi \textbackslash{}`e vera.\par\medskip
\noindent\textbf{quiz\_pagine\_441\_450\_028}: B. $|x|\sin x=O(x^2)$ e $f(0)=1$; $(f(x)-1)/x\to1$, dunque \textbackslash{}`e derivabile.\par\medskip
\noindent\textbf{quiz\_pagine\_451\_460\_002}: B. Derivando $\log|f(x)|=\log|x|+\arctan(1/x)$ si ottiene $f\prime/f=1/x-1/(x^2+1)=(x^2-x+1)/(x(x^2+1))$.\par\medskip
\noindent\textbf{quiz\_pagine\_451\_460\_004}: C. Per la regola della catena $g\prime(x)=e^{f(x)}f\prime(x)$; quindi $g\prime(1)=-2e^3$.\par\medskip
\noindent\textbf{quiz\_pagine\_451\_460\_005}: D. La derivabilit\textbackslash{}`a in $x_1$ garantisce una retta tangente, orizzontale se $f\prime(x_1)=0$. Un flesso pu\textbackslash{}`o avere derivata non nulla; una derivata nulla non implica estremo; $f\prime\to0$ non basta per un asintoto orizzontale.\par\medskip
\noindent\textbf{quiz\_pagine\_451\_460\_011}: A. $f(0)=4$ e $f(x)-4=e^{2x}+2e^x-3=(e^x-1)(e^x+3)$; dividendo per $x$ si ha A. Il valore $4$ \textbackslash{}`e invece il limite del rapporto incrementale.\par\medskip
\noindent\textbf{quiz\_pagine\_451\_460\_015}: C. $f\prime(x)=e^{x^3-ax^2}(1+3x^3-2ax^2)$. Per $x\to-\infty$ il polinomio \textbackslash{}`e negativo qualunque sia $a$, quindi la funzione non \textbackslash{}`e crescente su tutta $\mathbb R$.\par\medskip
\noindent\textbf{quiz\_pagine\_451\_460\_020}: B. Vicino a $0$ la funzione \textbackslash{}`e pari, quindi, essendo derivabile, ha derivata nulla nell\textbackslash{}textquotesingle origine.\par\medskip
\noindent\textbf{quiz\_pagine\_451\_460\_021}: A. Le derivate laterali in $0$ valgono entrambe $2$; inoltre la derivata del termine $|x|^{3/2}$ tende a zero. La derivata \textbackslash{}`e continua in $0$, dunque $f\in C^1$.\par\medskip
\noindent\textbf{quiz\_pagine\_451\_460\_022}: A. $f(0)=1$, $f\prime(0)=1$ e $f$ \textbackslash{}`e strettamente crescente; quindi $(f^{-1})\prime(1)=1/f\prime(0)=1$, valore non elencato.\par\medskip
\noindent\textbf{quiz\_pagine\_451\_460\_026}: A. $g\prime(x)=\cos(f(x))f\prime(x)$, perci\textbackslash{}`o $g\prime(0)=g\prime(1)=0$. Inoltre $g(0)=g(1)$: per Rolle esiste almeno un punto interno in cui $g\prime=0$.\par\medskip
\noindent\textbf{quiz\_pagine\_451\_460\_032}: E. La disuguaglianza sui valori non impone segni alle derivate: ad esempio $f(x)=e^{2x}+C+\sin(kx)$ con $C\ge1$ pu\textbackslash{}`o avere derivata negativa e seconda derivata negativa in alcuni punti, scegliendo $k$ grande.\par\medskip
\noindent\textbf{quiz\_pagine\_461\_470\_009}: C. Per la regola del quoziente, $g\prime(x)=[xf\prime(x)-f(x)]/x^2$ e quindi $g\prime(1)=2-f(1)$.\par\medskip
\noindent\textbf{quiz\_pagine\_461\_470\_020}: D. Le derivate di $\sin x$ sono periodiche di periodo quattro: per $k=4n$, $f^{(k)}(\pi/2)=\sin(\pi/2)=1$. Anche B \textbackslash{}`e formalmente vera se si ammettono ordini negativi, ma tali derivate non sono definite: D \textbackslash{}`e la scelta pertinente.\par\medskip
\noindent\textbf{quiz\_pagine\_461\_470\_023}: E. La regola della catena applicata al logaritmo d\textbackslash{}`a $f\prime=[3h^2h\prime]/(h^3+1)$, nei punti del dominio.\par\medskip
\noindent\textbf{quiz\_pagine\_471\_480\_012}: D. $f'(x)=3+1/(1+x^2)$, dunque $f'(0)=4$ e $(f^{-1})'(1)=1/f'(0)=1/4$.\par\medskip
\noindent\textbf{quiz\_pagine\_471\_480\_016}: D. Per la regola della catena $f'(x)=e^{\sin(1/(x+1))}\cos(1/(x+1))[-1/(x+1)^2]$.\par\medskip
\noindent\textbf{quiz\_pagine\_471\_480\_026}: A. $f'(x)=(4x+9x^2)/(2x^2+3x^3)=(4+9x)/(x(2+3x))$ nel dominio.\par\medskip
\noindent\textbf{quiz\_pagine\_481\_490\_018}: E. La funzione deve essere continua su $[a,7]$: occorre $a>5$. Inoltre $f(7)=\log2>0$ e $f(a)<0$ richiede $|a-5|<1$, cio\textbackslash{}`e $a<6$. Il caso $a=6$ ha gi\textbackslash{}`a uno zero all\textbackslash{}'estremo, ma non assicura uno zero interno tramite il criterio dei segni opposti.\par\medskip
\noindent\textbf{quiz\_pagine\_501\_510\_008}: D. La derivata di $\log|u|$ \textbackslash{}`e $u\prime/u$: con $u=\sin2x$ si ottiene $2\cos2x/\sin2x$.\par\medskip
\noindent\textbf{quiz\_pagine\_501\_510\_028}: B. La derivazione logaritmica d\textbackslash{}`a $f\prime/f=-4\sin x\log(3\cos x)-4\sin x$; in $x=0$ entrambi i termini sono nulli.\par\medskip
\noindent\textbf{quiz\_pagine\_501\_510\_029}: D. Entrambi i rami tendono a zero; la derivata destra del ramo $2\sqrt x$ tende a $+\infty$, mentre quella sinistra vale $2$.\par\medskip
\noindent\textbf{quiz\_pagine\_501\_510\_030}: A. $h\prime=-f\prime/f^2$, quindi $h\prime(0)=-3/16$.\par\medskip
\noindent\textbf{quiz\_pagine\_511\_520\_008}: A. Per $x>0$, $f=\frac12\log(1+x^2)-\log x-1/(2x^2)$. Quindi $f'=x/(1+x^2)-1/x+1/x^3=1/[x^3(1+x^2)]$.\par\medskip
\noindent\textbf{quiz\_pagine\_511\_520\_012}: E. $f(x)\sim2\sqrt{x}$ e $f'(x)=\sqrt{x}+(x+2)/(2\sqrt{x})\to+\infty$ per $x\to0^+$: tangente verticale.\par\medskip
\noindent\textbf{quiz\_pagine\_511\_520\_029}: D. Derivando $a^{u(x)}$ si ottiene $a^{u(x)}\log a\,u'(x)$, con $u'=\sin x+x\cos x$.\par\medskip
\noindent\textbf{quiz\_pagine\_511\_520\_030}: A. Il rapporto incrementale di $g$ in $x_0$ \textbackslash{}`e $f(x)$, che tende a $f(x_0)=0$ perch'e un punto angoloso implica continuit\textbackslash{}`a. Quindi $g'(x_0)=0$.\par\medskip
\noindent\textbf{quiz\_pagine\_511\_520\_031}: C. La formula per la derivata dell'inversa d\textbackslash{}`a $(f^{-1})'(f(2))=1/f'(2)=1/3$.\par\medskip
\subsubsection*{4.2 -- I teoremi fondamentali del calcolo differenziale}
\noindent\textbf{quiz\_pagine\_001\_010\_038}: A. Spiegazione non ancora disponibile.\par\medskip
\noindent\textbf{quiz\_pagine\_011\_020\_012}: non verificata / nessuna alternativa corretta. $f(2)=2>0$, $f\prime(2)=-3<0$, $f\prime\prime(2)=-10<0$: positiva, decrescente e concava. Nessuna opzione coincide.\par\medskip
\noindent\textbf{quiz\_pagine\_011\_020\_015}: non verificata / nessuna alternativa corretta. La funzione \textbackslash{}`e $e+O(x^4)$, quindi il polinomio di ordine 3 \textbackslash{}`e costante $e$. Verificare la formula stampata nel PDF: l\textbackslash{}'estrazione non chiarisce l\textbackslash{}'esponente.\par\medskip
\noindent\textbf{quiz\_pagine\_011\_020\_021}: C. $1/\cos x=1+x^2/2+o(x^2)$ e $\sin x=x+o(x^2)$.\par\medskip
\noindent\textbf{quiz\_pagine\_011\_020\_023}: D. $3x^2(1-3x+9x^2)-3x^2-3x^4=-9x^3+24x^4+o(x^4)$.\par\medskip
\noindent\textbf{quiz\_pagine\_011\_020\_025}: A. $f(x)=2+\log5-x^2/5+o(x^2)$: massimo locale stretto.\par\medskip
\noindent\textbf{quiz\_pagine\_011\_020\_026}: E. $\sinh(\log(1+x))=((1+x)-(1+x)^{-1})/2=x-x^2/2+x^3/2+o(x^3)$.\par\medskip
\noindent\textbf{quiz\_pagine\_011\_020\_029}: C. $f(7)=-3$ e $f\prime(7)=2$, quindi la tangente \textbackslash{}`e $y=-3+2(x-7)$.\par\medskip
\noindent\textbf{quiz\_pagine\_011\_020\_030}: A. $\cosh x=1+x^2/2+O(x^4)$, dunque $\log\cosh x=x^2/2+O(x^4)$.\par\medskip
\noindent\textbf{quiz\_pagine\_011\_020\_032}: A. $\sinh x=x+x^3/6+o(x^3)$; in $\log(1+u)=u-u^2/2+u^3/3+\cdots$ il coefficiente cubico \textbackslash{}`e $1/6+1/3=1/2$.\par\medskip
\noindent\textbf{quiz\_pagine\_011\_020\_033}: D. $2e^{-x}-x^2=2-2x+O(x^3)$ e $(1+x^2)^{-1}=1-x^2+o(x^2)$.\par\medskip
\noindent\textbf{quiz\_pagine\_011\_020\_034}: E. Il termine dominante $-4(x-2)^6$ d\textbackslash{}`a un massimo piatto in $2$, non un flesso crescente.\par\medskip
\noindent\textbf{quiz\_pagine\_011\_020\_035}: C. $(1+2x^2)(1+x^2+x^4+\cdots)=1+3x^2+3x^4+\cdots$.\par\medskip
\noindent\textbf{quiz\_pagine\_021\_030\_001}: B. $f(1)=e^{-1/3}$, $f\prime(1)=-2e^{-1/3}/3$, $f\prime\prime(1)=-2e^{-1/3}/9$, dunque il coefficiente quadratico \textbackslash{}`e $-e^{-1/3}/9$.\par\medskip
\noindent\textbf{quiz\_pagine\_021\_030\_004}: C. $\cosh x=1+x^2/2+x^4/24+o(x^4)$; $e^{\cosh x}=e[1+x^2/2+(1/24+1/8)x^4]+o(x^4)$.\par\medskip
\noindent\textbf{quiz\_pagine\_021\_030\_012}: D. $3^{2\sin x}=e^{2\log3(x+o(x^2))}=1+2x\log3+2x^2\log^23+o(x^2)$.\par\medskip
\noindent\textbf{quiz\_pagine\_021\_030\_028}: B. $(1+x^4)^{1/3}=1+x^4/3+o(x^4)$, $(1+x^3)^{-1/4}=1-x^3/4+o(x^4)$.\par\medskip
\noindent\textbf{quiz\_pagine\_041\_050\_002}: E. Risposta ufficiale indicata nella tabella della decima simulazione.\par\medskip
\noindent\textbf{quiz\_pagine\_041\_050\_003}: D. Risposta ufficiale indicata nella tabella della decima simulazione.\par\medskip
\noindent\textbf{quiz\_pagine\_041\_050\_006}: D. Per il teorema di Rolle, poiché f(-3)=f(-2)=0, esiste c interno con derivata nulla.\par\medskip
\noindent\textbf{quiz\_pagine\_041\_050\_017}: E. Risposta ufficiale indicata nella tabella della decima simulazione.\par\medskip
\noindent\textbf{quiz\_pagine\_051\_060\_005}: E. Il teorema di Rolle richiede continuità sul chiuso, derivabilità nell'aperto e uguaglianza dei valori agli estremi.\par\medskip
\noindent\textbf{quiz\_pagine\_051\_060\_006}: B. La risposta ufficiale indica B; le ipotesi riportate non specificano una forma indeterminata per applicare de l'Hôpital.\par\medskip
\noindent\textbf{quiz\_pagine\_051\_060\_012}: E. La risposta ufficiale è E, ma A e B risultano soddisfatte: la funzione è continua sul chiuso e derivabile nell'aperto.\par\medskip
\noindent\textbf{quiz\_pagine\_061\_070\_023}: D. Per Rolle, f(-2)=f(0)=0 implica un punto c in (-2,0) con derivata nulla.\par\medskip
\noindent\textbf{quiz\_pagine\_071\_080\_003}: E. Nello sviluppo in $x=3$ il coefficiente di $(x-3)^2$ è $3$, dunque $f''(3)=6$. La tabella ufficiale indica E per la derivata terza; la sua esistenza richiede cautela con il solo resto di Peano.\par\medskip
\noindent\textbf{quiz\_pagine\_071\_080\_014}: C. Lo sviluppo centrato in $-4$ dà $f'(-4)=0$ e $f''(-4)=6>0$: $-4$ è un minimo relativo.\par\medskip
\noindent\textbf{quiz\_pagine\_081\_090\_003}: B. $\cos x=1-x^2/2+O(x^4)$ e $\cosh x=1+x^2/2+O(x^4)$: il termine quadratico si annulla, quindi $f(x)=1+o(x^2)$. La risposta C è uno sviluppo corretto di ordine superiore.\par\medskip
\noindent\textbf{quiz\_pagine\_081\_090\_019}: E. $e^x-1=x+O(x^2)$, quindi $(e^x-1)^2=x^2+O(x^3)$ e $\sin x=x+O(x^3)$; il polinomio è $x+x^2$. L\textbackslash{}textquotesingle alternativa A rappresenta invece lo sviluppo con resto.\par\medskip
\noindent\textbf{quiz\_pagine\_091\_100\_002}: B. Il termine principale è $2(x+3)^3$, che cambia segno attraversando $-3$. Non è lecito dedurre $f(0)$ dallo sviluppo locale; $f^{\prime\prime\prime}(-3)=12$ se la derivata esiste.\par\medskip
\noindent\textbf{quiz\_pagine\_091\_100\_012}: D. Dallo sviluppo $f(\pi)=\pi$ e $f\prime(\pi)=2$; la tangente è $y=\pi+2(x-\pi)$.\par\medskip
\noindent\textbf{quiz\_pagine\_091\_100\_015}: C. Per il teorema di Rolle esiste almeno un $c\in(0,3)$ con $f\prime(c)=0$; il numero esatto di zeri non è determinato.\par\medskip
\noindent\textbf{quiz\_pagine\_091\_100\_019}: non verificata / nessuna alternativa corretta. Poiché $\cos x=1-x^2/2+O(x^4)$ e $\cosh x=1+x^2/2+O(x^4)$, il prodotto è $1+o(x^2)$. Anche B è formalmente vera: $x^4=o(x^2)$, quindi la scelta del coefficiente di $x^4$ non cambia il resto $o(x^2)$. La tabella ufficiale indica E.\par\medskip
\noindent\textbf{quiz\_pagine\_101\_110\_010}: C. Da $\sin x=x-x^3/6+o(x^3)$ e $x^3\cos x=x^3+o(x^3)$ si ottiene $1+x-7x^3/6+o(x^3)$.\par\medskip
\noindent\textbf{quiz\_pagine\_101\_110\_020}: E. Ponendo $u=\sin x/3=x/3-x^3/18+o(x^3)$ e usando $\log(1+u)=u-u^2/2+u^3/3+o(u^3)$, si trova $x/3-x^2/18-7x^3/162+o(x^3)$.\par\medskip
\noindent\textbf{quiz\_pagine\_111\_120\_008}: E. Moltiplicando $\cos x=1-x^2/2+O(x^4)$ e $\cosh x=1+x^2/2+O(x^4)$, i termini quadratici si cancellano: il polinomio di grado al piu' due e' $1$.\par\medskip
\noindent\textbf{quiz\_pagine\_111\_120\_016}: A. $4^x=e^{x\log4}$ e $\log4=2\log2$; applicando lo sviluppo esponenziale si ottengono i coefficienti dell'opzione a).\par\medskip
\noindent\textbf{quiz\_pagine\_111\_120\_033}: D. $2e^{-x}-x^2=2-2x+o(x^2)$ e $(1+x^2)^{-1}=1-x^2+o(x^2)$; il polinomio e' $2-2x-2x^2$.\par\medskip
\noindent\textbf{quiz\_pagine\_111\_120\_034}: C. Per il teorema di Lagrange, $f'(c)=(f(5)-f(1))/(5-1)=1$.\par\medskip
\noindent\textbf{quiz\_pagine\_111\_120\_036}: A. $(f-g)'=0$ su un intervallo; per Lagrange $f-g$ e' costante, non necessariamente nulla.\par\medskip
\noindent\textbf{quiz\_pagine\_121\_130\_011}: A. La risposta ufficiale e\textbackslash{}'e A. Attenzione: da uno sviluppo $f(x)=-1+2x+3x^2+o(x^2)$, senza ulteriori ipotesi sulla regolarita\textbackslash{}'a del resto, non segue necessariamente $f'(x)=2+6x+o(x)$. Si deduce comunque $f'(0)=2$, se lo sviluppo e\textbackslash{}'e un Taylor valido.\par\medskip
\noindent\textbf{quiz\_pagine\_121\_130\_014}: D. Poiche\textbackslash{}'e $\sinh t=t+t^3/6+o(t^3)$ e $t=x-x^3/6$, i termini cubici si cancellano: $f(x)=x+o(x^3)$.\par\medskip
\noindent\textbf{quiz\_pagine\_131\_140\_009}: E. $f(-5)=-10$, $f\prime(-5)=-2$; la tangente e\textbackslash{}'e $y=-10-2(x+5)=-2x-20$. Non vi e\textbackslash{}'e estremo locale.\par\medskip
\noindent\textbf{quiz\_pagine\_141\_150\_006}: B. $\cos(x^3)=1-x^6/2+o(x^6)$, percio\textbackslash{}'e $e^{\cos(x^3)}=e-(e/2)x^6+o(x^6)$.\par\medskip
\noindent\textbf{quiz\_pagine\_141\_150\_009}: A. $g\prime(x)=\cos(f(x))f\prime(x)$ si annulla in $0$ e $1$; Rolle, applicato a $g$ su $[0,1]$, fornisce almeno un ulteriore zero interno.\par\medskip
\noindent\textbf{quiz\_pagine\_141\_150\_025}: D. Per Lagrange $g\prime(c)=(g(2)-g(0))/2=(e^{-3}-1)/2$.\par\medskip
\noindent\textbf{quiz\_pagine\_151\_160\_017}: E. $g(11)=4$, $g(0)=0$; Lagrange da\textbackslash{}'a $g\prime(c)=(4-0)/11$.\par\medskip
\noindent\textbf{quiz\_pagine\_151\_160\_022}: E. Per $x\to0$, $f(x)/x=(ax+3x^2)\sin x=ax^2+3x^3+o(x^3)$; scegliendo $a=3$ si ottiene $h$.\par\medskip
\noindent\textbf{quiz\_pagine\_151\_160\_025}: B. $f(x)=x^4(1+o(x))$ e $f(0)=0$: per $x$ sufficientemente vicino a $0$, $f(x)\geq0$. Nulla garantisce un minimo globale.\par\medskip
\noindent\textbf{quiz\_pagine\_161\_170\_005}: D. $f'(x)=5x^4+1>0$, dunque $f([-2,2])=[-44,24]$: $32$ non e\textbackslash{}'e assunto. Le altre condizioni sono soddisfatte.\par\medskip
\noindent\textbf{quiz\_pagine\_161\_170\_008}: B. La funzione si annulla in $0,1,2,3$; per Rolle $f'$ ha almeno uno zero in ciascuno degli intervalli $(0,1),(1,2),(2,3)$.\par\medskip
\noindent\textbf{quiz\_pagine\_161\_170\_009}: A. $g(1)=\sin(2\pi)=0=g(2)$; Rolle garantisce un $c$ interno con $g'(c)=0$. Anche (c) e\textbackslash{}'e vera: per continuita\textbackslash{}'e $f$ assume $3\pi/2$, ma non necessariamente un multiplo intero di $\pi$ all'interno; quindi non garantisce $g(c)=0$.\par\medskip
\noindent\textbf{quiz\_pagine\_161\_170\_019}: A. $f(0)=f(\pi)=f(2\pi)=0$; Rolle sui due intervalli $[0,\pi]$ e $[\pi,2\pi]$ assicura due punti critici distinti.\par\medskip
\noindent\textbf{quiz\_pagine\_171\_180\_001}: E. $g(0)=f(0)$, $g'(0)=0$, $g''(0)=2f'(0)$, quindi $T_2g=f(0)+f'(0)x^2$.\par\medskip
\noindent\textbf{quiz\_pagine\_171\_180\_004}: D. $\sin(x^6)=x^6+o(x^6)$ e $\log(2+u)=\log2+u/2+o(u)$.\par\medskip
\noindent\textbf{quiz\_pagine\_171\_180\_013}: E. $f(0)=f(\pi)=0$ e $f$ e' derivabile, quindi Rolle assicura almeno un punto con $f'(c)=0$ in $(0,\pi)$.\par\medskip
\noindent\textbf{quiz\_pagine\_171\_180\_030}: D. Rolle applicato a $h=f-g$ garantisce un punto $c$ con $f'(c)=g'(c)$, non necessariamente entrambi nulli. L'ultima opzione usa il simbolo di non esistenza nell'originale.\par\medskip
\noindent\textbf{quiz\_pagine\_171\_180\_031}: D. Rolle vale su $[0,2]$ e $[2,5]$. Su $(0,1)$, $f$ e' negativa, quindi l'integrale in (e) e' negativo: anche (e) e' falsa. Quesito con piu' risposte non vere; non si puo' scegliere una soluzione univoca.\par\medskip
\noindent\textbf{quiz\_pagine\_181\_190\_021}: A. $\cos(\sqrt2|x|)=1-x^2+O(x^4)$, $\log(1+\sqrt3x)=\sqrt3x-\frac32x^2+\sqrt3x^3+O(x^4)$; i termini cubici si cancellano.\par\medskip
\noindent\textbf{quiz\_pagine\_191\_200\_010}: C. Il dominio e' $(0,+\infty)$: l'affermazione C pretende un segno anche su $(-\infty,0]$, dove la funzione non e' definita.\par\medskip
\noindent\textbf{quiz\_pagine\_191\_200\_016}: C. Integrando due volte $f''=0$ si ottiene $f'=a$ e $f=ax+b$, con $a,b$ arbitrari.\par\medskip
\noindent\textbf{quiz\_pagine\_211\_220\_041}: E. Per Lagrange $g'(c)=[g(3)-g(0)]/3=(e^{-2}-1)/3$.\par\medskip
\noindent\textbf{quiz\_pagine\_221\_230\_022}: D. Per il teorema di Lagrange esiste $c\in(-1,6)$ con $f'(c)=(3-0)/(6-(-1))=3/7$.\par\medskip
\noindent\textbf{quiz\_pagine\_221\_230\_023}: A. $g'=-\sin(f)f'$ e si annulla agli estremi $-5,9$; inoltre Rolle fornisce un punto interno dove $f'=0$.\par\medskip
\noindent\textbf{quiz\_pagine\_221\_230\_026}: D. Il polinomio e' regolare; $f(a)=f(b)$ quando $a,b$ sono simmetrici rispetto al vertice $x=2$: $-1$ e $5$.\par\medskip
\noindent\textbf{quiz\_pagine\_221\_230\_030}: A. $f'(x)=e^{-x}(x-1)>0$ su $[2,4]$, quindi il minimo e' in $2$ e il massimo in $4$.\par\medskip
\noindent\textbf{quiz\_pagine\_231\_240\_019}: B. $[f(0)-f(-3)]/3=-18/3=-6$; $f\prime(c)=2c-3=-6$ fornisce $c=-3/2$.\par\medskip
\noindent\textbf{quiz\_pagine\_241\_250\_014}: C. $\sin x=x-x^3/6+o(x^3)$ e $e^{\sin x}=1+x+x^2/2+0\cdot x^3+o(x^3)$.\par\medskip
\noindent\textbf{quiz\_pagine\_241\_250\_017}: E. Uno sviluppo locale attorno a 3 non determina $f(0)$, esterno all'intorno. Le altre proprieta' derivano dai coefficienti di Taylor.\par\medskip
\noindent\textbf{quiz\_pagine\_241\_250\_023}: D. $2e^{-x}-x^2=2-2x+0x^2-x^3/3+o(x^3)$; moltiplico per $1-x^2+o(x^3)$ e ottengo $2-2x-2x^2+5x^3/3+o(x^3)$.\par\medskip
\noindent\textbf{quiz\_pagine\_251\_260\_001}: E. Per $x>0$, $f(x)=(2x)^{4x^2}$ e $f'(x)=4x(2x)^{4x^2}(2\ln(2x)+1)$. Le prime quattro scritture sono equivalenti, la quinta no.\par\medskip
\noindent\textbf{quiz\_pagine\_251\_260\_003}: B. $f(x)=-2x$ per $x<-4$, $f(x)=8$ per $-4\leq x\leq4$, $f(x)=2x$ per $x>4$. I punti di non derivabilita' sono due, $-4$ e $4$. La notazione $f(-4,0)$ nell'opzione C e' riportata come nell'originale ed e' ambigua.\par\medskip
\noindent\textbf{quiz\_pagine\_251\_260\_004}: A. $\cosh x=1+x^2/2+O(x^4)$; pertanto $\log(\cosh x)=x^2/2+O(x^4)=x^2/2+o(x^3)$.\par\medskip
\noindent\textbf{quiz\_pagine\_251\_260\_006}: A. Usando $\sinh x=x+x^3/6+o(x^3)$ e $\log(1+t)=t-t^2/2+t^3/3+o(t^3)$ si ottiene $x-x^2/2+x^3/2+o(x^3)$.\par\medskip
\noindent\textbf{quiz\_pagine\_251\_260\_007}: D. $2e^{-x}-x^2=2-2x+o(x^2)$ e $(1+x^2)^{-1}=1-x^2+o(x^2)$, dunque il polinomio di grado due e' $2-2x-2x^2$; l'opzione A comprende anche il resto.\par\medskip
\noindent\textbf{quiz\_pagine\_251\_260\_008}: E. Il primo termine non costante e' $-4(x-2)^6$: $x=2$ e' un massimo locale con derivata nulla, non un flesso crescente. Inoltre $f^{(7)}(2)=3\cdot7!$.\par\medskip
\noindent\textbf{quiz\_pagine\_251\_260\_011}: C. Poiche' $f(0)<0$, per continuita' $g=-f$ vicino a zero, quindi $g'(0)=-f'(0)=4$.\par\medskip
\noindent\textbf{quiz\_pagine\_251\_260\_012}: D. $g(3)=(-4)^3=-64$; il rapporto incrementale e' $[g(x)-g(3)]/(x-3)=[f^3(x)+4^3]/(x-3)$.\par\medskip
\noindent\textbf{quiz\_pagine\_251\_260\_014}: B. Per $x\to0^+$, $\sqrt[5]{x+x^2}-\sqrt[5]x=\frac15x^{6/5}+o(x^{6/5})$; domina $x^{5/6}$, che ha esponente minore.\par\medskip
\noindent\textbf{quiz\_pagine\_251\_260\_015}: E. Per $x\to+\infty$, i primi due termini sono di ordini $x^{2/5}$ e $x^{1/5}$, entrambi trascurabili rispetto a $x^{5/6}$. La parte principale e' $x^{5/6}$, non presente nelle alternative.\par\medskip
\noindent\textbf{quiz\_pagine\_251\_260\_018}: E. Il limite finito di $f'(x)$ fuori da $x_0$ non basta se manca la continuita' di $f$ in $x_0$: si puo' assegnare un valore isolato diverso a $f(x_0)$.\par\medskip
\noindent\textbf{quiz\_pagine\_261\_270\_001}: A. L'esponente e' $(3x)^x=3^xx^x$, non $3x^2$: la A e' falsa. Per derivare si usa la derivazione logaritmica.\par\medskip
\noindent\textbf{quiz\_pagine\_261\_270\_002}: D. $(x^{3x})^x=x^{3x^2}$; la D ha esponente $(3x)^x$ ed e' in generale diversa.\par\medskip
\noindent\textbf{quiz\_pagine\_261\_270\_023}: B. Per definizione $f\in C^k(I)$ quando le derivate fino all'ordine $k$ esistono e sono continue su $I$.\par\medskip
\noindent\textbf{quiz\_pagine\_261\_270\_028}: E. Senza continuita' in $x_0$, l'esistenza del limite di $f'$ non basta: si puo' modificare il solo valore $f(x_0)$ creando una discontinuita'.\par\medskip
\noindent\textbf{quiz\_pagine\_271\_280\_009}: D. La derivata di $2^{u(x)}$ e' $2^{u(x)}\log2\,u'(x)$ con $u=x\sin x$.\par\medskip
\noindent\textbf{quiz\_pagine\_271\_280\_010}: A. $g(x_0)=0$ e $\frac{g(x)-g(x_0)}{x-x_0}=f(x)\to f(x_0)=0$: dunque $g'(x_0)=0$.\par\medskip
\noindent\textbf{quiz\_pagine\_271\_280\_011}: C. La derivata dell'inversa in $f(2)=1$ vale $1/f'(2)=1/3$.\par\medskip
\noindent\textbf{quiz\_pagine\_271\_280\_014}: D. Per il teorema di Lagrange esiste $c\in(-1,6)$ con $f'(c)=(f(6)-f(-1))/(6-(-1))=3/7$.\par\medskip
\noindent\textbf{quiz\_pagine\_271\_280\_017}: C. $2x/(1+x)=2x-2x^2+2x^3+o(x^3)$ e $-1/(x-1)=1+x+x^2+x^3+o(x^3)$.\par\medskip
\noindent\textbf{quiz\_pagine\_281\_290\_004}: C. $\cosh(x^2)=1+x^4/2+O(x^8)$, percio' $f(x)=1+x^4/8+O(x^8)$: le derivate di ordine 1, 2 e 3 si annullano, la quarta no.\par\medskip
\noindent\textbf{quiz\_pagine\_281\_290\_005}: A. La funzione $h(x)=f(x)-3x+4$ ha tre zeri distinti. Per Rolle, tra due zeri consecutivi vi sono almeno due punti distinti con $h'=0$, cioe' $f'=3$.\par\medskip
\noindent\textbf{quiz\_pagine\_281\_290\_020}: E. $\log(1+2x)=2x+O(x^2)$ e $\cos u=1-u^2/2+O(u^4)$, quindi $f(x)=1-2x^2+O(x^3)$.\par\medskip
\noindent\textbf{quiz\_pagine\_281\_290\_023}: D. $e^x=1+x+x^2/2+x^3/6+o(x^3)$, $x\sin^2x=x^3+o(x^3)$: il polinomio e' $1+x+x^2/2-5x^3/6$, assente tra le opzioni.\par\medskip
\noindent\textbf{quiz\_pagine\_291\_300\_021}: A. $g(0)=1$ e $g(2)=0$, dunque per Lagrange $g'(c)=(0-1)/(2-0)=-1/2$.\par\medskip
\noindent\textbf{quiz\_pagine\_291\_300\_040}: B. $f$ si annulla in $0,1,2,3$; per Rolle, $f'$ ha almeno uno zero in ciascuno dei tre intervalli adiacenti.\par\medskip
\noindent\textbf{quiz\_pagine\_291\_300\_041}: A. $g(1)=\sin2\pi=0=g(2)=\sin\pi$; per Rolle esiste $c$ interno con $g'(c)=0$.\par\medskip
\noindent\textbf{quiz\_pagine\_301\_310\_001}: C. Per il teorema di Lagrange esiste $c$ con $f\prime(c)=(8-3)/(4-1)=5/3$.\par\medskip
\noindent\textbf{quiz\_pagine\_311\_320\_008}: D. $f\prime(x)=5x^4+\cosh x>0$: $f$ e strettamente crescente e continua. Il valore richiesto e minore di $f(10)$, quindi viene assunto in un unico $c<10$.\par\medskip
\noindent\textbf{quiz\_pagine\_311\_320\_009}: A. $g\prime=3f^2f\prime$ si annulla in $0$ e $1$ e, per Rolle, in almeno un punto di $(0,1)$.\par\medskip
\noindent\textbf{quiz\_pagine\_311\_320\_017}: B. Per Lagrange esiste $c\in(4,6)$ con $f\prime(c)=(0-6)/(6-4)=-3$.\par\medskip
\noindent\textbf{quiz\_pagine\_311\_320\_023}: E. $f\prime(x)=5x^4+3e^x>0$, quindi esiste un unico zero; $f(0)=1>0$ e $f(x)\to-\infty$ a sinistra, percio lo zero e negativo.\par\medskip
\noindent\textbf{quiz\_pagine\_311\_320\_034}: D. $f\prime(x)=3(x-1)(x-3)$; in $[-2,2]$ l\textbackslash{}'unico critico e 1. I valori sono $f(-2)=-49$, $f(1)=5$, $f(2)=3$.\par\medskip
\noindent\textbf{quiz\_pagine\_321\_330\_008}: D. $3^{2\sin x}=\exp(2\log3\sin x)$; usando $\sin x=x+o(x^2)$ si ottiene lo sviluppo D.\par\medskip
\noindent\textbf{quiz\_pagine\_321\_330\_011}: C. Se $g$ \textbackslash{}`e costante, la derivata \textbackslash{}`e nulla ovunque; altrimenti un valore non nullo \textbackslash{}`e accompagnato da un estremo interno rispetto ai limiti nulli agli infiniti.\par\medskip
\noindent\textbf{quiz\_pagine\_321\_330\_043}: C. Da $f(t)=f(0)+f\prime(0)t+o(t)$, sostituendo $t=x^2$ si ottiene $g(x)=f(0)+f\prime(0)x^2+o(x^2)$.\par\medskip
\noindent\textbf{quiz\_pagine\_331\_340\_001}: B. Per Lagrange le due derivate assumono in punti (eventualmente diversi) lo stesso valore medio, poich\textbackslash{}`e $f(6)-f(0)=g(6)-g(0)$. Anche Rolle su $f-g$ assicura un punto con $f\prime(c)=g\prime(c)$, non necessariamente nullo. La tabella ufficiale riporta ambiguamente ``ab``; solo B \textbackslash{}`e necessariamente vera.\par\medskip
\noindent\textbf{quiz\_pagine\_331\_340\_002}: B. Si usa $(1+x^4)^{1/3}=1+x^4/3+o(x^4)$ e $(1+x^3)^{-1/4}=1-x^3/4+o(x^4)$; il prodotto fornisce B.\par\medskip
\noindent\textbf{quiz\_pagine\_331\_340\_012}: D. Dallo sviluppo segue $f(5)=3$ e $f\prime(5)=2$: la tangente \textbackslash{}`e $y=3+2(x-5)$.\par\medskip
\noindent\textbf{quiz\_pagine\_331\_340\_014}: D. $F\prime=\arctan x$ e $F\prime\prime=1/(1+x^2)>0$ su tutto $\mathbb R$.\par\medskip
\noindent\textbf{quiz\_pagine\_331\_340\_020}: C. $f\prime(x)=9x^8+e^x>0$; inoltre $f$ \textbackslash{}`e continua e suriettiva su $\mathbb R$ per i limiti agli estremi: ogni valore \textbackslash{}`e assunto una sola volta.\par\medskip
\noindent\textbf{quiz\_pagine\_331\_340\_025}: A. Lo sviluppo corretto \textbackslash{}`e $1+x+x^2/2+x^3/3+o(x^3)$, ossia A. La soluzione ufficiale non compare nelle pagine fornite (la simulazione prosegue nel blocco successivo).\par\medskip
\noindent\textbf{quiz\_pagine\_341\_350\_003}: D. Il coefficiente di $(x+2)^4$ nello sviluppo e $f^{(4)}(-2)/4!=-4$. Il termine negativo determina un massimo relativo.\par\medskip
\noindent\textbf{quiz\_pagine\_341\_350\_017}: E. $\sin x=x-x^3/6+o(x^3)$; espandendo $e^t$, i termini cubici $-x^3/6$ e $+x^3/6$ si cancellano.\par\medskip
\noindent\textbf{quiz\_pagine\_341\_350\_018}: D. $2^{\cos x}=2\exp[(\cos x-1)\log2]=2[1-(\log2)x^2/2+o(x^2)]$.\par\medskip
\noindent\textbf{quiz\_pagine\_341\_350\_021}: D. $f(-3)=f(-2)=0$ e $f$ e continua e derivabile: Rolle garantisce $c\in(-3,-2)$ con $f\prime(c)=0$.\par\medskip
\noindent\textbf{quiz\_pagine\_341\_350\_024}: E. $f(-5)=-2$, $f\prime(-5)=1>0$ e $f\prime\prime(-5)=-6<0$. Le proprieta valgono in un intorno sufficientemente piccolo.\par\medskip
\noindent\textbf{quiz\_pagine\_341\_350\_032}: E. Fermat richiede un estremo locale interno e derivabilita nel punto: allora la derivata e zero. C e una tautologia, non l\textbackslash{}textquotesingle enunciato del teorema.\par\medskip
\noindent\textbf{quiz\_pagine\_351\_360\_003}: D. Il coefficiente di $(x+2)^4$ nello sviluppo e $f^{(4)}(-2)/4!=-4$. Il termine negativo determina un massimo relativo.\par\medskip
\noindent\textbf{quiz\_pagine\_351\_360\_017}: E. $\sin x=x-x^3/6+o(x^3)$; espandendo $e^t$, i termini cubici $-x^3/6$ e $+x^3/6$ si cancellano.\par\medskip
\noindent\textbf{quiz\_pagine\_351\_360\_018}: D. $2^{\cos x}=2\exp[(\cos x-1)\log2]=2[1-(\log2)x^2/2+o(x^2)]$.\par\medskip
\noindent\textbf{quiz\_pagine\_351\_360\_021}: D. $f(-3)=f(-2)=0$ e $f$ e continua e derivabile: Rolle garantisce $c\in(-3,-2)$ con $f\prime(c)=0$.\par\medskip
\noindent\textbf{quiz\_pagine\_351\_360\_024}: E. $f(-5)=-2$, $f\prime(-5)=1>0$ e $f\prime\prime(-5)=-6<0$. Le proprieta valgono in un intorno sufficientemente piccolo.\par\medskip
\noindent\textbf{quiz\_pagine\_351\_360\_032}: E. Fermat richiede un estremo locale interno e derivabilita nel punto: allora la derivata e zero. C e una tautologia, non l\textbackslash{}textquotesingle enunciato del teorema.\par\medskip
\noindent\textbf{quiz\_pagine\_361\_370\_005}: E. Rolle richiede continuita sull\textbackslash{}textquotesingle intervallo chiuso, derivabilita nell\textbackslash{}textquotesingle intervallo aperto e uguaglianza dei valori agli estremi. Anche A e C sono formulazioni con ipotesi piu forti che implicano la conclusione; la risposta ufficiale e E.\par\medskip
\noindent\textbf{quiz\_pagine\_361\_370\_012}: E. $f$ e continua su $[0,2]$, derivabile su $(0,2)$ e $f(0)=f(2)=0$. $f\prime(x)=\frac{2-3x}{2\sqrt x}$ si annulla in $2/3$.\par\medskip
\noindent\textbf{quiz\_pagine\_361\_370\_030}: D. $e^{2x+x^2}=1+(2x+x^2)+\frac12(2x)^2+o(x^2)=1+2x+3x^2$. Le opzioni A e B sono identiche nella fonte.\par\medskip
\noindent\textbf{quiz\_pagine\_371\_380\_004}: A. Usando $\log(1+u)=u-u^2/2+u^3/3+o(u^3)$ con $u=x+x^2$ si ottiene $x+x^2/2-2x^3/3+o(x^3)$.\par\medskip
\noindent\textbf{quiz\_pagine\_371\_380\_012}: E. Il valore iniziale e $f(0)=e$. Sviluppando $\sqrt{1-x}=1-x/2-x^2/8-x^3/16+\cdots$ e poi l\textbackslash{}textquotesingle esponenziale si ottiene $e(1-x/2-x^3/48+o(x^3))$. La risposta E nel PDF appare senza il fattore $e$: possibile refuso tipografico.\par\medskip
\noindent\textbf{quiz\_pagine\_371\_380\_013}: D. $f(-2)=f(0)=0$ e $f$ e continua e derivabile su $[-2,0]$; Rolle garantisce un $c\in(-2,0)$ con $f\prime(c)=0$.\par\medskip
\noindent\textbf{quiz\_pagine\_371\_380\_018}: E. Il coefficiente di $(x-3)^3$ e nullo e il resto e $o((x-3)^3)$, quindi $f\prime\prime\prime(3)=0$ se la derivata terza esiste. La scrittura implica anche $f\prime\prime(3)=6$, non $3/2!$.\par\medskip
\noindent\textbf{quiz\_pagine\_371\_380\_026}: E. Il teorema degli zeri richiede continuita e segni opposti agli estremi, e garantisce uno zero interno. Anche A e corretta nella formulazione con prodotto non positivo e zero eventualmente agli estremi: il quiz ha piu risposte valide.\par\medskip
\noindent\textbf{quiz\_pagine\_391\_400\_004}: E. $(e^x-1)^2=x^2+o(x^2)$ e $\sin x=x+o(x^2)$, dunque il polinomio di Taylor di grado due e $x+x^2$ (senza resto).\par\medskip
\noindent\textbf{quiz\_pagine\_391\_400\_014}: B. Il termine dominante $2(x+3)^3$ cambia segno attraversando $-3$. La derivata terza, se esiste, sarebbe $12$ e non $2$.\par\medskip
\noindent\textbf{quiz\_pagine\_401\_410\_002}: A. La radice cubica vale $x(1+x)^{1/3}=x+\frac13x^2+o(x^2)$, quindi $f(x)\sim\frac23x^2$.\par\medskip
\noindent\textbf{quiz\_pagine\_401\_410\_004}: D. Dallo sviluppo si leggono $f(\pi)=\pi$ e $f\prime(\pi)=2$, perci\textbackslash{}`o la tangente \textbackslash{}`e $y=\pi+2(x-\pi)$.\par\medskip
\noindent\textbf{quiz\_pagine\_401\_410\_008}: A. Lo sviluppo di Maclaurin di $\cos x$ contiene solo potenze pari; tutte le derivate di ordine dispari in zero si annullano. Attenzione: anche D \textbackslash{}`e vera, ma descrive solo un sottoinsieme degli ordini dispari, quindi il quesito non \textbackslash{}`e univoco in senso stretto.\par\medskip
\noindent\textbf{quiz\_pagine\_401\_410\_011}: E. Il prodotto $(1-x^2/2+O(x^4))(1+x^2/2+O(x^4))=1+O(x^4)=1+o(x^2)$. Anche B, letta letteralmente, e corretta perche il resto $o(x^2)$ puo assorbire un termine di ordine quattro; la risposta ufficiale E e la forma canonica.\par\medskip
\noindent\textbf{quiz\_pagine\_411\_420\_010}: C. $\sin x=x-x^3/6+o(x^3)$ e $x^3\cos x=x^3+o(x^3)$, quindi il coefficiente cubico \textbackslash{}`e $-7/6$.\par\medskip
\noindent\textbf{quiz\_pagine\_411\_420\_020}: E. Si scrive $\log(1+\sin x/3)$ e si sviluppa: $x/3-x^2/18-7x^3/162+o(x^3)$.\par\medskip
\noindent\textbf{quiz\_pagine\_411\_420\_027}: A. $f(-2\pi)=\pi$ e $f\prime(-2\pi)=-2$, quindi $y-\pi=-2(x+2\pi)$, ossia $y=-2x-3\pi$.\par\medskip
\noindent\textbf{quiz\_pagine\_411\_420\_028}: E. $\cos x=1-x^2/2+O(x^4)$, $\cosh x=1+x^2/2+O(x^4)$: il polinomio di secondo grado \textbackslash{}`e $1$. Anche B rappresenta correttamente lo sviluppo con resto, pur non essendo il solo polinomio.\par\medskip
\noindent\textbf{quiz\_pagine\_411\_420\_036}: A. $4^x=e^{x\log4}$: lo sviluppo \textbackslash{}`e $1+(\log4)x+(\log4)^2x^2/2+(\log4)^3x^3/6+o(x^3)$.\par\medskip
\noindent\textbf{quiz\_pagine\_421\_430\_016}: D. $2e^{-x}-x^2=2-2x+o(x^2)$ e $(1+x^2)^{-1}=1-x^2+o(x^2)$, quindi il polinomio \textbackslash{}`e $2-2x-2x^2$. Anche C \textbackslash{}`e lo sviluppo completo con resto.\par\medskip
\noindent\textbf{quiz\_pagine\_431\_440\_003}: E. $(e^x-1)=x+x^2/2+x^3/6+o(x^3)$ e $1/\cos x=1+x^2/2+o(x^3)$; coefficiente cubico $1/6+1/2=2/3$.\par\medskip
\noindent\textbf{quiz\_pagine\_431\_440\_007}: A. Il termine cubico \textbackslash{}`e $\sqrt3/2-(\sqrt3)^3/6=0$. Pertanto $f(x)=\sqrt3x+\sqrt3x^2+o(x^3)$.\par\medskip
\noindent\textbf{quiz\_pagine\_431\_440\_026}: A. La risposta ufficiale A deriva dalla derivazione formale dello sviluppo. In generale derivare un resto $o(x^2)$ non garantisce un resto $o(x)$ senza ipotesi aggiuntive: quesito matematicamente non rigoroso.\par\medskip
\noindent\textbf{quiz\_pagine\_431\_440\_029}: D. $\sinh t=t+t^3/6+o(t^3)$ con $t=x-x^3/6$; i termini cubici si cancellano. Anche C ed E sono affermazioni vere ma meno precise: ambiguit\textbackslash{}`a della fonte.\par\medskip
\noindent\textbf{quiz\_pagine\_441\_450\_002}: A. $\sin^2 t=t^2-t^4/3+o(t^4)$, dunque $\sin^2(3x)=9x^2-27x^4+o(x^4)$.\par\medskip
\noindent\textbf{quiz\_pagine\_441\_450\_006}: D. Teorema di prolungamento della derivata: continuit\textbackslash{}`a nel punto e limite finito delle derivate implicano derivabilit\textbackslash{}`a con tale derivata. Anche B, se il limite esiste in senso esteso, richiede precisazioni.\par\medskip
\noindent\textbf{quiz\_pagine\_441\_450\_020}: A. $f$ \textbackslash{}`e continua su $[-1,1]$ e quindi soddisfa Weierstrass; non \textbackslash{}`e derivabile in $0$ ($f\prime_-(0)=1$, $f\prime_+(0)=-1$). NOTA: anche D risulta vera se per teorema degli zeri si intendono continuit\textbackslash{}`a e valori agli estremi di segno opposto? Qui $f(-1)=f(1)=0$, quindi non di segno opposto.\par\medskip
\noindent\textbf{quiz\_pagine\_451\_460\_001}: E. $f(-5)=-10$ e $f\prime(-5)=-2$: la tangente \textbackslash{}`e $y=-10-2(x+5)=-2x-20$, non quella proposta. La derivata non nulla esclude estremi locali e crescita in un intorno.\par\medskip
\noindent\textbf{quiz\_pagine\_451\_460\_012}: D. $\log(1+x)=x-x^2/2+O(x^3)$ e $\sin^2\sqrt{x}=x-x^2/3+O(x^3)$, perci\textbackslash{}`o il prodotto \textbackslash{}`e $x^2-5x^3/6+O(x^4)$. Sottraendo $x^2$ resta $-5x^3/6$.\par\medskip
\noindent\textbf{quiz\_pagine\_451\_460\_023}: B. $\cos(x^3)=1-x^6/2+O(x^{12})$; quindi $e^{\cos(x^3)}=e(1-x^6/2)+O(x^{12})$.\par\medskip
\noindent\textbf{quiz\_pagine\_451\_460\_027}: E. Per $x<0$, $\sqrt{4x^2+2x^3}=-2x\sqrt{1+x/2}=-2x+O(x^2)$, quindi $f(x)=x+O(x^2)$.\par\medskip
\noindent\textbf{quiz\_pagine\_461\_470\_008}: B. $f$ \textbackslash{}`e continua sui segmenti $[-\pi/2,\pi/2]$ e $[\pi/2,3\pi/2]$ e si annulla agli estremi. Rolle garantisce uno zero di $f\prime$ in ciascuno dei due intervalli aperti, entrambi contenuti nell\textbackslash{}textquotesingle intervallo dato.\par\medskip
\noindent\textbf{quiz\_pagine\_461\_470\_010}: D. Per Lagrange $g\prime(c)=[g(2)-g(0)]/2=(e^{-3}-1)/2$. ATTENZIONE: la soluzione ufficiale indica C, ma la sua stessa dimostrazione calcola $(e^{-3}-1)/2$, cio\textbackslash{}`e D.\par\medskip
\noindent\textbf{quiz\_pagine\_471\_480\_008}: E. Sviluppando $\sin x=x+O(x^3)$ si ha $f(x)/x=ax^2+3x^3+O(x^4)$. Per $a=3$ il polinomio di terzo grado \textbackslash{}`e $3x^2+3x^3$.\par\medskip
\noindent\textbf{quiz\_pagine\_471\_480\_011}: B. Poich'e $o(x^5)/x^4=o(x)\to0$, vicino a zero $f(x)=x^4(1+o(x))\ge0$ e $f(0)=0$, dunque \textbackslash{}`e minimo locale.\par\medskip
\noindent\textbf{quiz\_pagine\_481\_490\_024}: C. $f(0)=2$ e $f\prime(0)=3+1=4$. Per il teorema della derivata dell\textbackslash{}'inversa, $(f^{-1})\prime(2)=1/4$.\par\medskip
\noindent\textbf{quiz\_pagine\_501\_510\_009}: B. Poich\textbackslash{}'e $f(0)=1>0$, $\log f$ \textbackslash{}`e definita vicino a zero e $h\prime(0)=f\prime(0)/f(0)=2$. Non \textbackslash{}`e garantita la positivit\textbackslash{}`a di $f$ su tutto $[0,4]$.\par\medskip
\noindent\textbf{quiz\_pagine\_501\_510\_010}: E. Per Lagrange $g\prime(c)=[g(3)-g(0)]/3=(e^{-2}-1)/3$.\par\medskip
\noindent\textbf{quiz\_pagine\_501\_510\_034}: A. $g\prime(x)=4f(x)^3f\prime(x)$, quindi si annulla in $0$ e in $1$. Per Rolle applicato a $f$ esiste inoltre $c\in(0,1)$ con $f\prime(c)=0$: almeno tre zeri distinti.\par\medskip
\noindent\textbf{quiz\_pagine\_511\_520\_009}: C. Il coefficiente della secante \textbackslash{}`e $(f(3)-f(-1))/4=9$; $f'(c)=3c^2+2=9$ d\textbackslash{}`a $c=\pm\sqrt{7/3}$, ma solo quello positivo appartiene a $(-1,3)$.\par\medskip
\noindent\textbf{quiz\_pagine\_511\_520\_015}: E. Per il binomio generalizzato $(1+x)^{1/2}=1+x/2-x^2/8+x^3/16+o(x^3)$. L'item A riporta il solo polinomio senza resto.\par\medskip
\noindent\textbf{quiz\_pagine\_511\_520\_016}: C. La risposta ufficiale \textbackslash{}`e C: il termine cubico implica un andamento della derivata seconda come $-42x$, quindi decrescente, se \textbackslash{}`e lecito derivare il resto. Nota: dal solo resto $o(x^4)$ non segue in generale la monotonia locale di $f''$; il quesito sottintende regolarit\textbackslash{}`a aggiuntiva.\par\medskip
\noindent\textbf{quiz\_pagine\_511\_520\_034}: D. Il teorema di Lagrange garantisce $f'(c)=[f(6)-f(-1)]/[6-(-1)]=3/7$ per almeno un $c\in(-1,6)$.\par\medskip
\noindent\textbf{quiz\_pagine\_511\_520\_037}: C. $2x/(1+x)=2x-2x^2+2x^3+o(x^3)$ e $-1/(x-1)=1+x+x^2+x^3+o(x^3)$; sommando si ottiene C.\par\medskip
\subsubsection*{4.3 -- Proprietà delle funzioni: monotonia e punti estremali}
\noindent\textbf{quiz\_pagine\_001\_010\_013}: A. Spiegazione non ancora disponibile.\par\medskip
\noindent\textbf{quiz\_pagine\_001\_010\_014}: B. Spiegazione non ancora disponibile.\par\medskip
\noindent\textbf{quiz\_pagine\_001\_010\_024}: A. Spiegazione non ancora disponibile.\par\medskip
\noindent\textbf{quiz\_pagine\_001\_010\_032}: E. Spiegazione non ancora disponibile.\par\medskip
\noindent\textbf{quiz\_pagine\_011\_020\_024}: B. $(f^{-1})\prime(f(2))=1/f\prime(2)=7$.\par\medskip
\noindent\textbf{quiz\_pagine\_021\_030\_003}: D. $f\prime=3(x-1)(x-3)$: nell\textbackslash{}'intervallo il punto critico \textbackslash{}`e $1$. $f(-2)=-49$, $f(1)=5$, $f(2)=3$.\par\medskip
\noindent\textbf{quiz\_pagine\_031\_040\_004}: D. Spiegazione non ancora disponibile.\par\medskip
\noindent\textbf{quiz\_pagine\_041\_050\_009}: E. Risposta ufficiale indicata nella tabella della decima simulazione.\par\medskip
\noindent\textbf{quiz\_pagine\_041\_050\_019}: C. Risposta ufficiale indicata nella tabella della decima simulazione.\par\medskip
\noindent\textbf{quiz\_pagine\_041\_050\_043}: non verificata / nessuna alternativa corretta. Soluzione non riportata nel PDF; da verificare.\par\medskip
\noindent\textbf{quiz\_pagine\_051\_060\_011}: D. Il termine $1/(1+x^2)$ raggiunge il massimo $1$ in $x=0$.\par\medskip
\noindent\textbf{quiz\_pagine\_051\_060\_013}: C. Poiché $1/(1+x^2)>0$ e tende a zero per $|x|	o\infty$, l'estremo inferiore è $\log2$.\par\medskip
\noindent\textbf{quiz\_pagine\_061\_070\_024}: non verificata / nessuna alternativa corretta. Classificazione basata su estremi relativi; soluzione non riportata nelle pagine fornite. Notazione dei limiti con segni accoppiati da verificare.\par\medskip
\noindent\textbf{quiz\_pagine\_081\_090\_020}: B. $f\prime(x)=3^x\log3+3x^2>0$, dunque $f$ è strettamente crescente. Per continuità e perché l\textbackslash{}textquotesingle immagine è $\mathbb R$, il valore indicato è assunto esattamente una volta.\par\medskip
\noindent\textbf{quiz\_pagine\_091\_100\_007}: A. $f\prime(x)=5x^4+1>0$, quindi ha un unico zero; poiché $f(-2)<0<f(2)$, lo zero è in $(-2,2)$.\par\medskip
\noindent\textbf{quiz\_pagine\_091\_100\_011}: D. $f\prime(x)=7x^6+7e^{7x}>0$ e $f$ è suriettiva su $\mathbb R$. Poiché $\bar y=f(5)+2>f(5)$, la soluzione unica soddisfa $x_0>5$.\par\medskip
\noindent\textbf{quiz\_pagine\_101\_110\_017}: non verificata / nessuna alternativa corretta. La tabella indica A e i coefficienti suggeriscono $f(2)=2$, $f\prime(2)=6$, $f\prime\prime(2)=2$. Tuttavia il solo resto $o((x-2)^2)$ non garantisce convessità né monotonia in un intero intorno senza ulteriori ipotesi di regolarità: soluzione rigorosamente non determinabile.\par\medskip
\noindent\textbf{quiz\_pagine\_111\_120\_005}: C. Se il minimo esiste, non puo' superare il valore $f(0)=-6$. Senza continuita' le altre proprieta' non sono garantite.\par\medskip
\noindent\textbf{quiz\_pagine\_111\_120\_011}: A. Poiche' $f(x)=-18x^3+o(x^3)$, si ha $f'(0)=f''(0)=0$ e $f'''(0)<0$: flesso a tangente orizzontale decrescente.\par\medskip
\noindent\textbf{quiz\_pagine\_111\_120\_015}: B. Dallo sviluppo si ricavano $f(-1)=e>0$, $f'(-1)=e>0$, $f''(-1)=8e>0$: localmente positiva, crescente e convessa.\par\medskip
\noindent\textbf{quiz\_pagine\_111\_120\_024}: E. Sul tratto $(1,6)$, i limiti $+\infty$ ai due estremi assicurano un minimo interno per continuita'. Questo e' un minimo relativo alla componente, non necessariamente un minimo globale sull'intero dominio: il testo originale e' ambiguo.\par\medskip
\noindent\textbf{quiz\_pagine\_111\_120\_035}: E. $f(2)=e>0$, $f'(2)=e>0$, $f''(2)=-4e<0$: localmente positiva, crescente e concava (assumendo regolarita' sufficiente per dedurre il segno della concavita').\par\medskip
\noindent\textbf{quiz\_pagine\_131\_140\_004}: D. Per $x\neq0$ si ha $e^{|x|}>1$, con estremo inferiore $1$ non raggiunto; un minimo esiste esattamente quando $k\le1$.\par\medskip
\noindent\textbf{quiz\_pagine\_131\_140\_023}: non verificata / nessuna alternativa corretta. $f\prime=e^{x^3-ax^2}(1+3x^3-2ax^2)$; il polinomio cubico tende a $-\infty$ per $x\to-\infty$, quindi la funzione non puo\textbackslash{}'e essere crescente su tutta $\mathbb R$ per alcun $a$. Nessuna risposta corretta: verificare la formula nel PDF.\par\medskip
\noindent\textbf{quiz\_pagine\_141\_150\_015}: E. Il vincolo $f(x)\geq e^{2x}$ non impone segni alle derivate: ad esempio si puo\textbackslash{}'e sommare a $e^{2x}$ una funzione liscia non negativa e oscillante.\par\medskip
\noindent\textbf{quiz\_pagine\_141\_150\_021}: A. Per $x<0$ i valori sono in $(0,1]$ e l\textbackslash{}'estremo inferiore $0$ non e\textbackslash{}'e raggiunto. Il minimo esiste esattamente se $a\leq0$.\par\medskip
\noindent\textbf{quiz\_pagine\_141\_150\_022}: D. $\cosh x$ ha minimo assoluto in $x=0$; il dominio $[-2,a)$ contiene $0$ esattamente quando $a>0$.\par\medskip
\noindent\textbf{quiz\_pagine\_171\_180\_002}: A. Per continuita' $f$ si annulla tra $-1$ e $1$; quindi $f^2\geq0$ raggiunge il minimo assoluto $0$. Le altre proprieta' non sono generali.\par\medskip
\noindent\textbf{quiz\_pagine\_171\_180\_007}: A. Fermat: $F'(x_0)=0$ per un massimo interno; poiche' $F'=f$, si ottiene $f(x_0)=0$.\par\medskip
\noindent\textbf{quiz\_pagine\_171\_180\_016}: E. La definizione di massimo locale e' espressa da (e). Anche (c) e' logicamente conseguenza della definizione, perche' escludere $x_0$ non invalida la disuguaglianza: il quiz presenta due affermazioni vere.\par\medskip
\noindent\textbf{quiz\_pagine\_181\_190\_007}: B. La funzione e' strettamente crescente per $x\geq1$ (tutti i fattori positivi e crescenti), dunque ha minimo in $1$, ma non raggiunge il supremo al bordo escluso $2$.\par\medskip
\noindent\textbf{quiz\_pagine\_181\_190\_015}: A. Per $x\geq2$, $f(x)=-(x-2)e^{-x/3}$ e $f'(x)=e^{-x/3}(x-5)/3$: decresce fino a $5$, poi cresce. L'intervallo massimale e' $[5,+\infty)$.\par\medskip
\noindent\textbf{quiz\_pagine\_181\_190\_016}: non verificata / nessuna alternativa corretta. $f''(x)=30x^4-60x^2=30x^2(x^2-2)$ cambia segno in $\pm\sqrt2$, non in $0$. Nelle alternative B e C il PDF riporta $\sqrt2$ due volte senza il segno meno: l'intenzione sembra essere $\{-\sqrt2,\sqrt2\}$; la soluzione ufficiale B non coincide letteralmente con il testo.\par\medskip
\noindent\textbf{quiz\_pagine\_181\_190\_022}: A. $f'(x)=[2+1/x-2\log x]/(2x+1)^2$. Il numeratore e' strettamente decrescente, passa da positivo a negativo, dunque esiste un unico massimo.\par\medskip
\noindent\textbf{quiz\_pagine\_191\_200\_001}: C. $\cosh x=1+x^2/2+O(x^4)$ e $x^3\cos x=x^3+O(x^5)$; il polinomio richiesto e' $2+x^2/2+x^3$. La A contiene anche il resto, non e' soltanto il polinomio.\par\medskip
\noindent\textbf{quiz\_pagine\_191\_200\_014}: E. $x\sinh3x=3x^2+\frac{27}6x^4+o(x^4)$, $e^{-3x^2}=1-3x^2+\frac92x^4+o(x^4)$; somma $1+9x^4+o(x^4)$.\par\medskip
\noindent\textbf{quiz\_pagine\_191\_200\_015}: E. $f(x)=1+9x^4+o(x^4)$: $f(0)=1$, dunque la funzione non e' infinitesima; lo e' $f(x)-1$.\par\medskip
\noindent\textbf{quiz\_pagine\_191\_200\_019}: non verificata / nessuna alternativa corretta. $\frac{3x}{1+4x}=3x-12x^2+48x^3+O(x^4)$; sommando gli altri termini si ottiene $68x^3+O(x^4)$, non $28x^3$. ATTENZIONE: nessuna risposta coincide; la risposta ufficiale D sembra errata.\par\medskip
\noindent\textbf{quiz\_pagine\_201\_210\_002}: A. La radice vale $1+2/x+O(x^{-2})$, mentre $\cos(1/x)=1+O(x^{-2})$. Il prodotto e' $2x+O(1)\to+\infty$.\par\medskip
\noindent\textbf{quiz\_pagine\_201\_210\_005}: A. $f(x)=2+\log5-x^2/5+O(x^4)$, percio' $x=0$ e' un massimo locale.\par\medskip
\noindent\textbf{quiz\_pagine\_201\_210\_008}: B. $f(x)=2+\log5-x^3/5+O(x^6)$: non tende a zero, dunque non e' infinitesima.\par\medskip
\noindent\textbf{quiz\_pagine\_201\_210\_011}: E. Per $x\to0^+$, $x^3\sin(1/x)=o(\sqrt x)$ e $\sqrt{3\sin x}\sim\sqrt{3x}$. L'espressione e' reale solo dove $\sin x\ge0$ vicino a zero.\par\medskip
\noindent\textbf{quiz\_pagine\_201\_210\_012}: D. $\cos^2x=1-x^2+x^4/3+O(x^6)$ e $\sqrt{1-2x^2}=1-x^2-x^4/2+O(x^6)$: differenza $5x^4/6+O(x^6)$.\par\medskip
\noindent\textbf{quiz\_pagine\_201\_210\_014}: B. Il numeratore e' equivalente a $x^{1/4}$, il rapporto a $x^{-5/4}\to+\infty$.\par\medskip
\noindent\textbf{quiz\_pagine\_201\_210\_016}: D. $\sqrt[3]{1+5x}=1+5x/3+O(x^2)$, dunque il limite e' $(5/3)/7=5/21$.\par\medskip
\noindent\textbf{quiz\_pagine\_201\_210\_019}: A. Per $x\to\pm\infty$, $\sqrt[3]{x^3+x^4}\sim x^{4/3}$; per $f_3$ si usa $\sqrt[3]{x^3+x^2}=x+1/3+o(1)$; per $x\to0$ gli sviluppi danno $f_4\sim x$ e $f_5\sim-x^2/3$.\par\medskip
\noindent\textbf{quiz\_pagine\_231\_240\_005}: D. $f\prime(x)=-1/x-2x<0$ su $(1,2)$. Gli estremi dell'intervallo sono esclusi, quindi non si raggiungono massimo o minimo.\par\medskip
\noindent\textbf{quiz\_pagine\_231\_240\_013}: A. La stretta crescenza implica la crescenza non stretta. Non implica limitatezza ne' particolari proprieta' di concavita'.\par\medskip
\noindent\textbf{quiz\_pagine\_231\_240\_021}: C. La derivata di $e^{-3x}$ e' $-3e^{-3x}<0$. Le altre tre funzioni sono crescenti su $[-3,-1]$.\par\medskip
\noindent\textbf{quiz\_pagine\_241\_250\_003}: C. $f\prime\prime(x)=e^x(x^2+4x+2)$ e' negativa tra le radici $-2\pm\sqrt2$.\par\medskip
\noindent\textbf{quiz\_pagine\_241\_250\_007}: D. Teorema di Fermat: per un estremo locale interno di funzione derivabile la derivata si annulla. Non e' sufficiente: $f(x)=x^3$.\par\medskip
\noindent\textbf{quiz\_pagine\_241\_250\_008}: A. La riflessione orizzontale inverte la monotonia, mentre preserva la concavita'.\par\medskip
\noindent\textbf{quiz\_pagine\_241\_250\_027}: E. ATTENZIONE: il correttore originale indica C, ma la condizione non implica monotonia neppure vicino a zero. Controesempio continuo: $f(0)=0$, $f(x)=x(2+\sin(1/x))$ per $x>0$ e $f(x)=0$ per $x<0$. E' positiva per $x>0$ ma oscilla in ogni intorno destro di 0. Risposta matematica E, difforme dalla chiave originale C.\par\medskip
\noindent\textbf{quiz\_pagine\_251\_260\_002}: E. Le due branche hanno immagini $(-\infty,0]$ e $(0,+\infty)$: l'immagine e' $\mathbb R$. Inoltre le derivate laterali in $-5$ valgono entrambe $10$; dunque e' suriettiva e l'affermazione E e' falsa.\par\medskip
\noindent\textbf{quiz\_pagine\_261\_270\_003}: E. Per $|x|\leq1$, $f(x)=-2$, mentre per $|x|>1$, $f(x)=-2|x|$; quindi E e' falsa.\par\medskip
\noindent\textbf{quiz\_pagine\_261\_270\_004}: C. I due rami hanno segno opposto e non sono speculari rispetto alla retta $x=-5$. Le derivate laterali nel raccordo valgono entrambe zero.\par\medskip
\noindent\textbf{quiz\_pagine\_261\_270\_009}: E. $f'(x)=5x^4+3e^x>0$, $f(0)=1>0$ e $f(x)\to-\infty$ per $x\to-\infty$.\par\medskip
\noindent\textbf{quiz\_pagine\_261\_270\_016}: E. $\sin t=t$ per $t\geq0$ ha l'unica soluzione $t=0$. Pertanto l'unico zero e' $x=0$.\par\medskip
\noindent\textbf{quiz\_pagine\_261\_270\_021}: E. La sola condizione $f(x)>f(0)$ per $x>0$ non implica monotonia neppure vicino a zero: ad esempio $f(x)=x^2(2+\sin(1/x))$ per $x>0$, $f(0)=0$, con prolungamento continuo per $x<0$. La chiave ufficiale indica C, ma e' falsa.\par\medskip
\noindent\textbf{quiz\_pagine\_281\_290\_017}: E. Essendo derivabile e strettamente crescente, $g([a,b])=[-1,1]$: non assume il valore 2.\par\medskip
\noindent\textbf{quiz\_pagine\_291\_300\_004}: B. $f'(x)=9x^8+4>0$, quindi lo zero e' unico; $f(-1)=-2<0$ e $f(0)=3>0$.\par\medskip
\noindent\textbf{quiz\_pagine\_291\_300\_039}: D. La funzione e' crescente e $f(2)=24<32$, quindi non assume $32$ su $(-2,2)$.\par\medskip
\noindent\textbf{quiz\_pagine\_301\_310\_007}: B. Per Weierstrass il massimo assoluto esiste e, per definizione, e' almeno $f(2)$.\par\medskip
\noindent\textbf{quiz\_pagine\_301\_310\_012}: E. Il punto di minimo e' un elemento $x_0$ del dominio in cui il valore della funzione e' minore o uguale a tutti gli altri valori.\par\medskip
\noindent\textbf{quiz\_pagine\_321\_330\_010}: A. $g\prime=-\sin(f)f\prime$. Ci sono i tre zeri di $f$ e almeno due zeri distinti di $f\prime$ per Rolle; inoltre $\sin f=0$ pu\textbackslash{}`o annullarsi in altri punti.\par\medskip
\noindent\textbf{quiz\_pagine\_321\_330\_029}: E. Poich\textbackslash{}`e $0\in[-7,1]$ e $x^2\geq0$, il minimo \textbackslash{}`e zero.\par\medskip
\noindent\textbf{quiz\_pagine\_331\_340\_022}: D. $g\prime(x)=2f(x)f\prime(x)$, dunque $g\prime(0)=g\prime(10)=0$; per Rolle applicato a $g$ esiste anche $c\in(0,10)$ con $g\prime(c)=0$.\par\medskip
\noindent\textbf{quiz\_pagine\_361\_370\_011}: D. Il denominatore e minimo in $x=0$; dunque $f(0)=\log3$ e il massimo.\par\medskip
\noindent\textbf{quiz\_pagine\_361\_370\_013}: C. $f(x)>\log2$ per ogni $x$ e $f(x)\to\log2$ per $|x|\to\infty$.\par\medskip
\noindent\textbf{quiz\_pagine\_371\_380\_014}: E. La tabella ufficiale indica E. Le condizioni asintotiche vanno analizzate sui tre intervalli separati dai punti esclusi; da sole non determinano univocamente il numero dei massimi relativi.\par\medskip
\noindent\textbf{quiz\_pagine\_371\_380\_029}: C. Ponendo $h=x+4$, $f(-4+h)-f(-4)=3h^2+h^3+o(h^3)>0$ per $h$ piccolo e non nullo, quindi $-4$ e minimo locale.\par\medskip
\noindent\textbf{quiz\_pagine\_391\_400\_005}: B. $f\prime(x)=3+3x^2>0$ per ogni $x$ e $f$ ha immagine $\mathbb R$, quindi ogni valore e assunto esattamente una volta.\par\medskip
\noindent\textbf{quiz\_pagine\_391\_400\_018}: E. ATTENZIONE: risposta ufficiale E, ma anche A e falsa: vicino a $x=2$, $(x-2)^4\geq0$ e $(5-x)\arctan x>0$, dunque $x=2$ e un minimo locale, non un massimo. Il quesito non ha risposta univoca.\par\medskip
\noindent\textbf{quiz\_pagine\_391\_400\_019}: A. $f(-2)=-33<0$, $f(2)=35>0$; per Bolzano esiste uno zero in $(-2,2)$. Inoltre $f\prime(x)=5x^4+1>0$, quindi e unico.\par\medskip
\noindent\textbf{quiz\_pagine\_401\_410\_003}: D. $f\prime(x)=7x^6+7e^{7x}>0$; $f$ \textbackslash{}`e strettamente crescente e $\bar y=f(5)+2>f(5)$, perci\textbackslash{}`o la soluzione unica \textbackslash{}`e maggiore di $5$.\par\medskip
\noindent\textbf{quiz\_pagine\_401\_410\_016}: B. Per $x\to-\infty$, $e^x\to0$ e $\cos x$ oscilla; per continuit\textbackslash{}`a esistono infiniti cambi di segno e quindi infiniti zeri.\par\medskip
\noindent\textbf{quiz\_pagine\_411\_420\_017}: A. Dallo sviluppo si ricavano $f(2)=2$, $f\prime(2)=6$ e $f\prime\prime(2)=2$. La convessit\textbackslash{}`a in un intero intorno non segue tuttavia dal solo resto $o((x-2)^2)$ senza ulteriori ipotesi: la fonte assume implicitamente regolarit\textbackslash{}`a.\par\medskip
\noindent\textbf{quiz\_pagine\_411\_420\_025}: C. Se il minimo esiste, non pu\textbackslash{}`o superare il valore $f(0)=-6$. Senza continuit\textbackslash{}`a non si pu\textbackslash{}`o applicare il teorema dei valori intermedi.\par\medskip
\noindent\textbf{quiz\_pagine\_411\_420\_031}: A. Poich\textbackslash{}'e $f(x)=-18x^3+o(x^3)$, la tangente \textbackslash{}`e orizzontale e la funzione attraversa lo zero con andamento decrescente.\par\medskip
\noindent\textbf{quiz\_pagine\_411\_420\_035}: B. $f(-1)=e>0$, $f\prime(-1)=e>0$, $f\prime\prime(-1)=8e>0$. La convessit\textbackslash{}`a in un intorno richiede per\textbackslash{}`o regolarit\textbackslash{}`a ulteriore rispetto al solo resto $o((x+1)^2)$.\par\medskip
\noindent\textbf{quiz\_pagine\_421\_430\_007}: E. Sul tratto $(1,6)$ la funzione continua tende a $+\infty$ presso entrambi gli estremi; quindi raggiunge un minimo interno.\par\medskip
\noindent\textbf{quiz\_pagine\_421\_430\_017}: C. Il teorema di Lagrange d\textbackslash{}`a $f\prime(c)=(7-3)/(5-1)=1$.\par\medskip
\noindent\textbf{quiz\_pagine\_421\_430\_018}: E. Dallo sviluppo: $f(2)=e$, $f\prime(2)=e$, $f\prime\prime(2)=-4e$. L\textbackslash{}'affermazione di concavit\textbackslash{}`a in un intorno richiede regolarit\textbackslash{}`a aggiuntiva (per esempio continuit\textbackslash{}`a di $f\prime\prime$).\par\medskip
\noindent\textbf{quiz\_pagine\_431\_440\_008}: C. $f\prime(x)=8(x\log x)^7(\log x+1)$, negativa per $e^{-1}<x<1$. La risposta ufficiale C \textbackslash{}`e un intervallo chiuso, bench\textbackslash{}'e il testo chieda un intervallo aperto: l\textbackslash{}'intervallo aperto corretto sarebbe $(e^{-1},1)$.\par\medskip
\noindent\textbf{quiz\_pagine\_431\_440\_016}: D. Occorrono continuit\textbackslash{}`a su $[a,e]$ (quindi $a>0$) e segni opposti agli estremi: $\log a<0<\log e=1$, cio\textbackslash{}`e $0<a<1$.\par\medskip
\noindent\textbf{quiz\_pagine\_431\_440\_020}: D. Su $[0,1]$ l\textbackslash{}'immagine \textbackslash{}`e $[0,1]$ e su $[1,8]$ \textbackslash{}`e $[0,e^8-e]$; quest\textbackslash{}'ultimo contiene $[0,1]$, quindi le immagini coincidono.\par\medskip
\noindent\textbf{quiz\_pagine\_431\_440\_021}: A. Il teorema dei valori intermedi garantisce ogni valore strettamente compreso fra $1$ e $10$, in particolare $2$.\par\medskip
\noindent\textbf{quiz\_pagine\_441\_450\_001}: D. Per monotonia $f(-1)\le f(x)\le f(2)$; la continuit\textbackslash{}`a non \textbackslash{}`e garantita, quindi non necessariamente si assumono tutti i valori intermedi.\par\medskip
\noindent\textbf{quiz\_pagine\_441\_450\_014}: E. $f\prime(0)=0$, $f\prime\prime(0)=-2<0$: massimo locale stretto. Il coefficiente di $x^4$ \textbackslash{}`e nullo.\par\medskip
\noindent\textbf{quiz\_pagine\_441\_450\_024}: D. Per $x\ne0$, $e^{|x|}>1$ con infimo $1$ non raggiunto; il minimo esiste esattamente se $k\le1$.\par\medskip
\noindent\textbf{quiz\_pagine\_451\_460\_003}: A. La funzione vale $-3x-1$ per $x<-1$, $1-x$ per $-1\le x\le0$ e $3x+1$ per $x>0$. \textbackslash{}`E strettamente decrescente fino a $0$ e poi crescente: sul dominio $(-2,a)$ \textbackslash{}`e iniettiva esattamente per $a\le0$.\par\medskip
\noindent\textbf{quiz\_pagine\_451\_460\_008}: C. Una funzione strettamente crescente \textbackslash{}`e iniettiva, quindi invertibile sulla propria immagine $f(A)$. Non deve essere continua o suriettiva su $\mathbb R$, n\textbackslash{}`e avere limite finito agli estremi.\par\medskip
\noindent\textbf{quiz\_pagine\_451\_460\_009}: B. Serve $x\ne0$ e $f(x)\ne0$; quest\textbackslash{}textquotesingle ultima condizione equivale a $|x|\ne1$. Quindi $D=\mathbb R\setminus\{-1,0,1\}$.\par\medskip
\noindent\textbf{quiz\_pagine\_451\_460\_014}: B. La funzione \textbackslash{}`e strettamente crescente su $(3,4)$, quindi invertibile l\textbackslash{}`i. ATTENZIONE: la risposta B nell\textbackslash{}textquotesingle originale riporta $[3,+\infty)$, che esce dal dominio $x<4$ e non \textbackslash{}`e letteralmente valida. La chiave ufficiale indica B, intendendo evidentemente $[3,4)$.\par\medskip
\noindent\textbf{quiz\_pagine\_451\_460\_025}: E. Una funzione costante \textbackslash{}`e continua e decrescente (non strettamente): la sua immagine \textbackslash{}`e un singoletto, che confuta A, B e C. La formula D ha inoltre gli estremi invertiti per funzioni decrescenti.\par\medskip
\noindent\textbf{quiz\_pagine\_461\_470\_005}: B. Per $a=0$ la funzione \textbackslash{}`e continua e l\textbackslash{}textquotesingle immagine \textbackslash{}`e $[-1,1]$. Per $a\neq0$ manca il valore $0$, interno all\textbackslash{}textquotesingle intervallo tra minimo e massimo.\par\medskip
\noindent\textbf{quiz\_pagine\_461\_470\_006}: A. Per $x\in[-1,0)$, $-\sqrt[3]{x}>0$ e tende a $0$ senza raggiungerlo. Il minimo esiste solo se $f(0)=a\le0$.\par\medskip
\noindent\textbf{quiz\_pagine\_461\_470\_007}: D. $\cosh x$ \textbackslash{}`e decrescente per $x<0$ e ha minimo in $x=0$. Tale punto appartiene a $[-2,a)$ esattamente se $a>0$.\par\medskip
\noindent\textbf{quiz\_pagine\_461\_470\_015}: C. Una funzione affine $ax+b$ ha il medesimo grafico dell\textbackslash{}textquotesingle inversa solo in casi particolari, come $a=1,b=0$ oppure $a=-1$. La proposta C include affini generiche che non hanno questa propriet\textbackslash{}`a.\par\medskip
\noindent\textbf{quiz\_pagine\_471\_480\_003}: E. Il teorema di Lagrange per $g=f^2$ d\textbackslash{}`a $g'(c)=[g(11)-g(0)]/11=(4-0)/11$.\par\medskip
\noindent\textbf{quiz\_pagine\_471\_480\_006}: B. Il valore minimo $0$ viene raggiunto solo se $0\in(-1,a)$, cio\textbackslash{}`e $a>0$.\par\medskip
\noindent\textbf{quiz\_pagine\_471\_480\_017}: D. $f$ \textbackslash{}`e crescente, $f(2)=24$, quindi $32$ non appartiene all'immagine. La E \textbackslash{}`e vera come identit\textbackslash{}`a $f'(0)=1$, indipendentemente da $c$.\par\medskip
\noindent\textbf{quiz\_pagine\_471\_480\_020}: B. La funzione si annulla in $0,1,2,3$; per Rolle la derivata ha almeno uno zero in ciascuno degli intervalli $(0,1)$, $(1,2)$ e $(2,3)$.\par\medskip
\noindent\textbf{quiz\_pagine\_471\_480\_021}: A. $g(1)=\sin(2\pi)=0=g(2)=\sin\pi$, dunque Rolle garantisce $g'(c)=0$. ATTENZIONE: anche C \textbackslash{}`e vera per continuit\textbackslash{}`a di $f$ e passaggio per $3\pi/2$? In realt\textbackslash{}`a $\sin(3\pi/2)=-1$, non zero; per f tra $\pi$ e $2\pi$ non ci sono zeri interni obbligati.\par\medskip
\noindent\textbf{quiz\_pagine\_481\_490\_004}: A. $f(0)=f(\pi)=f(2\pi)=0$. Per Rolle, esiste un punto stazionario in $(0,\pi)$ e uno in $(\pi,2\pi)$.\par\medskip
\noindent\textbf{quiz\_pagine\_481\_490\_011}: A. Per il teorema degli zeri, $f$ si annulla in un punto di $(-1,1)$. Quindi $f^2\ge0$ raggiunge il minimo assoluto $0$.\par\medskip
\noindent\textbf{quiz\_pagine\_481\_490\_022}: E. $f(0)=f(\pi)=0$ e $f$ \textbackslash{}`e continua su $[0,\pi]$ e derivabile all\textbackslash{}'interno. Per Rolle esiste $c\in(0,\pi)$ con $f\prime(c)=0$.\par\medskip
\noindent\textbf{quiz\_pagine\_481\_490\_025}: E. La definizione di massimo locale \textbackslash{}`e esattamente l\textbackslash{}'alternativa E. Anche C \textbackslash{}`e vera, perch\textbackslash{}`e restringe la stessa propriet\textbackslash{}`a ai soli $x\ne x_0$: il testo presenta due alternative corrette.\par\medskip
\noindent\textbf{quiz\_pagine\_511\_520\_010}: C. La funzione $|x|^{2/3}$ \textbackslash{}`e continua sul compatto, ha minimo assoluto $0$ in $x=0$ e massimo assoluto in $x=4$.\par\medskip
\noindent\textbf{quiz\_pagine\_511\_520\_013}: A. Il test della derivata seconda \textbackslash{}`e sufficiente per un massimo locale stretto; non \textbackslash{}`e necessario, come mostra $f(x)=-x^4$ in $0$.\par\medskip
\noindent\textbf{quiz\_pagine\_511\_520\_033}: C. La continuit\textbackslash{}`a in $0$ permette di raccordare i due tratti decrescenti senza salti: per $x<0<y$, $f(x)\ge f(0)\ge f(y)$.\par\medskip
\subsubsection*{4.4 -- Lo studio di funzione}
\noindent\textbf{quiz\_pagine\_001\_010\_007}: C. Spiegazione non ancora disponibile.\par\medskip
\noindent\textbf{quiz\_pagine\_001\_010\_023}: D. Spiegazione non ancora disponibile.\par\medskip
\noindent\textbf{quiz\_pagine\_011\_020\_006}: B. $\log(3+2e^x)=x+\log(2+3e^{-x})=x+\log2+o(1)$.\par\medskip
\noindent\textbf{quiz\_pagine\_041\_050\_011}: B. La formula centrale è stata ricostruita dal testo estratto: verificare visivamente sul PDF.\par\medskip
\noindent\textbf{quiz\_pagine\_061\_070\_006}: B. f(x)-x=log x diverge, quindi non esiste intercetta finita per un asintoto obliquo.\par\medskip
\noindent\textbf{quiz\_pagine\_071\_080\_006}: B. Per $x\to+\infty$, $\log(e^{3x}+4)=3x+o(1)$ e $(\cos x-1)/x=o(1)$; l'asintoto è $y=7x$.\par\medskip
\noindent\textbf{quiz\_pagine\_091\_100\_006}: non verificata / nessuna alternativa corretta. La tabella ufficiale indica E. La funzione è non negativa in $(0,5)$ e $x=2$ è uno zero di ordine pari, dunque è un minimo locale, non un massimo: anche A risulta falsa. Quesito originale problematico, con almeno due risposte NON soddisfatte.\par\medskip
\noindent\textbf{quiz\_pagine\_091\_100\_024}: B. Per $x\to-\infty$, $e^x\to0$ mentre $\cos x$ oscilla fra $-1$ e $1$. Su infiniti intervalli negativi cambia segno, quindi la funzione ha infiniti zeri.\par\medskip
\noindent\textbf{quiz\_pagine\_181\_190\_023}: non verificata / nessuna alternativa corretta. Per $x\to-\infty$, la frazione e' $x+1+o(1)$, mentre $\sqrt{x^2+6x}=-x-3+o(1)$; dunque $f(x)\to-2$, non $1$. La risposta ufficiale D appare incompatibile con la formula stampata: nessuna alternativa e' corretta.\par\medskip
\noindent\textbf{quiz\_pagine\_211\_220\_037}: E. Per $x\to-\infty$, $\log(3+4/x^2)\to\log3$.\par\medskip
\noindent\textbf{quiz\_pagine\_251\_260\_017}: B. $f(x)=ax+b+o(1)$, $g(x)=cx+d+o(1)$, quindi $h(x)=(a-c)x+(b-d)+o(1)$: asintoto obliquo esattamente quando $a\ne c$.\par\medskip
\noindent\textbf{quiz\_pagine\_291\_300\_007}: C. Dividendo $x^2$ per $3x-2$ si ottiene $x/3+2/9$ e un resto trascurabile; anche $2/x$ e' trascurabile.\par\medskip
\noindent\textbf{quiz\_pagine\_301\_310\_013}: B. Dal polinomio di Taylor centrato in zero si ottengono $f(0)=-1$, $f\prime(0)=0$, $f\prime\prime(0)=4$ e $f\prime\prime\prime(0)=-18$.\par\medskip
\noindent\textbf{quiz\_pagine\_301\_310\_018}: A. $\cos(\sin x)=1-\frac12\sin^2x+o(x^2)=1-\frac12x^2+o(x^2)$, dunque i termini fino al secondo ordine si cancellano.\par\medskip
\noindent\textbf{quiz\_pagine\_321\_330\_039}: B. Per $x\to+\infty$, $\log(e^x+x)=x+\log(1+xe^{-x})=x+o(1)$. Il dominio non \textbackslash{}`e tutto $\mathbb R$.\par\medskip
\noindent\textbf{quiz\_pagine\_331\_340\_019}: D. Per definizione di asintoto $f(x)-(2x+6)\to0$ e $g(x)-(4x-3)\to0$; i rapporti con le rispettive rette tendono a 1.\par\medskip
\noindent\textbf{quiz\_pagine\_361\_370\_026}: B. $f(x)/x\to1$, ma $f(x)-x=\log x\to+\infty$, quindi manca l\textbackslash{}textquotesingle intercetta finita di un asintoto obliquo.\par\medskip
\noindent\textbf{quiz\_pagine\_371\_380\_021}: B. Per $x\to+\infty$, $\log(e^{3x}+4)=3x+o(1)$ e $(\cos x-1)/x=o(1)$: $f(x)=7x+o(1)$.\par\medskip
\noindent\textbf{quiz\_pagine\_381\_390\_007}: D. Dallo sviluppo $f(e)=e$ e $f\prime(e)=2$. La tangente e $y-e=2(x-e)$, cioe $y=2x-e$.\par\medskip
\noindent\textbf{quiz\_pagine\_381\_390\_008}: B. Il prodotto $(1-x^2/2+O(x^4))(1+x^2/2+O(x^4))$ ha coefficiente quadratico nullo, quindi e $1+o(x^2)$.\par\medskip
\noindent\textbf{quiz\_pagine\_441\_450\_010}: A. $\cos(1/x)\to1$ ed $e^{-x}\to0$, dunque $f(x)-(2x-1)\to0$.\par\medskip
\noindent\textbf{quiz\_pagine\_441\_450\_026}: E. $f(x)=(2+\sin x)/x^2\,(1+o(1))$, quindi $f\to0$ e ha asintoto orizzontale $y=0$, ma non \textbackslash{}`e equivalente a $2/x^2$.\par\medskip
\noindent\textbf{quiz\_pagine\_501\_510\_006}: E. Per $x\to-\infty$, $\log(3+4/x^2)\to\log3$; dunque $f(x)-(2x+\log3)\to0$.\par\medskip
\subsection*{Capitolo 5}
\subsubsection*{5.1 -- L’integrale di Riemann}
\noindent\textbf{quiz\_pagine\_121\_130\_008}: B. Con $u=3e^x$, $du=3e^x\,dx$, percio\textbackslash{}'o $\int e^x\sin(3e^x)\,dx=-\frac13\cos(3e^x)+C$.\par\medskip
\noindent\textbf{quiz\_pagine\_161\_170\_016}: D. Con $u=2+x^2$, $du=2x\,dx$, si ha $\int x\sqrt{2+x^2}\,dx=\frac13(2+x^2)^{3/2}+C$. Nel PDF le alternative omettono la costante.\par\medskip
\noindent\textbf{quiz\_pagine\_271\_280\_018}: D. Con $u=\sqrt{x+2}$ si ha $dx=2u\,du$, dunque $F=2e^u+C$. Da $F(-1)=0$ segue $C=-2e$.\par\medskip
\noindent\textbf{quiz\_pagine\_321\_330\_028}: B. La primitiva \textbackslash{}`e $\frac23(x+3)^{3/2}$: tra $-3$ e $0$ vale $2\sqrt3$.\par\medskip
\noindent\textbf{quiz\_pagine\_321\_330\_034}: A. $[\log(x+3)]_0^6=\log9-\log3=\log3$.\par\medskip
\noindent\textbf{quiz\_pagine\_341\_350\_012}: B. L\textbackslash{}textquotesingle integrabilita di Riemann non implica derivabilita: $|x|$ su un intervallo contenente zero e integrabile ma non derivabile in zero.\par\medskip
\noindent\textbf{quiz\_pagine\_351\_360\_012}: B. L\textbackslash{}textquotesingle integrabilita di Riemann non implica derivabilita: $|x|$ su un intervallo contenente zero e integrabile ma non derivabile in zero.\par\medskip
\noindent\textbf{quiz\_pagine\_421\_430\_003}: E. Ponendo $t=\sin x$ si ottiene $\int2t^2\,dt=\frac23t^3+C$; poich\textbackslash{}'e $\sin^2x+\cos^2x=1$, E \textbackslash{}`e una primitiva.\par\medskip
\noindent\textbf{quiz\_pagine\_421\_430\_005}: A. Una primitiva \textbackslash{}`e $e^{\sin x}$: agli estremi vale $1-1=0$.\par\medskip
\noindent\textbf{quiz\_pagine\_421\_430\_019}: A. Per il teorema di Lagrange $(f-g)\prime=0$ su un intervallo implica $f-g=C$.\par\medskip
\noindent\textbf{quiz\_pagine\_421\_430\_024}: D. La sola limitatezza non basta: la funzione di Dirichlet, $1$ sui razionali e $0$ sugli irrazionali, non \textbackslash{}`e Riemann-integrabile.\par\medskip
\noindent\textbf{quiz\_pagine\_421\_430\_025}: C. La continuit\textbackslash{}`a \textbackslash{}`e sufficiente ma non necessaria: le funzioni con un numero finito di discontinuit\textbackslash{}`a limitate sono integrabili.\par\medskip
\noindent\textbf{quiz\_pagine\_421\_430\_028}: C. La primitiva \textbackslash{}`e $2\sqrt{e^x-1}$; il valore \textbackslash{}`e $2(3-1)=4$.\par\medskip
\noindent\textbf{quiz\_pagine\_431\_440\_001}: D. $\sqrt{1+\cos x}=\sqrt2\cos(x/2)$ su $[\pi/2,\pi]$. L\textbackslash{}'integrale vale $4\sqrt2[\sin(x/2)]_{\pi/2}^{\pi}=4(\sqrt2-1)$.\par\medskip
\noindent\textbf{quiz\_pagine\_431\_440\_010}: C. L\textbackslash{}'integranda \textbackslash{}`e dispari su un intervallo simmetrico, dunque l\textbackslash{}'integrale \textbackslash{}`e zero.\par\medskip
\noindent\textbf{quiz\_pagine\_431\_440\_011}: D. Con la sostituzione $t=x-2\pi$ e la periodicit\textbackslash{}`a $f(t+2\pi)=f(t)$, gli integrali in D coincidono. Nell\textbackslash{}'originale il limite inferiore di B \textbackslash{}`e tipograficamente corrotto.\par\medskip
\noindent\textbf{quiz\_pagine\_431\_440\_014}: C. L\textbackslash{}'integrale di una funzione generica pu\textbackslash{}`o essere negativo; C non vale senza $f\geq0$ e $a<b$.\par\medskip
\noindent\textbf{quiz\_pagine\_431\_440\_015}: D. Le intersezioni sono $x=0$ e $x=3/2$. L\textbackslash{}'area \textbackslash{}`e $\int_0^{3/2}(3x/2-x^2)\,dx=9/16$.\par\medskip
\noindent\textbf{quiz\_pagine\_431\_440\_023}: B. Ponendo $t=3e^x$ si ha $dt=3e^x\,dx$ e $\int e^x\sin(3e^x)\,dx=-\frac13\cos(3e^x)+C$.\par\medskip
\noindent\textbf{quiz\_pagine\_481\_490\_016}: A. Essendo $F\prime=f$ e $F$ derivabile, per Fermat $F\prime(x_0)=f(x_0)=0$.\par\medskip
\subsubsection*{5.2 -- Proprietà dell’integrale}
\noindent\textbf{quiz\_pagine\_011\_020\_018}: A. $F\prime(x)=x^4\cos^4x\ge0$, quindi $F$ \textbackslash{}`e non decrescente su $\mathbb R$.\par\medskip
\noindent\textbf{quiz\_pagine\_031\_040\_007}: B. Spiegazione non ancora disponibile.\par\medskip
\noindent\textbf{quiz\_pagine\_041\_050\_010}: A. Risposta ufficiale indicata nella tabella della decima simulazione.\par\medskip
\noindent\textbf{quiz\_pagine\_041\_050\_038}: non verificata / nessuna alternativa corretta. Soluzione non riportata nel PDF; da verificare.\par\medskip
\noindent\textbf{quiz\_pagine\_041\_050\_040}: non verificata / nessuna alternativa corretta. Soluzione non riportata nel PDF; da verificare.\par\medskip
\noindent\textbf{quiz\_pagine\_051\_060\_002}: E. Una potenza dispari di una funzione dispari resta dispari e ha integrale nullo su un intervallo simmetrico.\par\medskip
\noindent\textbf{quiz\_pagine\_051\_060\_003}: D. Il teorema fondamentale del calcolo richiede la continuità di $f$ per garantire $F\prime=f$ nell'interno.\par\medskip
\noindent\textbf{quiz\_pagine\_051\_060\_009}: C. Per $F(x)=\int_x^a f(t)dt$ vale $F\prime(x)=-f(x)\le0$; inoltre $F(a)=0$.\par\medskip
\noindent\textbf{quiz\_pagine\_051\_060\_019}: C. Poiché $F\prime(x)=-f(x)\le0$, la funzione integrale è non crescente.\par\medskip
\noindent\textbf{quiz\_pagine\_061\_070\_008}: E. DISCREPANZA: risposta ufficiale E, ma integrando una funzione pari le opzioni A e D sono vere; E è falsa in generale (esempio f=1).\par\medskip
\noindent\textbf{quiz\_pagine\_061\_070\_011}: B. F è strettamente crescente e F(0)=0, perciò F(-1)<0.\par\medskip
\noindent\textbf{quiz\_pagine\_061\_070\_018}: A. Integrando le disuguaglianze strette 2<f<4 e dividendo per b-a>0 si ottiene 2<m<4.\par\medskip
\noindent\textbf{quiz\_pagine\_071\_080\_005}: A. Per $x\ge x_0$ l'integrale di una funzione non negativa è non negativo: $F(x)=\int_{x_0}^x f(t)\,dt\ge0$.\par\medskip
\noindent\textbf{quiz\_pagine\_091\_100\_021}: A. $F(x)=\log\cosh x\geq0$ per ogni $x$, e $F(x)>0$ anche per $x<0$; dunque A è falsa. $F\prime=\tanh x$ e $F^{\prime\prime}=1/\cosh^2x>0$.\par\medskip
\noindent\textbf{quiz\_pagine\_101\_110\_008}: D. Per il teorema fondamentale del calcolo $F\prime(x)=e^{-x^2}>0$: $F$ è strettamente crescente. La funzione integranda è pari, perciò $F$ è dispari.\par\medskip
\noindent\textbf{quiz\_pagine\_141\_150\_012}: C. Il teorema della media integrale vale per funzioni continue: $f(c)=\frac1{b-a}\int_a^b f$. Per $[4,5]$ il fattore e\textbackslash{}'e $1$.\par\medskip
\noindent\textbf{quiz\_pagine\_141\_150\_013}: E. Per il teorema fondamentale $F\prime(x)=-x^2/\cosh(x^2)\leq0$, con uguaglianza soltanto in $0$: $F$ e\textbackslash{}'e strettamente decrescente.\par\medskip
\noindent\textbf{quiz\_pagine\_151\_160\_002}: C. La media nulla implica $\int_0^4 f(t)\,dt=0$; per Newton--Leibniz $F(4)-F(0)=0$.\par\medskip
\noindent\textbf{quiz\_pagine\_151\_160\_003}: E. $F\prime=f\leq0$, dunque $F$ e\textbackslash{}'e non crescente. Attenzione: anche la (b) e\textbackslash{}'e vera, perche\textbackslash{}'e $F(0)=0$; il quesito ha due risposte valide. La media in (d) dovrebbe avere denominatore $2$.\par\medskip
\noindent\textbf{quiz\_pagine\_151\_160\_008}: C. La derivata di $\log|f(x)|$ e\textbackslash{}'e $f\prime(x)/f(x)$ nei punti in cui $f(x)\ne0$.\par\medskip
\noindent\textbf{quiz\_pagine\_151\_160\_012}: C. Per $x\in[0,4]$, $1\leq e^{\sqrt{x}}\leq e^2<8$; integrando si ha $4<I<4e^2<32$.\par\medskip
\noindent\textbf{quiz\_pagine\_161\_170\_023}: C. $f$ e\textbackslash{}'e dispari, quindi l'integrale su un intervallo simmetrico vale zero.\par\medskip
\noindent\textbf{quiz\_pagine\_171\_180\_025}: A. La media integrale puo' coincidere con un estremo se $f$ e' costante: (a), (c) ed (e) non sono necessariamente vere senza continuita' e stretta crescita. Il testo e' ambiguo: (a) ed (e) sono equivalenti.\par\medskip
\noindent\textbf{quiz\_pagine\_181\_190\_010}: E. La derivata del denominatore e' $2x+1$; la primitiva e' $\log(x^2+x+5)$, percio' il risultato e' $\log(11/7)$.\par\medskip
\noindent\textbf{quiz\_pagine\_191\_200\_006}: A. Con $u=2+x^2$, $du=2x\,dx$, si ottiene $\frac13(2+x^2)^{3/2}+C$. La A coincide con una primitiva valida: $\frac13(2+x^2)^{3/2}-2$. La costante additiva puo' essere qualsiasi numero reale.\par\medskip
\noindent\textbf{quiz\_pagine\_211\_220\_001}: B. $F'(x)=e^{-3x^2}\cos x\geq0$ su $[0,\pi/2]$ e $F(0)=0$.\par\medskip
\noindent\textbf{quiz\_pagine\_211\_220\_002}: A. $F(x)=x+o(x)$, dunque $F(x)/x^3\sim1/x^2\to+\infty$.\par\medskip
\noindent\textbf{quiz\_pagine\_211\_220\_003}: B. $F(x)=2x+o(x)$ per il teorema fondamentale del calcolo; $2xe^x\sim2x$.\par\medskip
\noindent\textbf{quiz\_pagine\_211\_220\_005}: D. $F'(x)=\arctan(x^5)+\sqrt[3]x$, quindi $F'(0)=0$.\par\medskip
\noindent\textbf{quiz\_pagine\_211\_220\_018}: B. $F(\pi)=\int_0^\pi e^{-3t^2}\sin t\,dt>0$.\par\medskip
\noindent\textbf{quiz\_pagine\_211\_220\_019}: B. $F'(x)=e^{-3x^2}\sin x$ cambia segno.\par\medskip
\noindent\textbf{quiz\_pagine\_211\_220\_020}: B. L'integranda $e^{-3t^2}\sin t$ e' dispari, quindi la funzione integrale da $0$ a $x$ e' pari.\par\medskip
\noindent\textbf{quiz\_pagine\_211\_220\_021}: A. Per $x>0$, l'integrale e' positivo: i contributi positivi prevalgono grazie al peso gaussiano decrescente.\par\medskip
\noindent\textbf{quiz\_pagine\_211\_220\_022}: B. La derivata cambia segno, dunque non e' iniettiva.\par\medskip
\noindent\textbf{quiz\_pagine\_211\_220\_023}: B. La derivata $e^{-3x^2}\sin x$ cambia segno.\par\medskip
\noindent\textbf{quiz\_pagine\_211\_220\_024}: B. $F(x)\sim x^2/2$, quindi il rapporto tende a $1/2$.\par\medskip
\noindent\textbf{quiz\_pagine\_211\_220\_025}: A. L'integrale $\int_0^{+\infty}e^{-3t^2}\sin t\,dt$ converge assolutamente.\par\medskip
\noindent\textbf{quiz\_pagine\_211\_220\_028}: A. $F(x)\geq0$ per ogni $x$ e $F(0)=0$.\par\medskip
\noindent\textbf{quiz\_pagine\_211\_220\_029}: B. $F''(0)=1>0$, dunque in zero non c'e' flesso.\par\medskip
\noindent\textbf{quiz\_pagine\_211\_220\_030}: B. $F(x)=x^2/2+O(x^4)$, quindi non ha termini lineari o cubici.\par\medskip
\noindent\textbf{quiz\_pagine\_211\_220\_031}: A. $e^{-3t^2}\sin t=t+O(t^3)$, dunque $F(x)=x^2/2+O(x^4)$.\par\medskip
\noindent\textbf{quiz\_pagine\_221\_230\_001}: C. Per $a<b$, integrando $\inf f\leq f(x)\leq\sup f$ e dividendo per $b-a$ si ottiene la disuguaglianza; non occorre continuita'.\par\medskip
\noindent\textbf{quiz\_pagine\_221\_230\_002}: C. Per il teorema fondamentale $F'(x)=-f(x)\leq0$: $F$ e' non crescente su tutta la retta.\par\medskip
\noindent\textbf{quiz\_pagine\_231\_240\_012}: B. $\int_0^5 x^{1/3}dx=[\frac34x^{4/3}]_0^5=\frac{15}{4}\sqrt[3]{5}$.\par\medskip
\noindent\textbf{quiz\_pagine\_241\_250\_006}: B. $\int_{-2}^0\frac{2x}{x^2+3}dx=\ln3-\ln7$; $\int_0^2(x^2+3x)dx=\frac{26}{3}$.\par\medskip
\noindent\textbf{quiz\_pagine\_261\_270\_006}: A. $\sinh x=x+x^3/6+o(x^3)$; sviluppando l'esponenziale il coefficiente cubico e' $1/6+1/6=1/3$.\par\medskip
\noindent\textbf{quiz\_pagine\_261\_270\_007}: E. $\log(1+x)=x-x^2/2+x^3/3+o(x^3)$ e $\sinh t=t+t^3/6+o(t^3)$.\par\medskip
\noindent\textbf{quiz\_pagine\_261\_270\_011}: C. Lo sviluppo e' centrato in $7$ (non in $0$, nonostante il termine 'Mac Laurin'); $f(7)=-3$ e $f'(7)=2$.\par\medskip
\noindent\textbf{quiz\_pagine\_261\_270\_026}: B. $\sqrt[5]{x+x^2}-\sqrt[5]{x}=\frac15x^{6/5}+o(x^{6/5})$; domina $x^{5/6}$ per $x\to0^+$.\par\medskip
\noindent\textbf{quiz\_pagine\_261\_270\_027}: E. Per $x\to+\infty$, $\sqrt[5]{x+x^2}\sim x^{2/5}$, $x^{5/6}$ domina; nell'originale la chiave indica E, ma l'alternativa corretta dovrebbe essere una forma $x^{5/6}$, assente. Attenzione: la fonte usa radice quinta di $x+x^2$.\par\medskip
\noindent\textbf{quiz\_pagine\_271\_280\_019}: A. L'integranda e' $1/(\sqrt x(1+x))$: e' confrontabile con $x^{-1/2}$ vicino a zero e con $x^{-3/2}$ all'infinito, quindi converge.\par\medskip
\noindent\textbf{quiz\_pagine\_281\_290\_013}: E. Se $f$ fosse costante, l'integrale varrebbe esattamente $(b-a)f(a)$, in contraddizione con l'ipotesi.\par\medskip
\noindent\textbf{quiz\_pagine\_281\_290\_014}: C. $2\min_{[1,3]}f\leq\int_1^3f\leq\int_1^2g\leq\max_{[1,2]}g$.\par\medskip
\noindent\textbf{quiz\_pagine\_281\_290\_015}: C. Per $0\leq x\leq1$, integrando $g(t)\leq t$ da 0 a $x$ si ottiene $\int_0^x g(t)dt\leq x^2/2$. La fonte non esplicita il dominio di $x$ nell'alternativa.\par\medskip
\noindent\textbf{quiz\_pagine\_281\_290\_016}: A. Per periodicita', $f(x+2\pi)=f(x)$ e il cambio di variabile $u=x-2\pi$ dimostra l'uguaglianza degli integrali.\par\medskip
\noindent\textbf{quiz\_pagine\_311\_320\_012}: C. Il valore assoluto del coseno ha quattro quarti d\textbackslash{}'onda con area unitaria: integrale uguale a 4.\par\medskip
\noindent\textbf{quiz\_pagine\_311\_320\_013}: A. $F\prime(x)=x^4\cos^4x\geq0$ per ogni reale, quindi $F$ e crescente (in senso non decrescente) su tutto $\mathbb R$.\par\medskip
\noindent\textbf{quiz\_pagine\_341\_350\_002}: D. Su $[-1,0]$ si ha $f(x)\leq x$ e dunque $\int_{-1}^0f(x)\,dx\leq\int_{-1}^0x\,dx=-1/2<0$.\par\medskip
\noindent\textbf{quiz\_pagine\_341\_350\_023}: D. Per $x\to+\infty$ l\textbackslash{}textquotesingle integranda e asintotica a $x^{-1/2}>0$, il cui integrale diverge. Vicino a zero l\textbackslash{}textquotesingle integranda e continua.\par\medskip
\noindent\textbf{quiz\_pagine\_351\_360\_002}: D. Su $[-1,0]$ si ha $f(x)\leq x$ e dunque $\int_{-1}^0f(x)\,dx\leq\int_{-1}^0x\,dx=-1/2<0$.\par\medskip
\noindent\textbf{quiz\_pagine\_351\_360\_023}: D. Per $x\to+\infty$ l\textbackslash{}textquotesingle integranda e asintotica a $x^{-1/2}>0$, il cui integrale diverge. Vicino a zero l\textbackslash{}textquotesingle integranda e continua.\par\medskip
\noindent\textbf{quiz\_pagine\_361\_370\_002}: E. Per $n$ dispari $f^n$ e dispari: l\textbackslash{}textquotesingle integrale su un intervallo simmetrico e nullo. Per $n$ pari $f^n$ e pari.\par\medskip
\noindent\textbf{quiz\_pagine\_361\_370\_003}: D. Il teorema fondamentale del calcolo richiede la continuita di $f$ per garantire $F\prime=f$ nei punti interni; agli estremi si intendono le derivate unilaterali.\par\medskip
\noindent\textbf{quiz\_pagine\_361\_370\_009}: C. $F\prime(x)=-f(x)\leq0$, quindi $F$ e non crescente; $F(a)=0$, $F(x)\leq0$ per $x\geq a$ e $F(x)\geq0$ per $x<a$. L\textbackslash{}textquotesingle alternativa ufficiale C usa ``positiva\textbackslash{}textquotesingle\textbackslash{}textquotesingle in senso stretto, non garantito se $f$ si annulla.\par\medskip
\noindent\textbf{quiz\_pagine\_361\_370\_019}: C. $F\prime(x)=-f(x)\leq0$, dunque $F$ e non crescente su tutto $\mathbb R$.\par\medskip
\noindent\textbf{quiz\_pagine\_361\_370\_028}: A. $f(x^2)$ e pari, dunque l\textbackslash{}textquotesingle integrale simmetrico vale il doppio dell\textbackslash{}textquotesingle integrale su $[0,1]$. Anche D e matematicamente equivalente ad A: la domanda ha due risposte corrette.\par\medskip
\noindent\textbf{quiz\_pagine\_371\_380\_001}: B. $F(0)=0$ e $F\prime=f>0$, dunque $F$ e strettamente crescente e $F(-1)<0$.\par\medskip
\noindent\textbf{quiz\_pagine\_371\_380\_008}: A. Integrando $2<f(x)<4$ su $[a,b]$ e dividendo per $b-a>0$ si ricava $2<m<4$.\par\medskip
\noindent\textbf{quiz\_pagine\_371\_380\_020}: A. Per $x\geq x_0$, l\textbackslash{}textquotesingle integrale di una funzione non negativa su $[x_0,x]$ e non negativo.\par\medskip
\noindent\textbf{quiz\_pagine\_401\_410\_013}: A. $F(x)=\log\cosh x\geq0$ per ogni $x$; dunque \textbackslash{}`e falsa la pretesa $F(x)<0$ per $x<0$. Inoltre $F\prime=\tanh x$ e $F\prime\prime=1/\cosh^2x>0$.\par\medskip
\noindent\textbf{quiz\_pagine\_411\_420\_008}: D. Per il teorema fondamentale $F\prime(x)=e^{-x^2}>0$, dunque $F$ \textbackslash{}`e strettamente crescente, non decrescente.\par\medskip
\noindent\textbf{quiz\_pagine\_441\_450\_016}: E. Teorema della media integrale per $a<b$.\par\medskip
\noindent\textbf{quiz\_pagine\_481\_490\_008}: C. La funzione \textbackslash{}`e dispari; l\textbackslash{}'integrale su un intervallo simmetrico rispetto all\textbackslash{}'origine vale zero.\par\medskip
\noindent\textbf{quiz\_pagine\_491\_500\_009}: C. La media integrale non deve essere assunta dalla funzione se manca la continuit\textbackslash{}`a. ATTENZIONE: per una funzione costante anche A ed E falliscono, poich\textbackslash{}'e $(f(a),f(b))$ \textbackslash{}`e vuoto; il quesito richiederebbe ipotesi aggiuntive.\par\medskip
\noindent\textbf{quiz\_pagine\_491\_500\_018}: E. Integrando $f(x)\ge x$ su $[0,1]$ si ottiene $\int_0^1f(x)\,dx\ge1/2$, e la lunghezza dell' intervallo \textbackslash{}`e $1$.\par\medskip
\noindent\textbf{quiz\_pagine\_491\_500\_020}: E. $f$ non \textbackslash{}`e definita in $0$; l' area richiede $|f|$ e la media richiede il fattore $1/2$. Inoltre $\mu=(-1/2+2/\pi)/2>0$, mentre $f(x)<0$ su $(-1,0)$.\par\medskip
\noindent\textbf{quiz\_pagine\_501\_510\_014}: C. La definizione di primitiva richiede derivabilit\textbackslash{}`a su tutto $I$ e $F\prime=f$ su $I$. La (a) ne d\textbackslash{}`a una costruzione particolare, non tutte le primitive.\par\medskip
\noindent\textbf{quiz\_pagine\_501\_510\_017}: C. Per una funzione integrabile il valor medio \textbackslash{}`e compreso tra infimo e supremo, anche se il minimo o massimo non esistono. L\textbackslash{}'uguaglianza $\mu=f(c)$ richiede ulteriori ipotesi, ad esempio continuit\textbackslash{}`a.\par\medskip
\noindent\textbf{quiz\_pagine\_501\_510\_018}: C. Per il teorema fondamentale del calcolo $F\prime(x)=-f(x)\le0$, quindi $F$ \textbackslash{}`e non crescente.\par\medskip
\noindent\textbf{quiz\_pagine\_501\_510\_037}: C. Il teorema della media integrale nella forma con $f(c)=\mu$ richiede la continuit\textbackslash{}`a su $[a,b]$.\par\medskip
\noindent\textbf{quiz\_pagine\_501\_510\_038}: B. $F\prime(x)=x^2\cosh(x^2)\ge0$ ed \textbackslash{}`e positiva per $x\ne0$: $F$ \textbackslash{}`e strettamente crescente su $\mathbb R$.\par\medskip
\subsubsection*{5.3 -- Integrali indefiniti e regole di integrazione}
\noindent\textbf{quiz\_pagine\_041\_050\_034}: non verificata / nessuna alternativa corretta. Soluzione non riportata nel PDF; da verificare.\par\medskip
\noindent\textbf{quiz\_pagine\_041\_050\_041}: non verificata / nessuna alternativa corretta. Soluzione non riportata nel PDF; da verificare.\par\medskip
\noindent\textbf{quiz\_pagine\_081\_090\_004}: A. Una primitiva è $-\frac32\cos 2x$, assente fra le alternative B--E.\par\medskip
\noindent\textbf{quiz\_pagine\_081\_090\_009}: D. Una primitiva è $-\cos(f(x))$; agli estremi si ottiene $\cos f(0)-\cos f(2)$.\par\medskip
\noindent\textbf{quiz\_pagine\_111\_120\_020}: E. Una primitiva e' $\frac23\sin^3x+C$; poiche' $\sin^2x+\cos^2x=1$, l'opzione e) e' equivalente con $C=1$.\par\medskip
\noindent\textbf{quiz\_pagine\_121\_130\_009}: B. Per $x\to+\infty$ l\textbackslash{}'integranda e\textbackslash{}'e equivalente a $x^{3/2}/x^2=1/\sqrt{x}$ ed e\textbackslash{}'e positiva definitivamente. L\textbackslash{}'integrale diverge a $+\infty$. In $0$ l\textbackslash{}'integranda e\textbackslash{}'e regolare.\par\medskip
\noindent\textbf{quiz\_pagine\_121\_130\_017}: D. La funzione $F(c)=\int_0^c f(x)\,dx$ e\textbackslash{}'e non decrescente e non negativa: il limite esiste in $[0,+\infty]$. La risposta ufficiale D esclude erroneamente $0$, possibile per $f\equiv0$; quindi nessuna alternativa e\textbackslash{}'e universalmente corretta come scritta.\par\medskip
\noindent\textbf{quiz\_pagine\_161\_170\_003}: E. Poiche\textbackslash{}'e $1-\pi<0$ e $x^{3-\pi}=1/x^{\pi-3}$ con $\pi-3<1$, l'integrale diverge negativamente.\par\medskip
\noindent\textbf{quiz\_pagine\_161\_170\_010}: C. Per parti, $\int xe^x\,dx=(x-1)e^x$; tra $-\infty$ e $0$ vale $-1$.\par\medskip
\noindent\textbf{quiz\_pagine\_181\_190\_001}: E. Dalla disuguaglianza $f(x)\geq x$, integrando su $[0,1]$, segue $\mu=\int_0^1f(x)\,dx\geq\int_0^1x\,dx=1/2$.\par\medskip
\noindent\textbf{quiz\_pagine\_181\_190\_003}: E. $f$ non e' definita in $0$. L'area richiede il valore assoluto, il valor medio richiede il fattore $1/2$ e $\mu=(-1/2+2/\pi)/2>0$, mentre $f(c)=c<0$ per $c\in(-1,0)$.\par\medskip
\noindent\textbf{quiz\_pagine\_181\_190\_027}: D. Quando le derivate esistono e sono integrabili, la derivata di ordine pari di una funzione dispari e' dispari, quindi il suo integrale su un intervallo simmetrico e' nullo. Il testo suppone solo integrabilita' di $f$, insufficiente per garantire che $f^{(n)}$ esista.\par\medskip
\noindent\textbf{quiz\_pagine\_211\_220\_045}: non verificata / nessuna alternativa corretta. Le opzioni A e C sono entrambe corrette (A fornisce una primitiva canonica, C la definizione); anche B implica derivabilita' se l'uguaglianza delle derivate vale in ogni punto. Quesito ambiguo.\par\medskip
\noindent\textbf{quiz\_pagine\_231\_240\_022}: A. Se $F\prime=f$, allora anche $(F+C)\prime=f$ per ogni costante reale $C$.\par\medskip
\noindent\textbf{quiz\_pagine\_241\_250\_022}: E. Pongo $u=x^2-4$, $du=2x\,dx$: $\int3x/(x^2-4)dx=(3/2)\ln|x^2-4|+C$.\par\medskip
\noindent\textbf{quiz\_pagine\_251\_260\_020}: B. Le primitive sono $\log|f(x)|+C$ negli intervalli dove $f\ne0$. I fattori costanti dentro il logaritmo aggiungono costanti; il fattore $2$ davanti al logaritmo ne cambia invece la derivata.\par\medskip
\noindent\textbf{quiz\_pagine\_321\_330\_004}: A. $\frac{x-4}{5-x}=-1-\frac1{x-5}$, dunque $F\prime(x)=1/(x-5)^2$ nell\textbackslash{}textquotesingle intervallo indicato.\par\medskip
\noindent\textbf{quiz\_pagine\_321\_330\_018}: E. Le prime quattro hanno derivata $2x/[3(x^2+1)]$. L\textbackslash{}textquotesingle ultima include il termine $-\log(2x)/3$ e non \textbackslash{}`e definita su tutta $\mathbb R$.\par\medskip
\noindent\textbf{quiz\_pagine\_321\_330\_019}: B. Una primitiva \textbackslash{}`e $\log|f(x)|+C$ negli intervalli dove $f\neq0$. Il fattore $2$ davanti al logaritmo ne raddoppia la derivata.\par\medskip
\noindent\textbf{quiz\_pagine\_321\_330\_021}: D. La fonte richiede che l\textbackslash{}textquotesingle intervallo $I$ sia connesso: $\mathbb R\setminus\{0\}$ non \textbackslash{}`e un intervallo. Tuttavia $\log|x|$ \textbackslash{}`e effettivamente una primitiva su questo dominio disconnesso in senso generale: quesito ambiguo.\par\medskip
\noindent\textbf{quiz\_pagine\_321\_330\_022}: E. Per $x>0$ la differenza $F-G$ vale $\pi/2$; per $x<0$ vale $-\pi/2$. La formula E non vale in tutto il dominio.\par\medskip
\noindent\textbf{quiz\_pagine\_331\_340\_008}: D. $\int\tan x\,dx=-\log|\cos x|+C$; il valore \textbackslash{}`e $\frac12\log2$, non $\log2$.\par\medskip
\noindent\textbf{quiz\_pagine\_341\_350\_025}: A. Per $x\leq x_0$, $F(x)=-\int_x^{x_0}f(t)\,dt\geq0$ dato che $f\leq0$.\par\medskip
\noindent\textbf{quiz\_pagine\_351\_360\_025}: A. Per $x\leq x_0$, $F(x)=-\int_x^{x_0}f(t)\,dt\geq0$ dato che $f\leq0$.\par\medskip
\subsubsection*{5.4 -- Tecniche di integrazione}
\noindent\textbf{quiz\_pagine\_001\_010\_034}: C. Spiegazione non ancora disponibile.\par\medskip
\noindent\textbf{quiz\_pagine\_001\_010\_037}: A. Spiegazione non ancora disponibile.\par\medskip
\noindent\textbf{quiz\_pagine\_011\_020\_017}: C. Per simmetria, $4\int_0^{\pi/2}\cos x\,dx=4$.\par\medskip
\noindent\textbf{quiz\_pagine\_021\_030\_008}: A. $\frac{x-4}{5-x}=-1-\frac1{x-5}$ e la derivata \textbackslash{}`e $1/(x-5)^2$.\par\medskip
\noindent\textbf{quiz\_pagine\_021\_030\_022}: E. La (e) non \textbackslash{}`e definita su tutta $\mathbb R$ e la sua derivata include $-1/(3x)$.\par\medskip
\noindent\textbf{quiz\_pagine\_021\_030\_025}: D. Il testo usa la definizione di primitiva su un intervallo: $\mathbb R\setminus\{0\}$ non \textbackslash{}`e un intervallo, pur essendo la derivata $1/x$ su ciascuna componente.\par\medskip
\noindent\textbf{quiz\_pagine\_041\_050\_035}: non verificata / nessuna alternativa corretta. Soluzione non riportata nel PDF; da verificare.\par\medskip
\noindent\textbf{quiz\_pagine\_061\_070\_005}: E. Integrazione per parti con u=x e dv=g\textbackslash{}prime(x)dx: 2g(2)-integrale di g.\par\medskip
\noindent\textbf{quiz\_pagine\_061\_070\_007}: C. Sostituendo u=3t+2 si ha dt=du/3, con estremi 2 e 8.\par\medskip
\noindent\textbf{quiz\_pagine\_061\_070\_012}: E. Con t=2x, dx=dt/2 e gli estremi diventano 0 e 2.\par\medskip
\noindent\textbf{quiz\_pagine\_061\_070\_017}: A. Per parti con dv=dx/x. DISCREPANZA: A e B sono equivalenti su [1,2]; la tabella ufficiale indica A.\par\medskip
\noindent\textbf{quiz\_pagine\_061\_070\_019}: D. Con t=log x e dt=dx/x gli estremi diventano 0 e log 2. DISCREPANZA: anche B ripete identicamente il membro sinistro.\par\medskip
\noindent\textbf{quiz\_pagine\_111\_120\_022}: A. Una primitiva e' $e^{\sin x}$; agli estremi vale $e^0-e^0=0$.\par\medskip
\noindent\textbf{quiz\_pagine\_131\_140\_018}: E. La risposta A indicata nella tabella ufficiale non e'e definita su alcun intervallo reale: $\sqrt{-x}>1$ richiede $x<-1$, ma $\log(x+1)$ richiede $x>-1$. Per $t=\sqrt{-x}$ si ottiene $-2\int dt/(1-t^2)$, con primitive logaritmiche sui singoli intervalli del dominio. La forma dell\textbackslash{}'opzione A e\textbackslash{}'e definita solo dove gli argomenti dei logaritmi sono positivi: risposta ufficiale da considerare con cautela.\par\medskip
\noindent\textbf{quiz\_pagine\_141\_150\_011}: A. $\frac{x+1}{x-1}=1+\frac2{x-1}$, quindi una primitiva e\textbackslash{}'e $x+2\log|x-1|=x+\log((1-x)^2)$.\par\medskip
\noindent\textbf{quiz\_pagine\_171\_180\_003}: C. Con $t=2x^2+3$ si ha $dt=4x\,dx$; gli estremi diventano $3$ e $5$.\par\medskip
\noindent\textbf{quiz\_pagine\_211\_220\_047}: E. Ponendo $u=2-\log x$, si ottiene $\log[(2-\log2)/(2-2\log2)]$. L'alternativa E e' equivalente per $\log(2a)=\log a+\log2$.\par\medskip
\noindent\textbf{quiz\_pagine\_221\_230\_004}: D. Il prodotto a destra si annulla per $x=0$, $x=2$ e $\sin x=1$; per la radice si richiede $x\geq0$. L'equazione e' separabile. La variabile $t$ non e' una soluzione $x(t)$.\par\medskip
\noindent\textbf{quiz\_pagine\_281\_290\_019}: C. Con $t=\log x$, $dt=dx/x$ e gli estremi diventano 0 e $\log2$.\par\medskip
\noindent\textbf{quiz\_pagine\_281\_290\_021}: D. Il numeratore e' meta' della derivata del denominatore. L'integrale e' $\frac12\log(9/4)=\log(3/2)$, mentre D vale $\log(3/4)$.\par\medskip
\noindent\textbf{quiz\_pagine\_321\_330\_016}: D. Con $u=2-3x^2$, $du=-6x\,dx$, quindi l\textbackslash{}textquotesingle integrale \textbackslash{}`e $-u^6/36+c$.\par\medskip
\noindent\textbf{quiz\_pagine\_331\_340\_017}: B. La monotonia d\textbackslash{}`a $f(a)\leq f(x)\leq f(b)$; integrando su $[a,b]$ si ottiene B.\par\medskip
\noindent\textbf{quiz\_pagine\_331\_340\_023}: C. In ciascuna semionda l\textbackslash{}textquotesingle area vale 2, quindi in un periodo completo $2+2=4$.\par\medskip
\noindent\textbf{quiz\_pagine\_361\_370\_025}: E. Per parti, con $u=x$ e $dv=g\prime(x)dx$, si ottiene $[xg(x)]_0^2-\int_0^2g(x)dx=2g(2)-\int_0^2g(x)dx$.\par\medskip
\noindent\textbf{quiz\_pagine\_361\_370\_027}: C. Con $u=3t+2$ si ha $du=3dt$ e gli estremi diventano $2$ e $8$.\par\medskip
\noindent\textbf{quiz\_pagine\_371\_380\_002}: E. Con $t=2x$, $dx=dt/2$ e gli estremi diventano $0$ e $2$.\par\medskip
\noindent\textbf{quiz\_pagine\_371\_380\_007}: A. Per parti con $u=f(x)$, $dv=dx/x$, si ha $[f(x)\log x]_1^2-\int_1^2 f\prime(x)\log x\,dx$. Le alternative A e B sono equivalenti poiche $x>0$.\par\medskip
\noindent\textbf{quiz\_pagine\_371\_380\_009}: D. Con $t=\log x$ si ha $dt=dx/x$ e gli estremi sono $0$ e $\log2$. La risposta B ripete letteralmente il membro sinistro ed e anch\textbackslash{}textquotesingle essa una uguaglianza vera: il quiz non e univoco.\par\medskip
\noindent\textbf{quiz\_pagine\_381\_390\_009}: A. Una primitiva e $-\frac32\cos2x+c$; nessuna delle alternative B--E coincide.\par\medskip
\noindent\textbf{quiz\_pagine\_381\_390\_014}: D. Per sostituzione $u=f(x)$, una primitiva e $-\cos f(x)$; il valore e $\cos f(0)-\cos f(2)$.\par\medskip
\noindent\textbf{quiz\_pagine\_441\_450\_004}: B. Pongo $u=2+x^4$, $du=4x^3dx$: $\int 2x^3/(u^2+1)dx=\frac12\arctan u+C$.\par\medskip
\noindent\textbf{quiz\_pagine\_471\_480\_027}: D. Ponendo $u=2+x^2$, $du=2x\,dx$, l'integrale vale $\frac12\int u^{1/2}du=\frac13u^{3/2}+C$. Le alternative omettono la costante arbitraria.\par\medskip
\noindent\textbf{quiz\_pagine\_481\_490\_012}: C. Con $t=2x^2+3$, $dt=4x\,dx$; gli estremi diventano $3$ e $5$.\par\medskip
\noindent\textbf{quiz\_pagine\_491\_500\_004}: E. La funzione \textbackslash{}`e positiva e, per $|x|\ge1$, $e^{-x^2}\le1/x^2$. Quindi l' integrale converge a un valore strettamente positivo.\par\medskip
\noindent\textbf{quiz\_pagine\_491\_500\_019}: E. $|\sin x\,e^x|\le e^x$ e $\int_{-\infty}^{-1}e^x\,dx=e^{-1}<\infty$, dunque converge assolutamente e anche semplicemente. L' alternativa B \textbackslash{}`e falsa.\par\medskip
\noindent\textbf{quiz\_pagine\_501\_510\_016}: E. Con $u=2-\log x$, l\textbackslash{}'integrale \textbackslash{}`e $\log[(2-\log2)/(2-\log4)]$. Poich\textbackslash{}'e $2-\log4=2(1-\log2)$, si ottiene la (e).\par\medskip
\noindent\textbf{quiz\_pagine\_501\_510\_036}: A. $f(x)=(3/2)x/(x^2+1)$; integrando con $u=x^2+1$ si ottiene $(3/4)\log(x^2+1)+C$.\par\medskip
\noindent\textbf{quiz\_pagine\_511\_520\_017}: D. Con $t=\sqrt{x}$, $dx=2t\,dt$, dunque $\int_0^2 2t/(1+t)\,dt=[2t-2\log(1+t)]_0^2=4-2\log3$.\par\medskip
\noindent\textbf{quiz\_pagine\_511\_520\_018}: C. Da $x=g(t)$ segue $dx=g'(t)\,dt$, quindi l'integrando trasformato \textbackslash{}`e $f(g(t))g'(t)$.\par\medskip
\noindent\textbf{quiz\_pagine\_511\_520\_038}: D. Con $t=\sqrt{x+2}$, $dx=2t\,dt$, e $\int e^t/t\,2t\,dt=2e^t+C$. Da $F(-1)=0$ segue $C=-2e$.\par\medskip
\subsubsection*{5.5 -- Integrali impropri}
\noindent\textbf{quiz\_pagine\_001\_010\_029}: A. Spiegazione non ancora disponibile.\par\medskip
\noindent\textbf{quiz\_pagine\_011\_020\_019}: A. Vicino a zero, $1-\cos x\sim x^2/2$ e l\textbackslash{}'integranda \textbackslash{}`e equivalente a $2/x$.\par\medskip
\noindent\textbf{quiz\_pagine\_041\_050\_008}: D. Risposta ufficiale indicata nella tabella della decima simulazione.\par\medskip
\noindent\textbf{quiz\_pagine\_051\_060\_007}: D. Per il criterio del confronto, se $0\le f\le g$ e l'integrale di $g$ converge, converge anche quello di $f$.\par\medskip
\noindent\textbf{quiz\_pagine\_051\_060\_017}: A. La presenza della singolarità $x=2$ impedisce la convergenza per intervalli che la contengono.\par\medskip
\noindent\textbf{quiz\_pagine\_051\_060\_018}: C. Vicino a zero l'integranda si comporta come $x^{1-\alpha}$, all'infinito come una costante per $x^\alpha$.\par\medskip
\noindent\textbf{quiz\_pagine\_051\_060\_020}: A. Vicino a zero $\sin x/(x\sqrt{x})\sim1/\sqrt{x}$; all'infinito converge per confronto assoluto.\par\medskip
\noindent\textbf{quiz\_pagine\_061\_070\_009}: A. Per alfa positivo, vicino a zero serve alfa<1 e all'infinito serve 2alfa>1.\par\medskip
\noindent\textbf{quiz\_pagine\_061\_070\_013}: D. Una primitiva è -3x\textasciicircum{}(-1/3); il limite superiore vale zero e resta 3/a\textasciicircum{}(1/3).\par\medskip
\noindent\textbf{quiz\_pagine\_061\_070\_015}: C. Vicino a zero sin t/t\textasciicircum{}alfa è equivalente a t\textasciicircum{}(1-alfa); diverge se alfa>=2.\par\medskip
\noindent\textbf{quiz\_pagine\_061\_070\_016}: A. La primitiva è arctan t e il salto fra gli estremi infiniti vale pi.\par\medskip
\noindent\textbf{quiz\_pagine\_061\_070\_020}: B. Poiché arctan x tende a pi/2, il comportamento è quello di x\textasciicircum{}alfa; convergenza per alfa<-1.\par\medskip
\noindent\textbf{quiz\_pagine\_071\_080\_002}: B. L'integranda è integrabile in $0$ e per $x\to+\infty$ è equivalente a $x^{-3/2}$; l'integrale converge ed è positivo, come indicato nella soluzione ufficiale.\par\medskip
\noindent\textbf{quiz\_pagine\_071\_080\_015}: A. Per $x\to+\infty$ il termine $1/\sqrt{x}$ ha integrale divergente e l'altro addendo è non negativo; diverge l'integrale totale.\par\medskip
\noindent\textbf{quiz\_pagine\_081\_090\_007}: D. Per $x\to0^+$, $1/(\sqrt{x}+x^2)\sim x^{-1/2}$, integrabile in $0$; l\textbackslash{}textquotesingle integranda è positiva, quindi il valore è positivo.\par\medskip
\noindent\textbf{quiz\_pagine\_081\_090\_018}: A. Per il criterio del confronto, $0\leq f\leq g$ e la convergenza di $\int g$ implicano la convergenza di $\int f$. La formulazione D è ambigua: un integrale improprio di una funzione con segno variabile può non convergere, ma il termine ``indeterminato non è definito.\par\medskip
\noindent\textbf{quiz\_pagine\_081\_090\_025}: C. L\textbackslash{}textquotesingle integranda è $(1-\pi)x^{3-\pi}$: poiché $3-\pi>-1$, l\textbackslash{}textquotesingle integrale converge e vale $(1-\pi)/(4-\pi)<0$.\par\medskip
\noindent\textbf{quiz\_pagine\_091\_100\_020}: B. Una primitiva è $3/[2(x-3)^2]$. Il valore $F(4)-F(3^+)$ diverge a $-\infty$.\par\medskip
\noindent\textbf{quiz\_pagine\_101\_110\_003}: C. Una primitiva è $\frac12(\log x)^2$, che tende a $+\infty$ per $x\to+\infty$.\par\medskip
\noindent\textbf{quiz\_pagine\_111\_120\_009}: C. Per $x>0$, $4/\sqrt x\leq(5+\sin(1/x))/\sqrt x\leq6/\sqrt x$. Integrando da $0$ a $1$ si ottiene $8\leq I\leq12$; le disuguaglianze sono strette per l'integrale.\par\medskip
\noindent\textbf{quiz\_pagine\_111\_120\_023}: B. Per $x\to+\infty$ l'integranda e' equivalente a $x^{-1/2}$, il cui integrale diverge positivamente.\par\medskip
\noindent\textbf{quiz\_pagine\_111\_120\_026}: D. Il minorante $x^{-10}$ e' integrabile, quindi non decide la convergenza: $f=x^{-10}$ converge, mentre $f=1$ diverge.\par\medskip
\noindent\textbf{quiz\_pagine\_111\_120\_040}: B. $f(x)\sin x\sim3x^4$, quindi l'integranda e' equivalente a $3x^{4-\alpha}$; converge presso $0$ se $4-\alpha>-1$, cioe' $\alpha<5$.\par\medskip
\noindent\textbf{quiz\_pagine\_131\_140\_014}: A. Vicino a $0$ l\textbackslash{}'integranda e\textbackslash{}'e equivalente a $x^{-1/3}$; vicino a $1/2$ e\textbackslash{}'e equivalente a $C(1/2-x)^{-1/3}$. Entrambe le singolarita\textbackslash{}'a sono integrabili.\par\medskip
\noindent\textbf{quiz\_pagine\_141\_150\_014}: D. L\textbackslash{}'integrale parziale e\textbackslash{}'e limitato superiormente da $1$; puo\textbackslash{}'e convergere oppure tendere a $-\infty$, non a $+\infty$.\par\medskip
\noindent\textbf{quiz\_pagine\_151\_160\_016}: B. $\arctan t\to\pi/2>0$, quindi l\textbackslash{}'integranda e\textbackslash{}'e equivalente a $(\pi/2)t^\alpha$; converge per $\alpha<-1$.\par\medskip
\noindent\textbf{quiz\_pagine\_151\_160\_018}: D. L\textbackslash{}'esponenziale domina ogni polinomio: $(2+e^x)/(1+x^5)\to+\infty$; l\textbackslash{}'integrale diverge a $+\infty$.\par\medskip
\noindent\textbf{quiz\_pagine\_171\_180\_012}: C. L'integranda e' positiva ed equivalente a $x^{-4}$ per $x\to+\infty$; l'integrale converge a un valore strettamente positivo. Non si possono dichiarare equivalenti due integrali come numeri.\par\medskip
\noindent\textbf{quiz\_pagine\_171\_180\_020}: E. L'integranda e' positiva; e' limitata su intervalli compatti e per $|x|\geq1$ e' maggiorata da $e^{-|x|}$, integrabile.\par\medskip
\noindent\textbf{quiz\_pagine\_181\_190\_002}: E. Si ha $|e^x\sin x|\leq e^x$ e $\int_{-\infty}^{-1}e^x\,dx<\infty$: converge assolutamente, quindi anche semplicemente. La risposta E e' vera, ma anche C e' falsa; il testo distingue impropriamente convergenza assoluta e semplice.\par\medskip
\noindent\textbf{quiz\_pagine\_181\_190\_011}: E. Integranda positiva e continua; per $x\to+\infty$ e' $O(x^{-3/2})$, quindi converge a un numero strettamente positivo.\par\medskip
\noindent\textbf{quiz\_pagine\_191\_200\_011}: D. Per $x\ge4$ il denominatore e' positivo; $f(x)\sim1/x^2$ e l'integrale improprio converge a un numero strettamente positivo.\par\medskip
\noindent\textbf{quiz\_pagine\_211\_220\_004}: C. Per $x\to+\infty$, $\cosh x\sim e^x/2$, quindi $f(x)\sim\sqrt2 e^{-3x/4}$: integrabile. Vicino a zero $f(0)=1$.\par\medskip
\noindent\textbf{quiz\_pagine\_211\_220\_006}: E. Il numeratore e' equivalente a $2x^4$, percio' l'integranda e' equivalente a $2x^{4-\alpha}$: converge se $4-\alpha>-1$.\par\medskip
\noindent\textbf{quiz\_pagine\_211\_220\_007}: A. Integrando per parti si ottiene $[-(t/2+1/4)e^{-2t}]_0^{+\infty}=1/4$.\par\medskip
\noindent\textbf{quiz\_pagine\_211\_220\_008}: B. L'integranda e' proporzionale a $x^{-1/3}$: l'integrale su $[a,+\infty)$ diverge perche' $1/3\leq1$.\par\medskip
\noindent\textbf{quiz\_pagine\_211\_220\_009}: D. Calcolo: $\int_0^a(5x)^{-1/3}dx=\frac3{2\cdot5^{1/3}}a^{2/3}=\frac3{10}(5a)^{2/3}$. Corrisponde alla risposta D.\par\medskip
\noindent\textbf{quiz\_pagine\_211\_220\_010}: B. $\sin x/x\to1$: converge.\par\medskip
\noindent\textbf{quiz\_pagine\_211\_220\_011}: A. Vicino a zero e' equivalente a $1/x$: diverge.\par\medskip
\noindent\textbf{quiz\_pagine\_211\_220\_012}: B. Numeratore $\sim x^2/2$, quindi integranda $\sim x^{-1/3}/2$: converge.\par\medskip
\noindent\textbf{quiz\_pagine\_211\_220\_013}: B. Integranda $\sim\frac12 x^{2-\pi}$, con $2-\pi<-1$: diverge.\par\medskip
\noindent\textbf{quiz\_pagine\_211\_220\_014}: A. Per $x\to4^-$, il rapporto e' equivalente a $\pi(4-x)^{1-\pi}$, non integrabile.\par\medskip
\noindent\textbf{quiz\_pagine\_211\_220\_015}: A. Per $x\to+\infty$, l'integranda tende a $1$.\par\medskip
\noindent\textbf{quiz\_pagine\_211\_220\_016}: A. Converge assolutamente per confronto con $1/x^\pi$.\par\medskip
\noindent\textbf{quiz\_pagine\_211\_220\_017}: B. Per $x\to+\infty$ il rapporto e' equivalente a $1/x$ (con $M$ mantissa limitata), quindi diverge.\par\medskip
\noindent\textbf{quiz\_pagine\_211\_220\_026}: B. $|e^{-3t^2}\sin t|\leq e^{-3t^2}$ e il gaussiano e' integrabile.\par\medskip
\noindent\textbf{quiz\_pagine\_211\_220\_027}: A. Converge assolutamente per confronto con $e^{-3t^2}$.\par\medskip
\noindent\textbf{quiz\_pagine\_221\_230\_003}: A. Vicino a zero $f(x)\sim x^{-1/2}$, integrabile; all'infinito $|f(x)|\leq x^{-3/2}$, anch'esso integrabile.\par\medskip
\noindent\textbf{quiz\_pagine\_251\_260\_016}: C. Il prodotto $fg$ e' identicamente nullo fuori da $[-1,1]$, quindi l'integrale improprio coincide con quello sull'intervallo compatto.\par\medskip
\noindent\textbf{quiz\_pagine\_281\_290\_001}: D. L'integrale gaussiano converge ed e' pari a $\sqrt\pi>0$.\par\medskip
\noindent\textbf{quiz\_pagine\_281\_290\_018}: A. La primitiva e' $\arctan x$; ai due estremi si ottiene $\pi/2-(-\pi/2)=\pi$.\par\medskip
\noindent\textbf{quiz\_pagine\_281\_290\_029}: D. Con $t=\sqrt x$, $dx=2t\,dt$ e l'integrale vale $2\int_0^1\frac{dt}{1+t}=2\log2>0$ (diverso da 2).\par\medskip
\noindent\textbf{quiz\_pagine\_311\_320\_014}: A. Per $x\to0^+$, $1-\cos x\sim x^2/2$, dunque $x/(1-\cos x)\sim2/x$; diverge positivamente.\par\medskip
\noindent\textbf{quiz\_pagine\_321\_330\_032}: C. Vicino a zero $\sin x\sim x$; l\textbackslash{}textquotesingle integrando \textbackslash{}`e equivalente a $x^{1-7\alpha}$, integrabile quando $1-7\alpha>-1$.\par\medskip
\noindent\textbf{quiz\_pagine\_321\_330\_035}: A. $\sin(x^6)\sim x^6$ per $x\to0$, quindi la condizione \textbackslash{}`e $6-\alpha>-1$, cio\textbackslash{}`e $\alpha<7$. Le opzioni D ed E sono identiche nella fonte.\par\medskip
\noindent\textbf{quiz\_pagine\_331\_340\_007}: B. $\int_1^A3x^{-3}\,dx=[-3/(2x^2)]_1^A\to3/2$.\par\medskip
\noindent\textbf{quiz\_pagine\_331\_340\_024}: B. Per $x\to+\infty$ l\textbackslash{}textquotesingle integranda \textbackslash{}`e asintotica a $x^{-1/2}>0$; l\textbackslash{}textquotesingle integrale diverge a $+\infty$.\par\medskip
\noindent\textbf{quiz\_pagine\_361\_370\_007}: D. Per il confronto $0\leq f\leq g$, la convergenza dell\textbackslash{}textquotesingle integrale di $g$ implica quella dell\textbackslash{}textquotesingle integrale di $f$ sul medesimo intervallo.\par\medskip
\noindent\textbf{quiz\_pagine\_361\_370\_017}: A. L\textbackslash{}textquotesingle integrale da $\alpha\leq2$ attraversa la singolarita $x=2$, dove $f(x)\sim1/(x-2)$ e non e integrabile. La singolarita in $x=3$ e integrabile.\par\medskip
\noindent\textbf{quiz\_pagine\_361\_370\_018}: C. Vicino a zero $\arctan x\sim x$, dunque serve $\alpha<2$; all\textbackslash{}textquotesingle infinito $\arctan x\to\pi/2$, dunque serve $\alpha>1$.\par\medskip
\noindent\textbf{quiz\_pagine\_361\_370\_020}: A. Per $x\to0^+$, $f(x)\sim x^{-1/2}$, integrabile; per $x\to+\infty$, $|f(x)|\leq x^{-3/2}$, integrabile.\par\medskip
\noindent\textbf{quiz\_pagine\_361\_370\_029}: A. Per $\alpha>0$, vicino a zero l\textbackslash{}textquotesingle integranda e $t^{-\alpha}$ (serve $\alpha<1$), all\textbackslash{}textquotesingle infinito e $t^{-2\alpha}$ (serve $\alpha>1/2$). Per $\alpha\leq0$ non converge all\textbackslash{}textquotesingle infinito.\par\medskip
\noindent\textbf{quiz\_pagine\_371\_380\_003}: D. Una primitiva e $-3x^{-1/3}$, dunque il valore e $3a^{-1/3}$.\par\medskip
\noindent\textbf{quiz\_pagine\_371\_380\_005}: C. Per $t\to0^+$, $\sin t/t^\alpha\sim t^{1-\alpha}$; diverge se $1-\alpha\leq-1$, cioe $\alpha\geq2$.\par\medskip
\noindent\textbf{quiz\_pagine\_371\_380\_006}: A. La primitiva e $\arctan t$ e la differenza tra i limiti e $\pi$.\par\medskip
\noindent\textbf{quiz\_pagine\_371\_380\_010}: B. Poiche $\arctan x\to\pi/2$, il criterio asintotico riduce a $\int_1^\infty x^\alpha dx$, convergente se $\alpha<-1$.\par\medskip
\noindent\textbf{quiz\_pagine\_371\_380\_017}: B. Per $x\to+\infty$ il numeratore e $O(x^{3/2})$, il denominatore e $O(x^3)$: coda assolutamente integrabile. L\textbackslash{}textquotesingle integrando e positivo per $x$ sufficientemente grande, ma il segno complessivo dell\textbackslash{}textquotesingle integrale richiede attenzione. La fonte indica B.\par\medskip
\noindent\textbf{quiz\_pagine\_371\_380\_030}: A. Il termine $1/\sqrt x$ e non negativo e ha integrale divergente all\textbackslash{}textquotesingle infinito, quindi diverge anche la somma con $(1-\cos x)/x^2\geq0$.\par\medskip
\noindent\textbf{quiz\_pagine\_381\_390\_012}: D. Per $x\to0^+$, $1/(\sqrt{x}+x^2)\sim x^{-1/2}$, integrabile in $0$; l\textbackslash{}textquotesingle integrando e positivo, dunque il valore e positivo. La disuguaglianza E e falsa perche $\sqrt{x}\geq x$ su $(0,1)$.\par\medskip
\noindent\textbf{quiz\_pagine\_391\_400\_003}: A. Il criterio del confronto per integrali impropri di funzioni non negative garantisce $0\leq I_f\leq I_g<+\infty$. Senza non negativita non si puo applicare.\par\medskip
\noindent\textbf{quiz\_pagine\_391\_400\_010}: C. La potenza $x^{3-\pi}$ e integrabile in $0$ perche $3-\pi>-1$; il fattore $1-\pi<0$ rende il valore negativo: $(1-\pi)/(4-\pi)$.\par\medskip
\noindent\textbf{quiz\_pagine\_401\_410\_012}: B. L\textbackslash{}textquotesingle integranda \textbackslash{}`e negativa e si comporta come $-(x-3)^{-3}$ in $3^+$; l\textbackslash{}textquotesingle integrale diverge a $-\infty$.\par\medskip
\noindent\textbf{quiz\_pagine\_411\_420\_003}: C. Una primitiva \textbackslash{}`e $(\log x)^2/2$, che diverge positivamente all\textbackslash{}'infinito.\par\medskip
\noindent\textbf{quiz\_pagine\_411\_420\_029}: C. Per $x>0$, $4\leq5+\sin(1/x)\leq6$; integrando $4/\sqrt x$ e $6/\sqrt x$ si ottiene $8\leq I\leq12$.\par\medskip
\noindent\textbf{quiz\_pagine\_421\_430\_006}: B. Per $x\to+\infty$ l\textbackslash{}'integranda \textbackslash{}`e equivalente a $1/\sqrt x$, il cui integrale diverge.\par\medskip
\noindent\textbf{quiz\_pagine\_421\_430\_009}: D. Il minorante $x^{-10}$ \textbackslash{}`e integrabile, ma ci\textbackslash{}`o non decide la convergenza di $f$: scegliere $f=x^{-10}$ oppure $f=1$.\par\medskip
\noindent\textbf{quiz\_pagine\_421\_430\_023}: B. Vicino a zero $f(x)\sin x\sim3x^4$, quindi l\textbackslash{}'integranda \textbackslash{}`e equivalente a $3x^{4-\alpha}$ e converge se $4-\alpha>-1$. Anche A \textbackslash{}`e vera, ma meno precisa: ambiguit\textbackslash{}`a nel quesito.\par\medskip
\noindent\textbf{quiz\_pagine\_431\_440\_005}: C. Scrivendo $\log|x^2/(x-1)|=2\log x-\log|x-1|$, l\textbackslash{}'integrale completo vale $2\log4-2$, dunque la met\textbackslash{}`a \textbackslash{}`e $\log4-1$. Le singolarit\textbackslash{}`a logaritmiche sono integrabili.\par\medskip
\noindent\textbf{quiz\_pagine\_431\_440\_024}: B. Per $x\to+\infty$, il numeratore \textbackslash{}`e asintotico a $x^{3/2}$ e il denominatore a $x^2$: integranda $\sim1/\sqrt{x}$, il cui integrale diverge positivamente.\par\medskip
\noindent\textbf{quiz\_pagine\_431\_440\_032}: D. Le primitive integrali parziali sono non decrescenti e hanno limite in $[0,+\infty]$. La risposta ufficiale D esclude impropriamente il caso $f\equiv0$, per cui il limite \textbackslash{}`e $0$: nessuna alternativa descrive esattamente tutti i casi.\par\medskip
\noindent\textbf{quiz\_pagine\_441\_450\_018}: D. Per $u=x+1\to0^+$, $f\sim-\log u/u^2$, la cui primitiva diverge all\textbackslash{}'estremo.\par\medskip
\noindent\textbf{quiz\_pagine\_471\_480\_002}: B. Per $t\to\infty$, $\arctan t\to\pi/2$, quindi l'integranda \textbackslash{}`e equivalente a $(\pi/2)t^\alpha$. L'integrale converge se $\alpha<-1$.\par\medskip
\noindent\textbf{quiz\_pagine\_471\_480\_004}: D. Il termine $e^x/x^5\to+\infty$: l'integranda \textbackslash{}`e positiva e definitivamente maggiore di $1$, perci\textbackslash{}`o l'integrale diverge a $+\infty$.\par\medskip
\noindent\textbf{quiz\_pagine\_471\_480\_015}: E. $1-\pi<0$ e $\pi-3\in(0,1)$: l'integrale della potenza diverge, con segno negativo.\par\medskip
\noindent\textbf{quiz\_pagine\_471\_480\_022}: C. Una primitiva \textbackslash{}`e $(x-1)e^x$; tra $-\infty$ e $0$ si ottiene $-1$.\par\medskip
\noindent\textbf{quiz\_pagine\_481\_490\_021}: C. Per $x\to\infty$, $\log(1+x^{-4})\sim x^{-4}$: l\textbackslash{}'integrale converge e, essendo l\textbackslash{}'integranda positiva, ha valore positivo. L\textbackslash{}'alternativa E usa impropriamente $\sim$ tra valori di integrali definiti, che non sono necessariamente equivalenti.\par\medskip
\noindent\textbf{quiz\_pagine\_501\_510\_019}: A. Vicino a zero $\sin x/(x\sqrt x)\sim1/\sqrt x$, integrabile. All\textbackslash{}'infinito $|f(x)|\le x^{-3/2}$, anch\textbackslash{}'esso integrabile: converge assolutamente su $(0,+\infty)$.\par\medskip
\noindent\textbf{quiz\_pagine\_501\_510\_039}: D. Le primitive parziali $\int_0^R f$ sono non crescenti in $R$; il loro limite esteso esiste ed \textbackslash{}`e finito non positivo oppure $-\infty$.\par\medskip
\noindent\textbf{quiz\_pagine\_511\_520\_019}: C. Per $x\to0^+$, $\arctan x\sim x$, quindi serve $\alpha<2$; per $x\to+\infty$, $\arctan x\to\pi/2$, quindi serve $\alpha>1$.\par\medskip
\noindent\textbf{quiz\_pagine\_511\_520\_039}: A. L'integranda \textbackslash{}`e $1/[\sqrt{x}(1+x)]$: vicino a $0$ \textbackslash{}`e dell'ordine di $x^{-1/2}$, all'infinito di $x^{-3/2}$; entrambi gli integrali convergono.\par\medskip
\subsection*{Capitolo 6}
\subsubsection*{6.1 -- Metodi risolutivi di alcune equazioni differenziali ordinarie}
\noindent\textbf{quiz\_pagine\_001\_010\_008}: D. Spiegazione non ancora disponibile.\par\medskip
\noindent\textbf{quiz\_pagine\_011\_020\_020}: D. Il polinomio caratteristico \textbackslash{}`e $(r-2)^2$, radice doppia $r=2$.\par\medskip
\noindent\textbf{quiz\_pagine\_021\_030\_021}: D. $y\prime=-\sin x$, quindi $\sin^2x+\cos^2x=1$.\par\medskip
\noindent\textbf{quiz\_pagine\_031\_040\_008}: A. Spiegazione non ancora disponibile.\par\medskip
\noindent\textbf{quiz\_pagine\_041\_050\_037}: non verificata / nessuna alternativa corretta. Soluzione non riportata nel PDF; da verificare.\par\medskip
\noindent\textbf{quiz\_pagine\_041\_050\_042}: non verificata / nessuna alternativa corretta. Soluzione non riportata nel PDF; da verificare.\par\medskip
\noindent\textbf{quiz\_pagine\_071\_080\_009}: B. L'equazione $x'=x-6$ ha soluzione $x(t)=6+Ce^t$; da $x(-1)=5$ segue $C=-e$, dunque $x(t)=6-e^{t+1}$.\par\medskip
\noindent\textbf{quiz\_pagine\_071\_080\_018}: A. L'equazione $y''-(2/3)y'=0$ ha soluzione $y=C_1+C_2e^{2x/3}$. Il vincolo $y(0)=0$ impone $C_1=-C_2$, lasciando infiniti valori di $C_2$.\par\medskip
\noindent\textbf{quiz\_pagine\_081\_090\_012}: C. Per $y=x$ e $x\neq0,1$ il secondo membro è $1$, uguale a $y\prime=1$; nei punti $y=0,1$ il denominatore si annulla.\par\medskip
\noindent\textbf{quiz\_pagine\_081\_090\_022}: C. Si scrive $y\prime-x^2y=0$ (lineare) e, per $y\neq0$, $dy/y=x^2dx$ (separabile). La soluzione nulla è anch\textbackslash{}textquotesingle essa ammessa.\par\medskip
\noindent\textbf{quiz\_pagine\_091\_100\_014}: D. Per linearità e omogeneità dell'equazione, ogni combinazione lineare di soluzioni è soluzione. Una costante additiva non lo è necessariamente.\par\medskip
\noindent\textbf{quiz\_pagine\_091\_100\_027}: C. Le soluzioni sono $y=A\cos(x/2)+B\sin(x/2)$. La sola soluzione costante è $y\equiv0$; imporre un unico valore iniziale non determina una soluzione unica.\par\medskip
\noindent\textbf{quiz\_pagine\_101\_110\_007}: B. Sostituendo $y=e^2\sin t$ si trova $-e^2\sin t=ae^2\sin t+b$, dunque $a=-1$, $b=0$.\par\medskip
\noindent\textbf{quiz\_pagine\_101\_110\_011}: D. L\textbackslash{}textquoteright equazione caratteristica è $\lambda(\lambda+1/5)=0$, perciò $y=C_1+C_2e^{-t/5}$, non una combinazione di seni e coseni.\par\medskip
\noindent\textbf{quiz\_pagine\_111\_120\_013}: D. L'equazione non e' lineare per $e^y$ ne' separabile nella forma data. Le funzioni proposte non soddisfano l'identita' differenziale.\par\medskip
\noindent\textbf{quiz\_pagine\_111\_120\_028}: D. Si riscrive $y'+\frac1{x+1}y=\frac{x}{x+1}$, equazione lineare del primo ordine.\par\medskip
\noindent\textbf{quiz\_pagine\_131\_140\_015}: C. Si riscrive $y\prime-\frac{3}{2x+3}y=\frac{2x}{2x+3}$: equazione lineare del primo ordine, per $x\neq-3/2$.\par\medskip
\noindent\textbf{quiz\_pagine\_151\_160\_004}: E. Il campo e\textbackslash{}'e definito per $x\geq0$. Per $x=2\pi$ il fattore $1-\cos x$ si annulla; la soluzione costante soddisfa l\textbackslash{}'equazione.\par\medskip
\noindent\textbf{quiz\_pagine\_151\_160\_009}: C. $\log t$ e\textbackslash{}'e definito per $t>0$, dove $x\equiv0$ risolve l\textbackslash{}'equazione.\par\medskip
\noindent\textbf{quiz\_pagine\_151\_160\_015}: D. Le radici caratteristiche sono $-3,-7$; ogni combinazione lineare di $e^{-3x},e^{-7x}$ tende a $0$.\par\medskip
\noindent\textbf{quiz\_pagine\_161\_170\_017}: B. $y'=e^xe^y$ e\textbackslash{}'e a variabili separabili, ma non lineare.\par\medskip
\noindent\textbf{quiz\_pagine\_171\_180\_022}: E. $y'=4e^{4t}$ e $y''=16e^{4t}=4y'$.\par\medskip
\noindent\textbf{quiz\_pagine\_191\_200\_008}: D. Per $y=x$ e $x\ne0$, $y'=1$ e il secondo membro e' $\sin1-\sin1+1=1$. Tuttavia l'equazione non e' definita in $x=0$, quindi il dato $y(0)=0$ e' formalmente problematico: la soluzione ufficiale D va intesa per prolungamento.\par\medskip
\noindent\textbf{quiz\_pagine\_191\_200\_020}: A. La funzione $g(y)=\arctan y+5y^2$ ha zeri in $0$ e in un punto negativo $\alpha$. Per $y_0>0$ la soluzione cresce; per $y_0=0$ e $y_0=\alpha$ e' costante; per $y_0<\alpha$ cresce anch'essa. Quindi C, indicata come ufficiale, non e' corretta come ``se e solo se``; A e' la scelta logicamente coerente.\par\medskip
\noindent\textbf{quiz\_pagine\_201\_210\_003}: E. Le prime quattro differiscono solo per una costante. L'ultima ha dominio diverso e derivata $\frac{2x}{3(x^2+1)}-\frac1{3x}$, dunque non e' una primitiva della funzione data.\par\medskip
\noindent\textbf{quiz\_pagine\_201\_210\_022}: B. La definizione richiede che $F$ sia derivabile in ogni punto di $I$ e che $F'=f$ su tutto $I$.\par\medskip
\noindent\textbf{quiz\_pagine\_201\_210\_023}: C. $f$ e' pari e strettamente positiva. La sua integrale da $0$ a $x$ e' dispari, non pari.\par\medskip
\noindent\textbf{quiz\_pagine\_201\_210\_024}: E. Con $u=1+e^{2x}$, $du=2e^{2x}dx$, dunque una primitiva e' $u^{3/2}/3+C$; la E corrisponde a $C=-8$.\par\medskip
\noindent\textbf{quiz\_pagine\_261\_270\_015}: D. Con $u=6+x^3$, $du=3x^2dx$, quindi l'integrale e' $\frac43[\log u]_6^7$.\par\medskip
\noindent\textbf{quiz\_pagine\_261\_270\_019}: A. Il numeratore e' la derivata del denominatore; $[\log(x^3+x^2+9)]_1^3=\log(45/11)$.\par\medskip
\noindent\textbf{quiz\_pagine\_261\_270\_030}: B. Le primitive differiscono per una costante additiva. $2\log|f|$ ha invece derivata $2f'/f$, quindi e' errata (salvo casi particolari).\par\medskip
\noindent\textbf{quiz\_pagine\_281\_290\_002}: E. Derivando $e^{4t}$ si ottiene $y'(t)=4e^{4t}=4y(t)$.\par\medskip
\noindent\textbf{quiz\_pagine\_281\_290\_008}: B. Per linearita' e omogeneita', ogni combinazione lineare di soluzioni e' soluzione.\par\medskip
\noindent\textbf{quiz\_pagine\_291\_300\_008}: D. $y'=-\sin x$ e $\sin^2x+\cos^2x=1$.\par\medskip
\noindent\textbf{quiz\_pagine\_301\_310\_014}: C. Per definizione una primitiva e' derivabile in ogni punto di $I$ e ha derivata uguale a $f$ ovunque. La sola esistenza della derivata non e' implicita nella formulazione B.\par\medskip
\noindent\textbf{quiz\_pagine\_301\_310\_017}: A. Con $u=1+e^x$ si ha $du=e^x dx$ e $\int\sqrt u\,du=\frac23u^{3/2}+C$.\par\medskip
\noindent\textbf{quiz\_pagine\_311\_320\_015}: D. Il polinomio caratteristico $(\lambda-2)^2=0$ ha radice doppia 2: soluzione $(C_1+C_2x)e^{2x}$.\par\medskip
\noindent\textbf{quiz\_pagine\_321\_330\_017}: D. Per $y=\cos x$ si ha $(y\prime)^2+y^2=\sin^2x+\cos^2x=1$.\par\medskip
\noindent\textbf{quiz\_pagine\_331\_340\_018}: A. Per $y=e^{2x}$ vale $y\prime=2e^{2x}=2y$, dunque $y\prime-2y=0$.\par\medskip
\noindent\textbf{quiz\_pagine\_341\_350\_010}: E. L\textbackslash{}textquotesingle equazione caratteristica $(\lambda-1)^2=0$ ha radice doppia $1$, quindi $y=(a+bx)e^x$.\par\medskip
\noindent\textbf{quiz\_pagine\_341\_350\_029}: A. $x\prime=3(x-2)$ ha soluzione generale $x(t)=2+Ce^{3t}$; la condizione $x(3)=2$ impone $C=0$.\par\medskip
\noindent\textbf{quiz\_pagine\_351\_360\_010}: E. L\textbackslash{}textquotesingle equazione caratteristica $(\lambda-1)^2=0$ ha radice doppia $1$, quindi $y=(a+bx)e^x$.\par\medskip
\noindent\textbf{quiz\_pagine\_351\_360\_029}: A. $x\prime=3(x-2)$ ha soluzione generale $x(t)=2+Ce^{3t}$; la condizione $x(3)=2$ impone $C=0$.\par\medskip
\noindent\textbf{quiz\_pagine\_391\_400\_007}: C. La forma $y\prime-x^2y=0$ e lineare del primo ordine; $dy/y=x^2dx$ (per $y\neq0$) e a variabili separabili; $y=0$ e soluzione aggiuntiva.\par\medskip
\noindent\textbf{quiz\_pagine\_411\_420\_007}: B. Derivando due volte $y=e^2\sin t$ si ottiene $y\prime\prime=-y$: $a=-1$, $b=0$.\par\medskip
\noindent\textbf{quiz\_pagine\_411\_420\_011}: D. L\textbackslash{}'equazione caratteristica \textbackslash{}`e $\lambda(\lambda+1/5)=0$: le soluzioni sono $C_1+C_2e^{-t/5}$, non oscillanti. Nota: anche B \textbackslash{}`e falsa se ``limitate in $\mathbb R$\textbackslash{}'\textbackslash{}' significa limitate su tutta la retta: solo le costanti lo sono, e sono infinite; quindi B \textbackslash{}`e vera.\par\medskip
\noindent\textbf{quiz\_pagine\_411\_420\_033}: D. Il termine $e^y$ rende l\textbackslash{}'equazione non lineare e non separabile; le funzioni proposte non soddisfano l\textbackslash{}'identit\textbackslash{}`a.\par\medskip
\noindent\textbf{quiz\_pagine\_421\_430\_011}: D. Riscrivendo $y\prime+\frac{1}{x+1}y=\frac{x}{x+1}$ si riconosce un\textbackslash{}'equazione lineare del primo ordine.\par\medskip
\noindent\textbf{quiz\_pagine\_431\_440\_025}: A. $y\prime-\frac{1}{x^2+1}y=\frac{x^2}{x^2+1}$ \textbackslash{}`e lineare del primo ordine.\par\medskip
\noindent\textbf{quiz\_pagine\_431\_440\_030}: E. Non \textbackslash{}`e lineare n\textbackslash{}'e separabile in $x,y$; $y=x$ darebbe $1=2$, falso. L\textbackslash{}'opzione D \textbackslash{}`e una condizione algebrica, non una soluzione costante.\par\medskip
\noindent\textbf{quiz\_pagine\_451\_460\_010}: A. ATTENZIONE: nel testo originale la risposta A contiene logaritmi senza valori assoluti, non definiti su tutto il dominio reale di $h$; B presenta i valori assoluti ed \textbackslash{}`e una primitiva sui componenti appropriati. La chiave ufficiale indica A. Inoltre per $x<-1$ la A \textbackslash{}`e definita, ma $\log(x+1)$ non lo \textbackslash{}`e: incongruenza del quesito. Verificare il dominio scelto.\par\medskip
\noindent\textbf{quiz\_pagine\_451\_460\_028}: A. $(x+1)/(x-1)=1+2/(x-1)$, dunque una primitiva \textbackslash{}`e $x+2\log|x-1|=x+\log((1-x)^2)$.\par\medskip
\noindent\textbf{quiz\_pagine\_451\_460\_029}: C. Il teorema della media integrale per funzioni continue garantisce $f(c)=\frac1{5-4}\int_4^5 f=\int_4^5 f$. La sola integrabilit\textbackslash{}`a non basta a garantire che il valor medio sia assunto.\par\medskip
\noindent\textbf{quiz\_pagine\_451\_460\_030}: E. $F\prime(x)=-x^2/\cosh(x^2)\le0$ per ogni $x$, quindi $F$ \textbackslash{}`e decrescente su $\mathbb R$.\par\medskip
\noindent\textbf{quiz\_pagine\_461\_470\_012}: C. La media integrale nulla equivale a $\int_0^4 f(t)\,dt=0$; per il teorema fondamentale, $F(4)-F(0)=0$.\par\medskip
\noindent\textbf{quiz\_pagine\_461\_470\_013}: E. $F\prime=f\le0$ implica che $F$ \textbackslash{}`e non crescente: E \textbackslash{}`e corretta. Anche B \textbackslash{}`e vera, perch\textbackslash{}'e $F(x)=\int_0^x f\le0$ per $x\ge0$. Il quesito presenta quindi almeno due alternative vere; la chiave ufficiale indica E.\par\medskip
\noindent\textbf{quiz\_pagine\_461\_470\_018}: C. Dove $f(x)\neq0$, la derivata di $\log|f(x)|$ \textbackslash{}`e $f\prime(x)/f(x)$; quindi la primitiva generale \textbackslash{}`e $\log|f(x)|+c$.\par\medskip
\noindent\textbf{quiz\_pagine\_461\_470\_022}: C. Per $0<x<4$ si ha $1<e^{\sqrt{x}}<e^2<8$; integrando $4<I<4e^2<32$.\par\medskip
\noindent\textbf{quiz\_pagine\_471\_480\_001}: D. Il polinomio caratteristico ha radici $-7$ e $-3$: $y=C_1e^{-7x}+C_2e^{-3x}\to0$ per ogni coppia di costanti.\par\medskip
\noindent\textbf{quiz\_pagine\_481\_490\_002}: B. $y\prime=e^x e^y$ separa le variabili: $e^{-y}\,dy=e^x\,dx$. Non \textbackslash{}`e lineare in $y$.\par\medskip
\noindent\textbf{quiz\_pagine\_481\_490\_009}: C. La soluzione del problema \textbackslash{}`e $y(x)=\cos(2x)$, che oscilla e non ha limite per $x\to+\infty$.\par\medskip
\noindent\textbf{quiz\_pagine\_491\_500\_006}: E. $y\prime=4e^{4t}$ e $y\prime\prime=16e^{4t}=4y\prime$.\par\medskip
\noindent\textbf{quiz\_pagine\_501\_510\_020}: D. Annullando il fattore dipendente da $x$ si ottengono le soluzioni costanti $x=0$, $x=2$ e $x=\pi/2+2k\pi$ con $x\ge0$. La risposta (d) segue la convenzione $\mathbb N=\{0,1,\dots\}$.\par\medskip
\noindent\textbf{quiz\_pagine\_501\_510\_040}: A. Le soluzioni sono $y=C_1e^x+C_2e^{-x}$: scegliendo $C_1=1$, $C_2=0$ si ha una soluzione illimitata; scegliendo $C_1=0$ si hanno soluzioni limitate.\par\medskip
\noindent\textbf{quiz\_pagine\_511\_520\_020}: A. L'equazione caratteristica \textbackslash{}`e $r(r+1)=0$; la soluzione generale \textbackslash{}`e $y=C_1+C_2e^{-x}$, con infinite soluzioni costanti scegliendo $C_2=0$.\par\medskip
\noindent\textbf{quiz\_pagine\_511\_520\_040}: C. L'equazione caratteristica \textbackslash{}`e $(r-2)(r-3)=0$: il forzamento $e^{2t}$ \textbackslash{}`e risonante con una radice semplice, quindi si cerca $x_p=Cte^{2t}$; $C=-1$.\par\medskip
\subsubsection*{6.2 -- Buona positura del problema di Cauchy}
\noindent\textbf{quiz\_pagine\_041\_050\_014}: A. Risposta ufficiale indicata nella tabella della decima simulazione.\par\medskip
\noindent\textbf{quiz\_pagine\_051\_060\_001}: A. Per $a=1$ il dato iniziale non appartiene al dominio reale $x\ge 2$ della radice.\par\medskip
\noindent\textbf{quiz\_pagine\_111\_120\_003}: B. Le soluzioni sono $y(x)=Ce^{4x/5}$; la condizione $y(1)=0$ impone $C=0$.\par\medskip
\noindent\textbf{quiz\_pagine\_121\_130\_010}: A. Riscrivendo $y'-\frac{1}{x^2+1}y=\frac{x^2}{x^2+1}$ si ottiene un\textbackslash{}'equazione lineare del primo ordine.\par\medskip
\noindent\textbf{quiz\_pagine\_121\_130\_015}: E. Non e\textbackslash{}'e lineare ne\textbackslash{}'e separabile nella forma data. $y=x$ darebbe $1=\cos0+1=2$, impossibile. La condizione dell\textbackslash{}'opzione d) non definisce una soluzione costante.\par\medskip
\noindent\textbf{quiz\_pagine\_161\_170\_024}: C. La soluzione e\textbackslash{}'e $y(x)=\cos(2x)$, che oscilla e non ammette limite per $x\to+\infty$.\par\medskip
\noindent\textbf{quiz\_pagine\_181\_190\_024}: C. $y(t)=A\cos2t+B\sin2t$; $A=0$, $2B=2$, quindi $y(t)=\sin2t$ e $y(\pi)=0$.\par\medskip
\noindent\textbf{quiz\_pagine\_241\_250\_028}: A. La continuita' dei coefficienti $a,b$ in un intorno di $x_0$ garantisce esistenza e unicita' locale. Il testo originale parla impropriamente di intorno di $(x_0,y_0)$, mentre $a,b$ dipendono solo da $x$.\par\medskip
\noindent\textbf{quiz\_pagine\_261\_270\_013}: A. $1\leq2+\sin(1/x)\leq3$. Il confronto con $1/\sqrt{x}$ dà convergenza vicino a zero e divergenza a $+\infty$.\par\medskip
\noindent\textbf{quiz\_pagine\_261\_270\_014}: C. $1\leq2+\sin(1/x)\leq3$. Il confronto con $1/\sqrt{x}$ dà convergenza vicino a zero e divergenza a $+\infty$.\par\medskip
\noindent\textbf{quiz\_pagine\_261\_270\_024}: C. Il prodotto $fg$ e' nullo fuori dall'intervallo $[-1,1]$; l'integrale improprio coincide con quello su tale intervallo.\par\medskip
\noindent\textbf{quiz\_pagine\_271\_280\_020}: C. Il polinomio caratteristico e' $(r-2)(r-3)$: $e^{2t}$ e' soluzione omogenea, quindi per risonanza si prova $Cte^{2t}$.\par\medskip
\noindent\textbf{quiz\_pagine\_281\_290\_010}: A. Il polinomio caratteristico e' $(r+1)^2$. Con i dati iniziali si ricava $y(t)=te^{-t}\to0$. La fonte rappresenta l'alternativa D con un simbolo di limite inesistente.\par\medskip
\noindent\textbf{quiz\_pagine\_281\_290\_011}: A. $y'(0)=\pi^2>0$ e $y''(0)=3+2\pi^3>0$. La spiegazione originale contiene un refuso nel calcolo di $y'(0)$.\par\medskip
\noindent\textbf{quiz\_pagine\_281\_290\_012}: B. Separando le variabili, $y=C(t+1)$ per $t>-1$. Da $y(0)=1$ segue $C=1$ e $y(1)=2$.\par\medskip
\noindent\textbf{quiz\_pagine\_301\_310\_016}: E. Per il teorema fondamentale del calcolo, $F\prime(x)=-e^{\sin^2x}<0$ per ogni $x$.\par\medskip
\noindent\textbf{quiz\_pagine\_361\_370\_001}: A. La radice richiede $x\geq2$: il dato iniziale $x(0)=1$ non e ammissibile. La soluzione costante $x=2$ soddisfa solo il dato $a=2$.\par\medskip
\noindent\textbf{quiz\_pagine\_371\_380\_024}: B. La soluzione generale e $x(t)=6+Ce^t$; da $x(-1)=5$ segue $C=-e$, quindi $x(t)=6-e^{t+1}$.\par\medskip
\noindent\textbf{quiz\_pagine\_381\_390\_003}: A. L\textbackslash{}textquotesingle equazione caratteristica e $r(r-2/3)=0$: $y=C_1+C_2e^{2x/3}$. Da $y(0)=0$ segue $C_1=-C_2$, lasciando un parametro libero e infinite soluzioni.\par\medskip
\noindent\textbf{quiz\_pagine\_381\_390\_017}: C. Per $y=x$, il membro destro vale $(x-x^2)/(x-x^2)=1=y\prime$ dove il denominatore non si annulla, cioe per $x\neq0,1$.\par\medskip
\noindent\textbf{quiz\_pagine\_401\_410\_006}: D. Per linearit\textbackslash{}`a e omogeneit\textbackslash{}`a ogni combinazione lineare di soluzioni \textbackslash{}`e soluzione, in particolare $3\varphi_1-2\varphi_2$.\par\medskip
\noindent\textbf{quiz\_pagine\_401\_410\_019}: C. La soluzione generale \textbackslash{}`e $y=C_1\cos(x/2)+C_2\sin(x/2)$; la sola soluzione costante \textbackslash{}`e quella nulla. Una sola condizione iniziale non determina entrambe le costanti.\par\medskip
\noindent\textbf{quiz\_pagine\_411\_420\_023}: B. La soluzione generale \textbackslash{}`e $Ce^{4x/5}$; $y(1)=0$ implica $C=0$.\par\medskip
\noindent\textbf{quiz\_pagine\_441\_450\_009}: D. In $t=0$, $x\prime(0)=\sqrt2(1-1+2)=2\sqrt2\ne0$; esiste una soluzione locale per regolarit\textbackslash{}`a del campo, non costante.\par\medskip
\noindent\textbf{quiz\_pagine\_451\_460\_006}: A. Vicino a $0$ l\textbackslash{}textquotesingle integranda \textbackslash{}`e dell\textbackslash{}textquotesingle ordine $x^{-1/3}$; vicino a $1/2$ \textbackslash{}`e dell\textbackslash{}textquotesingle ordine $(1/2-x)^{-1/3}$. Entrambe le singolarit\textbackslash{}`a sono integrabili perch\textbackslash{}`e $1/3<1$.\par\medskip
\noindent\textbf{quiz\_pagine\_451\_460\_031}: D. Dato che $\int_0^{+\infty}e^{-x}dx$ converge, le parti positive di $f$ sono integrabili; la parte negativa pu\textbackslash{}`o avere integrale finito o infinito, quindi il totale converge oppure tende a $-\infty$.\par\medskip
\subsubsection*{6.3 -- Studi qualitativi}
\noindent\textbf{quiz\_pagine\_201\_210\_004}: B. Su ogni periodo di lunghezza $\pi$, l'integrale di $|\sin x|$ vale $2$: la somma diverge.\par\medskip
\noindent\textbf{quiz\_pagine\_201\_210\_009}: C. Per $2<x\le3<\pi$ si ha $\sin x>0$, quindi il segno e' costante. Inoltre $f(x)\sim\sin2/\sqrt{x-2}$.\par\medskip
\noindent\textbf{quiz\_pagine\_201\_210\_010}: D. La parte intera $[x]$ e' limitata e continua a tratti su $[0,5]$. ATTENZIONE: nell'originale $I_5$ ha una radice di $x(x+2)(x-12)$, negativa per $1\le x\le10$: l'integranda non e' reale. Il testo originale e' problematico; D resta vera.\par\medskip
\noindent\textbf{quiz\_pagine\_201\_210\_021}: D. Un integrale improprio e' divergente se il limite delle integrali parziali e' infinito, con segno definito. ATTENZIONE: il testo considera solo il limite all'infinito superiore e non esplicita la gestione del limite in $0^+$; C e' anche un caso particolare di D.\par\medskip
\noindent\textbf{quiz\_pagine\_301\_310\_009}: A. Un integrale improprio e' divergente positivamente o negativamente quando il limite degli integrali troncati e' rispettivamente $+\infty$ o $-\infty$. La mancata convergenza include anche il caso oscillante. Nota: la fonte usa la dizione "divergente" in questo senso ristretto.\par\medskip
\noindent\textbf{quiz\_pagine\_441\_450\_015}: C. $h(y)=\arctan y+5y^3$ ha il segno di $y$; per $y_0>0$ la soluzione cresce nel suo intervallo di esistenza, per $y_0=0$ \textbackslash{}`e costante.\par\medskip
\noindent\textbf{quiz\_pagine\_451\_460\_007}: C. Riscrivendo $y'=\frac{3}{2x+3}y+\frac{2x}{2x+3}$ si riconosce un'equazione lineare del primo ordine, non a variabili separabili in generale.\par\medskip
\subsection*{Capitolo 7}
\subsubsection*{7.1 -- Concetti di base}
\noindent\textbf{quiz\_pagine\_091\_100\_009}: C. Le radici quadrate di $i$ sono opposte: $z_1=e^{i\pi/4}$ e $z_2=-z_1$, quindi $z_1/z_2=-1\in\mathbb R$.\par\medskip
\noindent\textbf{quiz\_pagine\_091\_100\_025}: B. Poiché i coefficienti sono reali, anche $2-i$ è radice. Le tre radici distinte $0,2+i,2-i$ impongono grado almeno $3$.\par\medskip
\noindent\textbf{quiz\_pagine\_201\_210\_006}: C. La somma $(y')^2+y^2+1$ e' sempre almeno $1$ per funzioni reali.\par\medskip
\noindent\textbf{quiz\_pagine\_201\_210\_007}: E. L'ultima si riscrive $y'-y=\sin x$, lineare del primo ordine.\par\medskip
\noindent\textbf{quiz\_pagine\_201\_210\_017}: C. La soluzione generale e' $y=3+Ce^{2x}$; da $y(0)=4$ segue $C=1$.\par\medskip
\noindent\textbf{quiz\_pagine\_201\_210\_018}: C. $y'=2\cos2x$, quindi $(y')^2/4+y^2=\cos^22x+\sin^22x=1$. ATTENZIONE: la tabella ufficiale indica D, ma D e' falsa: il suo lato sinistro e' $4$, non $1$.\par\medskip
\noindent\textbf{quiz\_pagine\_261\_270\_022}: A. Per l'equazione lineare del primo ordine e' sufficiente la continuita' dei coefficienti in un intervallo contenente $x_0$. La notazione $(x_0,y_0)$ nella fonte e' imprecisa.\par\medskip
\noindent\textbf{quiz\_pagine\_461\_470\_014}: E. La forma $x\prime=h(x)(1-t)$ \textbackslash{}`e separabile. Per $x=2\pi$, $1-\cos(2\pi)=0$, quindi $x(t)\equiv2\pi$ \textbackslash{}`e soluzione; $x=-2\pi$ non appartiene al dominio reale di $\sqrt{x}$.\par\medskip
\noindent\textbf{quiz\_pagine\_461\_470\_019}: C. Il termine $\log t$ \textbackslash{}`e definito per $t>0$. Su tale dominio $x(t)=0$ soddisfa $x\prime=0=x\log t$.\par\medskip
\end{document}
